\section{extra}
***
JOY SERIES (14-16)
\subsection{快}

14. SW 10B 408: 014; DXB 217 (10B 10b); GL vol 11: 4660; Duàn 502 (10B 25b); TKJ 1441; Ozaki vol. 5: 1000.
HG:            (快), kuài, ``in good spirits"

G:	喜也。
SSF: 从心,
PIF: 夬聲。

is 'glad'. [[is '(a way of being) glad (scil. as a mood rather than a reaction)'.]]
It has xīn ``heart'' as a semantic constituent,
guài (GSR 312a: ``archer's thimble, ...") is the phonetic constituent.

PIF     夬[古賣切 (Gu(ngyùn: 古邁切); LH kuas, OCM kwrêts] phonetic in 快 [苦夬切; LH khuas, OCM khwrêts]
PR	Rhyme and homogenous initials.


15. SW 10B 408: 015; DXB 217 (10B 10b); GL vol 11: 4660; Duàn 502 (10B 25b); TKJ 1441; Ozaki vol. 5: 1000, see
variants in WGY: 447.
HG: y (愷), k&i, ``joyful"

G:	樂也。
SSF: 从心,
PIF: 豈聲。

is 'delighted'. [[is '(a way of being) delighted (scil. prototypically on major festive occasions or the like)'.]]
It has xīn ``heart'' as a semantic constituent,
q( (GSR 548a: ``joyous, ...") is the phonetic constituent.

HG One of the anomalies in Shuōwén is the recurrence of the same character under two different radicals: the character k(i 愷 reoccurs under the radical q+ 豈 (Shuōwén 5A 16a, TKJ: 663, k(i 愷康也。从心、豈，豈亦聲。[苦亥切]). These must count as cases of carelessness even when the meanings attributed are different. (Conceivably, it might be a matter of a copyist's carelessness). Note that many characters have a range of different meanings and are not for that reason listed several times. In the other entry for k(i 愷 the semantic gloss is close enough, but the graphologi- cal analysis is sharply different: q+ 豈 is said to be a semantic constituent, indeed the anomaly of the position of the radical in the second place is explained by the left to right order of the explanation of the graph. The semantic function of q+ 豈 is made to appear plausible by Xǔ Shèn's



gloss of the character in terms of triumphal music. It is almost as if, at the time Xǔ Shèn explained the graph k(i 愷 under the heart radical, he did not think of the rather special gloss he had for q+豈. Xǔ Shèn offends against two of his own rules:
1. There must be one and only one entry for each graph.
2. There should be one and only one correct graphological analysis for any given graph.
On the second point, one might add that Xǔ Shèn must surely have been considering many possible alternatives of graphological analysis much of the time, but he made a decision not to communicate his hesitations or deliberations to his readers. He reports his final results and not his supporting arguments most of the time.
PIF      Q+ 豈 [墟喜切; LH kh\#iB, kh\$iB, OCM kh*+ i) , kh\#i) ] phonetic in k(i 愷 [苦亥切; LH kh\#iB, OCM kh*+ i) ]
PR Pronunciation looks different only from a modern point of view, but much closer in LH and in OCM than in Middle Chinese.

16. SW 10B 408: 016; DXB 217 (10B 10b); GL vol 11: 4661; Duàn 502 (10B 25b); TKJ 1441; Ozaki vol. 5: 1000, see
variants in WGY: 447.


HG:
G: SSF: PIF:

(心),
快心。从心, 匧聲。

qiè,
'in a good mood'.
It has xīn ``heart'' as a semantic constituent,
qiè (SW: ``box, ...") is the phonetic constituent.

G Note the absence of the customary final y\% 也. Bisyllabic glosses do not regularly omit y\% 也. Shuōwén usage of y\% 也 is not always consistent. Thus, in two consecutive entries even, we have both the gloss 玉聲也 (1A 6.71, qiáng 䱗:玉聲也 'the sound given of by jade (when struck)'); and the gloss 玉聲 (1A 6.70, líng 玲:玉聲 'the sound given of by jade (when struck)').
    The presence versus absence of y\% 也 is by no means always arbitrary as the following sequence from the beginning of the fish radical shows:
qū 魼	魚也。从魚去聲。
nà [魚+納]	魚。似鼈，無甲，有尾，無足，口在腹下。从魚納聲。
    The absence of y\% 也 in the second case is clearly connected to the fact that the explanation continues in an elaborate way.
PIF    匧[苦叶切; LH khep, OCM khêp] phonetic in 心 [苦叶切; LH khep, OCM khêp]
PR	The phonetic constituent is homophonous with the character.
***



MINOR THINK SERIES (17-18)

17. SW 10B 408: 017; DXB 217 (10B 10b); GL vol 11: 4661b; Duàn 502 (10B 25b); TKJ 1441; Ozaki vol. 5: 1001.


HG:
G: SSF: PIF:

; (念),
常思也。从心,
今聲。


niàn,
is 'to be constantly preoccupied by something'. It has xīn ``heart'' as a semantic constituent,
jīn (GSR 651a: ``now, ...") is the phonetic constituent.

G   This is a non-discursive straightforward simple analytic gloss defining niàn 念 as a subtype of sī 思 with the added feature of long duration. Thus in instances like these, Xǔ Shèn himself performs for us the task we usually need to take upon ourselves in explanatory paraphrase (EP).
    Xǔ Shèn does not generally tend to gloss meanings in terms of genus and species. Under the heart radical such analytic ``Aristotelian'' definitions are rare, whereas – for example – the glosses for kinship terms are gratifyingly analytic throughout. There are certain other sections of the dictionary in which the analytic mode prevails to the point of being pervasive, as in the case of the
horse radical in which not a single kind of horse is given a generic gloss 馬也, although there are
some cases of glosses like 馬名 'designation of a horse'.
PIF     今[居音切; LH k\$m, OCM k\#m] phonetic in 念 [奴店切; LH nemC, OCM nǐms]
PR	On the face of it a clear case of f*i sh*ng 非聲. Note the case of 貪: 欲物.从貝今聲: 今 [居音切; ] phonetic in 貪 [他含切].

18. SW 10B 408: 018; DXB 217 (10B 11a); GL vol 11: 4661b; Duàn 503 (10B 26a); TKJ 1441; Ozaki vol. 5: 1001.


HG:
G: SSF: PIF:

! (䓛),
思也。从心, 付聲。


fū,
is 'think of something'.
It has xīn ``heart'' as a semantic constituent,
fù ((GSR 136a: ``give, ...") is the phonetic constituent.

G   Presumably also the rare word fū 䓛 would have had its specific features, but in cases like these we are unable to venture any reconstruction of the nuances involved. We provide our extended explanatory paraphrase (EP) only where we feel we have something to say.
PIF     付[方遇切; LH puoC, OCM poh] phonetic in 䓛 [甫無切; LH puo, OCM po]
***
19. SW 10B 408: 019; DXB 217 (10B 11a); GL vol 11: 4662; Duàn 503 (10B 26a); TKJ 1441; Ozaki vol. 5: 1001.


HG:
G: SSF1:

¦ (憲), 敏也。从心,

xiàn,
is 'skillful'.
It has xīn ``heart'' as a semantic constituent,




SSF2: PIF:

从目， 害省聲。

it (also) has mù ``eye'' as a semantic constituent.
It has an abbreviated form of hài (GSR 314a:"... to hurt ... ``) as the phonetic.

G Xǔ Shèn's gloss remains puzzling: no meaning of 憲 in the literature we know seems to come close to any meaning of 敏.
PIF The notation X sh\%ng sh*ng X省聲 'has X as an abbreviated phonetic' is an important analytical tool in Shuōwén. It allows Xǔ Shèn to suggest that a given form which has no phonetic value because it is never used to write any word must often be taken as a graphic abbreviated form of a fuller form which does have a phonetic value. Xǔ Shèn's judgements on what is an abbrevia- tion for what, often seem conjectural, and tend to remain conjectural rather than observational, but he needs such problematic analyses to allow him to explain structurally opaque characters according to his principles. According to his principle, on the one hand Xǔ Shèn often presents plausible suggestions, as in the present case. On the other hand, the freedom this procedure gives him can be abused. In any case, the procedure does allow him to greatly reduce the number of f*i sh*ng 非聲 that would be necessary without it. (Hé Línyí 1998: 898)
害[胡蓋切; LH gâs, OCM gâts] phonetic in 憲 [許建切; LH hianC, OCM hngans]
PR F*i sh*ng 非聲. Pre-final -n and -t are both dental. Also the initials are different, but homoge- neous.


20. SW 10B 408: 020; DXB 217 (10B 11a); GL vol 11: 4662b; Duàn 503 (10B 26a); TKJ 1441; Ozaki vol. 5: 1002.


HG:
G:

SSF: PIF:

(憕),
平也。
从心, 登聲。

chéng, ``serene"
is 'peaceful'. [[is '(a way of being) peaceful (scil. not in society but in one's mind)'.]]
It has xīn ``heart'' as a semantic constituent,
d*ng (GSR 883e: ``ascend, ...") is the phonetic constituent.

G	It is true enough that the word píng 平 has no strong psychological side to it and tends to be social in meaning.
PIF     登[都滕切; LH t\#ng, OCM t*+ ng] phonetic in 憕 [直陵切; LH dr\$ng, OCM dr\#ng]
PR	Homogeneous but different initials. One notes that consonantal diphthongs can correspond to simple consonants.

21. SW 10B 408: 021; DXB 217 (10B 11a); GL vol 11: 4663; Duàn 503 (10B 26a); TKJ 1442; Ozaki vol. 5: 1002.
HG:  º (戁), n&n, ``awe-struck"




G:	敬也。
SSF: 从心,
PIF: 難聲。

is 'respectful' [[is '(a way of feeling) respect (scil. coupled with awe-struck fear)'.]]
It has xīn ``heart'' as a semantic constituent,
nán (GSR 152d: ``difficult, ...") is the phonetic constituent.

HG Note that this word is also written with the woman radical in which form the word is glossed as g\$ngjìng 恭敬 in Hànyǔ dàcídi(n.
PIF     難 [那干切; LH nân, OCM nân (see TKJ: 522 [ꌅ] for which 難 is a graphic variant)] phonetic in 戁 [女版切; LH nanB, OCM nrân)]

22. SW 10B 408: 022; DXB 217 (10B 11a); GL vol 11: 4663b; Duàn 503 (10B 26a); TKJ 1442; Ozaki vol. 5: 1002.


HG:             (忻),
G:	䌞也。
SSF: 从心,


xīn,
is 'inspire'.
It has xīn ``heart'' as a semantic constituent,

PIF: IQ:

斤聲。
《司馬法》曰：

jīn (GSR 443a: ``axe, ...") is the phonetic constituent. The S(m+ f& says:

善者，忻民之善， ``As for the good person, he inspires goodness in the

閉民之惡。

people and blocks off wickedness in the people".

G	We follow Duàn in our interpretation of the rare character k(i 䌞.
IQ	We have not found this passage, but compare:古者賢王明民之德，盡民之善，故無廢德. (Lǐ Líng 1992, Sīm( F( yì zhù 司馬法譯注).
PIF     斤[舉欣切; LH, k\$n, OCM k\#n] phonetic in 忻 [許斤切; LH h\$n, OCM h\#n]

23. SW 10B 408: 023; DXB 217 (10B 11a); GL vol 11: 4664; Duàn 503 (10B 26a); TKJ 1442; Ozaki vol. 5: 1003.


HG:	(⤋),
G:	遟也。
SSF: 从心,
PIF: 重聲。

zhòng ``tardy"
is 'slow'. [[is '(a way of being) slow (scil. of mind??)'.]] It has xīn ``heart'' as a semantic constituent,
zhòng ((GSR 1188a: ``heavy, ...") is the phonetic constituent.

G	Usually Xǔ Shèn avoids rare graphic variants in his glosses but in this instance the Dà Xú b\%n
preserves such a variant for chí 遲.
PIF     重[柱用切; LH drongB, OCM drong)] phonetic in ⤋ [直隴切; LH dongC , OCM dôngh]
    重聲 might look ambiguous and therefore potentially misleading because the character 重 has two common readings. In fact there is no ambiguity whatsoever in this case. Xǔ Shèn is not




specifying a pronunciation for the head graph. What he is commenting on is the graphological function of the constituent 重 in the head graph. On this matter, he passes a perfectly unambiguous judgment: that function is phonetic, see note 10 p. 12.

            *** GENEROSITY SERIES (24-27)

24. SW 10B 408: 024; DXB 217 (10B 11a); GL vol 11: 4664; Duàn 503 (10B 26b); TKJ 1442; Ozaki vol. 5: 1003.



HG:

d (惲), yùn,

G:

SSF: PIF:

重厚也。
从心, 軍聲。

is 'very generous'. [[EP is (a way of being) generous (scil. perhaps in a 'warm-hearted' way)'.]]
It has xīn ``heart'' as a semantic constituent,
jūn (GSR 458a: ``troop, ...") is the phonetic constituent.

PIF     軍[舉云切; LH kun, OCM kun] phonetic in 惲 [於粉切; LH kun, OCM kun]

25. SW 10B 408: 025; DXB 217 (10B 11a); GL vol 11: 4664b; Duàn 503 (10B 26b); TKJ 1442; Ozaki vol. 5: 1003.



HG:
G: SSF: PIF:

  (惇/数),
厚也。从心, 数聲。

dūn,
is 'generous' / (morally/ psychologically) earnest'. It has xīn ``heart'' as a semantic constituent,
chún (SW: ``cooked, ...") is the phonetic constituent.

G	Compare modern Chinese dūnhòu 敦厚 ``honest, sincere".
PIF   For once, Xǔ Shèn chooses to specify as the phonetic the relevant element in the seal script!
数[常倫切; LH d´zuin, OCM dun] phonetic in 惇 [都昆切; LH tu\#n, tshuin, OCM tûn, tun]

26. SW 10B 408: 026; DXB 217 (10B 11a); GL vol 11: 4664b; Duàn 503 (10B 26b); TKJ 1442; Ozaki vol. 5: 1004, see
variants in WGY: 448.


HG:
G: SSF: PIF:

 (忼),
慨也。从心, 亢聲。


kàng/k+ng, is 'generous'.
It has xīn ``heart'' as a semantic constituent,
kàng (GSR 698a: ``overbearing, ...") is the phonetic con- stituent.




CAS: 一曰:
《易》:
忼龍有悔。

One source points out:
The Book of Changes (SSJZS: 14a) says:
"The excited (or generous) dragon: there will be regrets".

HG   This unusual graphic variant of k&ng 慷 is also attested in F( yán 法言.
G     The format of this entry suggests that Xǔ Shèn may be considering the binome kàng k(i 忼慨as a synonym compound. It is not, perhaps, a coincidence that Xǔ Shèn glosses the first element of this compound by the second whereas he glosses the second element by the whole compound. Since the compound has several well-attested meanings, Xǔ Shèn's gloss does little, in fact, to explain which meaning he had in mind.
    In this instance Duàn rewrites the gloss as 忼慨也 and declares that ``all editions have taken away the character kàng". In fact Duàn has simply conjecturally added a character into the text. When Duàn then goes on to move the phrase 壯士不得志於心也 from the succeeding entry, and adds the characters 於心 for good measure, basing himself on Yùpi&n 玉篇 chapter 8 and Wénxu(n, it is worth pointing out that the Sòng b\%n Yùpi&n 宋本玉篇 (1983: 150) does not in fact have those characters 於心. Thus what Duàn presents as an emendation based on Yùpi&n turns out to be only partially based on the Yùpi&n and to involve a conjectural addition by his own hand. No matter what one may think of the plausibility of Duàn's conjectures and emendations, one must deplore the carelessness of his philological procedures. (See our translation of the next entry.)
PIF 亢[古郎切; LH kâng, OCM kâng] phonetic in 忼 [苦浪切,又，口朗切; LH khângB, khângC, OCM khâng) /h]
PR The case of alternative fǎnqiè glosses 又 regularly involves the recognition that the DXB/ Tángyùn's sources were inconsistent and that they were not sure of their choice. Our early systematic source for readings of Chinese characters in context, the Jīngdi(n shìwén 經典釋文provides an extraordinary wealth of alternative readings which the current dictionaries since Qièyùn times have tended to disregard in their efforts towards standardisation.
CAS Very often, yī yu* 一曰 in Shuōwén is used as a formula equivalent to huò yu* 或曰 meaning something like ``an alternative account/source says", and introducing an alternative to what Xǔ Shèn is claiming in his main text. Quite often the alternative thus presented is a commonly held current view on the meaning of a character which Xǔ Shèn has chosen to disregard for his own graphological reasons. Duàn declares that yī yu* 一曰 is a mistake and ``an addition by an unsubtle mind qi(nrén 淺人", and at first sight one cannot fail to be surprised that the phrase yī yu* 一曰, in this instance, introduces not an alternative gloss or interpretation but a quotation.
See Yìjīng Hexagram Qián 乾.



27. SW 10B 408: 027; DXB 217 (10B 11b); GL vol 11: 4666; Duàn 503 (10B 26b); TKJ 1443; Ozaki vol. 5: 1004, see
variants in WGY: 448.


HG:
EBF: G:

SSF: PIF:

Ž (慨),
忼慨，
壯士不得志也。
从心, 既聲。

k&i,
(as in) k+ngk&i ``generous",
is a man being full of potential who does not get to achieve his ambitions.
It has xīn ``heart'' as a semantic constituent,
jì (GSR 515c: ``finish, ...") is the phonetic constituent.

EBF When Xǔ Shèn opens the gloss for a member of a current binome by the binome itself we translate this by the formula ``as in". We need to investigate whether the characters thus defined do or do not have independent meanings outside the binome.
Duàn conjecturally, and without explanations, adds the character y\% 也 after 忼慨. In fact, his
rewriting of the text only becomes necessary because he has removed from the Shuōwén text the phrase that followed 忼慨. It is indeed as if Duàn never contemplated the systematic significance of y\% 也 in Xǔ Shèn's notational system. If he had, he would surely have found it worth his while to justify his conjectural addition, and at the very least to record his changing of the textus receptus.
G    The problems raised by binomes of this sort are systematically damaging to Xǔ Shèn's project in that his graphological analysis in so far as it attributes meaning to independent graphs must assume a meaning for each graph even in those cases where an independent use of the component characters can not be documented.
The old Wénxu(n commentary (Liù Chén zhù Wénxu(n 六臣注文選) quotes Shuōwén: 說文
曰：慨，太息也, and the Zìlín 字林 has the gloss: 慨，壯士不得志也) ``k(i: for an adult not to have achieved one's ambitions'' (Gǔlín ibidem). Elsewhere we have four quotations of the type: 說文曰：慷慨，壯士不得志於心也. (See Wáng Guìyuán 王貴元 2002: 448).
PIF   既 [written 旣:居未切; LH k\$s, OCM k\#ts] phonetic in 慨 [古䰬切; LH k\#s, OCM k*+ ts]
***

28. SW 10B 408: 028; DXB 217 (10B 11b); GL vol 11: 4666b; Duàn 503 (10B 27a); TKJ 1443; Ozaki vol. 5: 1005, see
variants in WGY: 448.


HG:
G: SSF: PIF:

№ (悃), 愊也。 从心,
困聲。


kǔn,
is 'sincere'.
It has xīn ``heart'' as a semantic constituent,
kùn (GSR 420a: ``distress, ...") is the phonetic constituent.



G   Here as elsewhere Xǔ Shèn's gloss becomes fully comprehensible to us only when we go on to read the following entry. Thus this is another illustration of the fact that Shuōwén must be read as consecutive text and not merely as an accumulation of entries.
Duàn conjecturally rewrites the gloss as follows: 悃,悃䮤也. 至誠也. In fact he manages to
copy into the gloss of kǔn the gloss given for bì 䮤 in the Yùpi&n 玉篇 (p. 150). His annotation fails to make clear what he has done to the text:
1. He has added another copy of the character kǔn 悃 to Xǔ Shèn's gloss.
2. He has added a y\% 也 to the Yùpi&n entry on 䮤 after the first part of the gloss (Yùpi&n: 䮤:悃䮤至誠也).
3. He has taken the gloss for a different word from the Yùpi&n and added it into Xǔ Shèn's text.
4. By his punctuation he has suggested that Xǔ Shèn gave two parallel glosses with essentially the same meaning, something which Xǔ Shèn tends to avoid.
5. He has failed to notify his readers that he has radically rewritten Xǔ Shèn's text.
PIF     困 [苦悶切; LH khu\#nC, OCM khûns] phonetic in 悃 [苦本切; LH khu\#nB, OCM khûn)] Duàn rewrites kùn 困 with the hé 禾 radical for the mù 木, apparently basing himself on Xi(o
Xú b\%n 1987: 208, whereas the standard DXB edition provides the current unproblematic graph.


29. SW 10B 408: 029; DXB 217 (10B 11b); GL vol 11: 4667; Duàn 503 (10B 27a); TKJ 1443; Ozaki vol. 5: 1005, see
variants in WGY: 448.


HG:
G: SSF: PIF:

n (䮤), 誠志也。从心,
䛑聲。


bì,
is 'to be earnest in one's aspirations'.
It has xīn ``heart'' as a semantic constituent,
bì (compare fú (GSR 933a ``... wine vessel, ...")) is the phonetic constituent.

G Instead of inter-defining kǔn 悃 and bì 䮤, which can even form a poetic binome, Xǔ Shèn chooses to provide an analytic gloss of the latter.
    Duàn creates his own gloss for this graph: 悃䮤也. In this instance he does tell the reader what he is doing and what the textus receptus said. He even goes on to provide an argument for his emendation: 今依全書通例訂 ``In our edition we emend this according to the general practice of the whole book.'' We are not sure about the pervasive pattern t\$nglì 通例 he refers to, and how this pattern licences the mechanical introduction of a liánmián 連綿 binome as the gloss of the second member of that binome whenever the two members of the binome are adjacent and in the correct order in Xǔ Shèn's dictionary. But as noted above, Duàn has rewritten the gloss for the first member of the binome k&ngk(i 忼慨 as the whole binome k&ngk(i 忼慨. Thus on the one hand one must object to his rewriting such glosses at all, and on the other hand one must object that his rewriting of these glosses is in fact unsystematic and incoherent.



PIF 䛑[芳逼切](Gu(ngyùn also has 房六切; LH buk, OCM b\#k) phonetic in 䮤 [芳逼切; LH p\$k, OCM pr\#k]
PR 䛑 is said in Shuōwén to be phonetic in both 福 and 䮤. Thus this phonetic constituent (like a fair number of others) has two distinct phonetic values, and the phrase 䛑聲 cannot mean ``pronounced like 䛑", and must mean ``䛑 is phonetic in 䮤". The old Chinese pronunciations of words like fú 福 *p23 and bī 逼 *pr23 appear to have been phonetically close enough according to Zhèngzhāng shàngfáng 鄭張尚芳: (2003). By the time of Middle Chinese, the two pronuncia- tions have diverged. The ascription of the phonetic which is implausible according to the fǎnqiè spellings in Shuōwén turns out to have been plausible in the language of Xǔ Shèn. This illustrates the importance of our Old Chinese and Eastern Han Chinese reconstructions.


30. SW 10B 408: 030; DXB 217 (10B 11b); GL vol 11: 4667b; Duàn 503 (10B 27a); TKJ 1443; Ozaki vol. 5: 1005.
HG: | (愿), yuàn,

G:	謹也。
SSF: 从心,
PIF: 原聲。

is 'diligent'. [[is '(a way of being) diligent (scil. in a positive and warm-hearted manner).]]
It has xīn ``heart'' as a semantic constituent,
yuán (GSR 258a: ``spring, ...") is the phonetic constituent.

HG    Note that yuàn 願, though homophone, is a different etymon.
PIF Xǔ Shèn provides a phonetic constituent here which, though not an entry in Shuōwén, is actually recorded as an allograph of the more complex character 厵 (Shuōwén 11B 2b). One would like to know what motivated Xǔ Shèn's choice of the head graph in this instance, especially since the complex version of yuán 原 is rarely attested epigraphically or in the received literature.
    Xǔ Shèn does not call yuán 原 an abbreviated form of 厵 because 原, in his scheme of things, is simply a small seal graphic variant of 厵.
    原 [small seal of yuán 厵: 愚袁切; LH ngyan, OCM ngwan or OCM ngon] phonetic in 愿 [魚怨切; LH ngyanC, OCM ngons (or ngwans?)]

            *** INTELLIGENCE SERIES (31-33)

31. SW 10B 408: 031; DXB 217 (10B 11b); GL vol 11: 4668; Duàn 503 (10B 27a); TKJ 1443; Ozaki vol. 5: 1006.


HG:
G:

 (慧),
䫞也。


huì, ``clever'' is 'ingenious'.




SSF: PIF:

从心, 彗聲。

It has xīn ``heart'' as a semantic constituent,
huì (GSR 527a: ``broom, ...") is the phonetic constituent.

G Xǔ Shèn chooses to gloss a fairly common word by an exasperatingly rare one. No one can imagine that any speaker of Chinese would learn the meaning of huì 慧 using the gloss xu&n 䫞. It turns out that xu&n 䫞 is in turn glossed as huì 慧 under the man radical, but this fact does not make the present gloss any more illuminating.
PIF    彗[祥歲切; LH zuis, OCM s-wis] phonetic in 慧 [胡桂切; LH Gues, OCM wǐs]
PR The initials are different, and so are the rhymes. However there is good evidence to show that the two endings are taken to rhyme with each other in early literature.

32. SW 10B 408: 032; DXB 217 (10B 11b); GL vol 11: 4668; Duàn 503 (10B 27a); TKJ 1443; Ozaki vol. 5: 1006.
HG: ¢ (䓴), li&o, ``understand"

G:	慧也。
SSF: 从心,
PIF: 尞聲。

is 'clever'. [[is '(a way of being) clever (scil. with respect to a certain question or subject matter).]]
It has xīn ``heart'' as a semantic constituent,
liáo (GSR 1151b: ``burnt-offering, ...") is the phonetic constituent.

G The definiens, in this case, has just been defined in the preceding entry. This is a common way of linking semantic series together in Shuōwén. Conversely, we have found that the definiens can get to be glossed only in the succeeding entry as in the case of bì 䮤 above.
PIF    尞[力照切; LH liauC , OCM riauh] phonetic in 䓴 [力小切; LH leu, OCM riâu]

33. SW 10B 408: 033; DXB 217 (10B 11b); GL vol 11: 4668; Duàn 503 (10B 27a); TKJ 1444; Ozaki vol. 5: 1006.
HG:  / (䮟), xiáo, ``intelligent"

G:	䓴也。
SSF: 从心,
PIF: 交聲。

is 'understand'. [[is '(a way of having) understanding (regarding a wide range of tasks and problems).]]
It has xīn ``heart'' as a semantic constituent,
ji+o (GSR 1166a: ``to cross, ...") is the phonetic constituent.

G   The F&ngyán 方言 dictionary glosses this word in accordance with its use in Mencius as kuài 快 ``pleased, in good spirits'' (Zh\%u Zǔmó 1956: 21), but Xǔ Shèn follows a different tradition. (Incidentally, Wáng Lì's dictionary of classical Chinese (2000: 310) makes a point of claiming that the character is not found in Shuōwén.).



    Stringing together lexical entries in such a way that the definiens in the second lexical entry has been defined in the first is a standard procedure that gives coherence to certain parts of Shuōwén. We call these ``enchained definitions'' (thus we have the pattern: AB也, CA也, DC也).
PIF	ji&o 交 [古爻切; LH kau, OCM krâu (or OCM kriâu ?)] phonetic in xiáo 䮟 [下交切又，古了切, LH gauC , OCM grâuh (or griâuh ?)]
***

34. SW 10B 408: 034; DXB 217 (10B 11b); GL vol 11: 4668b; Duàn 503 (10B 27b); TKJ 1444; Ozaki vol. 5: 1007.


HG:
G: SSF: PIF:

½ (瘱), 靜也。从心,
旹聲。

yì,
is 'calm'.
It has xīn ``heart'' as a semantic constituent,
qiè (SW: ``have difficulty in breathing, ...") is the phonetic constituent.

PIF     qiè 旹 [苦叶切; LH khep, OCM khêp] phonetic in yì 瘱 [於計切; LH )ep, OCM )êp]!
    Xú Xuán 徐鉉 declared that the phonetic qiè 旹 is f*i sh*ng 非聲 ``the wrong phonetic constituent", and he adds the note wèi xiáng 未詳 ``I have no idea what is going on".
Duàn chooses to print the phonetic qiè 旹 in a playfully archaising form.
PR	Note the homogeneous initials Kh and ).
For convenient reference, we list here all the cases we have found in which Xú Xuán professes his ignorance or even openly declares a f*i sh*ng 非聲 ``wrong phonetic constituent":

dài
酄
及也。从隶枲聲。《詩》曰：'' 酄天之未陰雨。"臣鉉等曰：枲非聲。未詳。
jìn
¾
羊名。从羊執聲。汝南平輿有¾亭。讀若晉。臣鉉曰：執非聲，未詳。
qiū
ꇛ
禿ꇛ 也。从鳥尗聲。臣鉉等曰：尗非聲，未詳。
yu&n
ꇢ
鷙鳥也。从鳥䒇聲。臣鉉等曰：䒇非聲。
niè
籋
箝也。从竹爾聲。臣鉉等曰：爾非聲，未詳。
réng
த
驚聲也。从乃省，西聲。籀文த不省。或曰：த，往也。讀若仍。臣鉉等曰：

qián

伲
西非聲。未詳。
虎行皃。从虍文聲。讀若矜。臣鉉等曰：文非聲。未詳。
k(n
毱
i屬。从虎去聲。臣鉉等曰：去非聲。未詳。
gòng
贛
賜也。从貝，竷省聲。臣鉉等曰：竷非聲，未詳。
xù
旭
日旦出皃。从日九聲。若䆞。一曰明也。臣鉉等曰：九，非聲。未詳。
sh*n
曑
商，星也。从晶媎聲。臣鉉等曰：媎，非聲，未詳。
kù
玥
未練治纑也。从麻後聲。臣鉉等曰：後非聲，疑復字譌，當从復省乃得聲。
bó
ƒ
小瓜也。从瓜交聲。臣鉉等曰：交非聲。未詳。
ào
᱈
宛也；室之西南隅。从宀?聲。臣鉉等曰：?非聲未詳。
dài
代
更也。从人弋聲。臣鉉等曰：弋，非聲。




huái
褱
俠也。从衣眔聲。一曰橐。臣鉉等曰：眔非聲，未詳。
yì
裔
衣裾也。从衣富聲。臣鉉等曰：富非声，疑象衣裾之形。
duì
兌
說也。从儿p聲。臣鉉等曰：p，古文兖字，非聲。
jí
泭
目赤也。从見，‰省聲。臣鉉等曰：‰非聲，未詳。
\%
厄
科厄，木節也。从卪厂聲。賈侍中說以爲厄，裹也。一曰厄，蓋也。臣鉉等


曰：厂非聲，未詳。
hé
貈
似狐，善睡獸.从豸舟聲.《論語》曰："狐貈之厚以居."臣鉉等曰：舟非聲，未


bó	佀
huán   萈

néng	能

hè	䉄
chán	½
jìn	搐
y(n	殦
yì	瘱
yì	❰
xié	⥻
dào	悼
gǔn 鯀zhu(n 辶n(o     嵑

j+	戟

lòu	宔
f(ng	瓬

詳。
馬色不純。从馬爻聲。臣鉉等曰：爻非聲
山羊細角者。从兔足，勚聲。凡萈之屬皆从萈。讀若丸。Ⓒ字从此。臣鉉等 曰：勚，徒結切，非聲。
熊屬。足似鹿。从肉Q聲。能獸堅中，故稱賢能；而彊壯，稱能傑也。凡能之屬皆从能。臣鉉等曰：Q非聲。
火熱也。从火高聲。《詩》曰："多將䉄䉄。"臣鉉等曰：高非聲，當从䬮省。小熱也。从火干聲。《詩》曰："憂心½½。"臣鉉等曰：干非聲，未詳。
火餘也。从火聿聲。一曰薪也。臣鉉等曰：聿非聲，疑从聿省。火光也。从炎舌聲。臣鉉等曰：舌非聲，當从甛省。
靜也。从心旹聲。臣鉉等曰：旹，非聲。未详。
怒也。从心刀聲。讀若顡。李陽冰曰："刀非聲。当从刈省。怨恨也。从心录聲。讀若膎。臣鉉等曰：录非聲，未详。
懼也。陳楚謂懼曰悼。从心卓聲。臣鉉等曰：卓非聲，當从罩省。魚也。从魚系聲。臣鉉等曰：系非聲，疑从孫省。
開閉門利也。从門繇聲。一曰縷十紘也。臣鉉等曰：繇非聲，未詳。
有所恨也。从女䒱聲。今汝南人有所恨曰嵑。臣鉉等曰：䒱，古4字。非聲， 當从匘省。
有枝兵也。从戈、倝。《周禮》："戟，長丈六尺。"讀若棘。臣鉉等曰：倝非
聲，義當从榦省，榦，枝也。Even when Xǔ Shèn does not claim phonetic function...
側逃也。从匸丙聲。一曰箕屬。臣鉉等曰：丙非聲。義當从內。會意。
周家搏埴之工也。从瓦方聲。讀若怷破之怷。臣鉉等曰：怷音瓦，非聲，未 詳。




chén



pèi

辰 震也。三月，陽气動，䧠電振，民農時也。物皆生，从乙、匕，象佄達；厂， 聲也。辰，房星，天時也。从二，二，古文上字。凡辰之屬皆从辰。徐鍇曰："匕音化。乙，艸木萌初出曲卷也。"臣鉉等曰：三月陽气成，艸木生上徹於土，故从匕。厂，非聲。疑亦象物之出。
配 酒色也。从酉己聲。臣鉉等曰：己非聲。當从妃省。


35. SW 10B 408: 035; DXB 217 (10B 11b); GL vol 11: 4670; Duàn 503 (10B 27b); TKJ 1444; Ozaki vol. 5: 1007.



HG:
G: SSF: PIF:

C (悊),
敬也。从心, 折聲。


zhé,
is 'respectful'.
It has xīn ``heart'' as a semantic constituent,
zhé (GSR 287a: ``to break, ...") is the phonetic constituent.

HG This is another case where one and the same graph is explained under two radicals in Shuōwén. Under the graph zhé 哲 (2A 22: 059) the Shuōwén says 或从心作悊, so that the gloss for zhé 哲 would be 知也. Without giving any reasoning Duàn wrongly declares that this entry under 哲 is the addition of ``an unsubtle mind'' qi(nrén 淺人. Meanwhile epigraphic evidence of the Warring States period demonstrates that the standard graph for zhé 哲 was indeed 悊. Thus Xǔ Shèn confirms observations which Duàn summarily dismisses as concoctions of an ignoramus. Moreover, in Hàn shū Xíng f( zhì 漢書刑法志, a text much quoted by Duàn, 悊 is used for 哲 in the phrase 明悊之性.
PIF    折[食列切; LH tshat, OCM tet or tat ] phonetic in 悊 [陟列切; LH trat, OCM trat]

36. SW 10B 408: 036; DXB 217 (10B 12a); GL vol 11: 4670b; Duàn 503 (10B 27b); TKJ 1444; Ozaki vol. 5: 1008.


HG:
G: SSF: PIF:

L (悰),
樂也。从心, 宗聲。

cóng,
is 'delighted'.
It has xīn ``heart'' as a semantic constituent,
zǒng (GSR 1003a: ``clan, ...") is the phonetic constituent.

PIF     宗[作冬切; LH tsoung, OCM tsûng] phonetic in 悰 [藏宗切; LH dzoung, OCM dzûng]

37. SW 10B 408: 037; DXB 217 (10B 12a); GL vol 11: 4670b; Duàn 503 (10B 27b); TKJ 1444; Ozaki vol. 5: 1008, see
variants in WGY: 448.
HG: : (恬), tián, ``tranquil"




G:	安也。
SSF: 从心,

is 'peaceful'. [[is '(a way of being) peaceful (scil. in a profound psychological way).]]
It has xīn ``heart'' as a semantic constituent,

PIF:

甛省聲。 It has an abbreviated form of tián (SW:"savoury") as the phonetic.

HG In this case Duàn has decided to rewrite not only Xǔ Shèn's explanations but even the head character itself: he conjectures that Xǔ Shèn must at this point have intended to write not about the character tián 恬 in the received text but about a completely different character a which we are unable to find in HYDZD. [For an explanation for what he is doing one must look for under the character ti(n 䈅 (6A 53b, p. 264) where the received text reads: 䈅 炊伽木。从木舌聲。臣鉉等曰：當从甛省乃得聲。他念切。"Gu& is a wooden implement for cooking. It has 木 as a semantic constituent, 舌 is phonetic. Xuán and the others consider it should take its phonetic constituent from an abbreviated form of 甛. The initial is as in 他 and the final as in 念.'' Duàn goes on to comment: 按徐說非也. 䈅甛銛等字皆從婞聲.婞見ۑ部.轉寫譌為舌耳. ``Xú's explanation is wrong, the characters 䈅, 甛, 銛 , all have 婞 as the phonetic constituent. For 婞 see the radical ۑ. In rewriting it they mistakenly made it into a 舌.'' We are not at this point about to enter into discussion whether Duàn's hypothesis on the graphological structure of these characters is correct or not. (One notes, however, that the right hand element of huà 話, ඖ, is distinguished carefully by Xǔ Shèn from the right hand element of 恬 which in fact corresponds to modern standard 婞. Thus, while Duàn is quite inconsistent when he rewrites the head character in this case and not in the other cases, his insistence on not taking the right hand-side of tián 恬 as the same as the right hand-side of huà 話 is correct. In fact, close to the character ti(n 䈅, Xǔ Shèn deals separately with the graph kuò 桰 which is often and even standardly miswritten as 䈅.)
    What matters in connection with the study of Xǔ Shèn is that the received text of Shuōwén never refers to Duàn's hypothetical graph. Indeed, we commend Xǔ Shèn for addressing the graph tián 恬 which is epigraphically attested on several Hàn seals and the Lóu shòu 婁壽碑 stone inscription (AD 174).
PIF 甛[徒兼切; LH dem, OCM lêm or dêm] phonetic in 恬 [徒兼切; LH dem, OCM lêm or dêm] perfect fǎnqiè match.
PR   The phonetic is homophonous with the character.


38. SW 10B 408: 038; DXB 217 (10B 12a); GL vol 11: 4671; Duàn 503 (10B 27b); TKJ 1444; Ozaki vol. 5: 1008.


HG:	(伭),
G:	大也。
SSF: 从心,
PIF: 灰聲。

huī, ``expansive"
is 'large'. [[is '(a way of being) large (scil. with a poetic stylistic nuance).]]
It has xīn ``heart'' as a semantic constituent,
huī (GSR 950a: ``ashes, ...") is the phonetic constituent.



PIF     灰[呼伭切; LH hu\#i, OCM hw*+ ] phonetic in 伭 [苦回切; LH khw\#, OCM khw*+ ]

39. SW 10B 408: 039; DXB 217 (10B 12a); GL vol 11: 4671; Duàn 503 (10B 27b); TKJ 1444; Ozaki vol. 5: 1009.

HG: ; (恭), gǒng, ``respectful"

G:	肅也。
SSF: 从心,
PIF: 共聲。

is 'serious'. [[is '(a way of being) serious (scil. about one's relation to seniors or superiors).]]
It has xīn ``heart'' as a semantic constituent,
gòng (GSR 1182c: ``join the hands, ...") is the phonetic constituent.

PIF     共 [渠用切; LH k\$ongB, OCM kong? ] phonetic in 恭 [倶容切; LH k\$ong, OCM k(r)ong]

***
AFFECTION SERIES (40 -(42)- 44)

40. SW 10B 408: 040; DXB 217 (10B 12a); GL vol 11: 4672; Duàn 504 (10B 28a); TKJ 1444; Ozaki vol. 5: 1009, see
variants in WGY: 449.


HG:
G: SSF1: SSF2: PIF:

© (䇥), 敬也。从心,
从敬， 敬亦聲。


j(ng, ``respectful'' is 'respectful'.
It has xīn ``heart'' as a semantic constituent,
it (also) has jìng ``respectful'' as a semantic constituent.
jìng (GSR 813a: ``respectful, ...") is at the same time the phonet- ic constituent.

G      There is only a tonal difference between the gloss and the head graph. Duàn asserts that j+ng䇥 refers to the psychological component of respect. The addition of the heart radical to characters in order to symbolise a psychological perspective on concepts has been shown by Páng Pǔ 龐樸2004 to have played a significant role in Gu\%dian 郭店 and Shàngbó 上海博物館 palaeography. However, in this instance the graph with the heart radical might be taken to be used to write a different word and not a different acceptation of the same word as in many of the cases noted by Páng Pǔ, since the pronunciations of the characters j+ng 䇥 and jìng 敬 differ.
PIF    敬[居慶切; LH k\$/ngC, OCM krengh] phonetic in 䇥 [居影切; LH k\$/ngB, OCM kreng)



41. SW 10B 408: 041; DXB 217 (10B 12a); GL vol 11: 4672b; Duàn 504 (10B 28a); TKJ 1444; Ozaki vol. 5: 1009.

HG: 0 (恕),
G:	仁也。
SSF: 从心,


shù, ``fair"
is 'kind-hearted'. [[is '(a way of being) kind-hearted (scil. as in an even-handed way, showing no unfair partial favour).]]
It has xīn ``heart'' as a semantic constituent,

PIF: SA:

如聲。
   [ 彚] ， 古文省。

rú (GSR 94G: ``... resemble, ...") is the phonetic constituent.  is an ancient style graph and it is abbreviated.

G  Xǔ Shèn subsumes the virtue shù 'fairness' under the general virtue rén 仁 'kind-heartedness'. He does the same thing for the character huì 惠 (SW 4B 2a). It appears that Xǔ Shèn takes rén 仁to be a virtue of a person of superior position. Rén 仁 in turn is glossed as 親也 ``to be kind towards", and the warm-heartedness here intended is clearly of the non-symmetric superior benevolent kind.
PIF    如[人諸切; LH ñaB, OCM na) ] phonetic in 恕 [商署切; in LH shaC, OCM nhah]
PR    Note the different initials as well as the different finals. Perhaps one should investigate whether rú and shù ever rhymed.

42. SW 10B 408: 042; DXB 217 (10B 12a); GL vol 11: 4673b; Duàn 504 (10B 28a); TKJ 1445; Ozaki vol. 5: 1010.
HG: № (怡), yí,

G:	和也。
SSF: 从心,
PIF: 台聲。

is 'harmonious'. [[is '(a way of being) harmonious (scil. in an elevated psychological manner).]]
It has xīn ``heart'' as a semantic constituent,
yí (GSR 976p: ``I, me, ...") is the phonetic constituent.

HG This entry is an intrusion into the ``affection'' series.
G Duàn writes 龢 for 和 in the definition. He provides no good reason for interfering with the DXB text.
PIF Yí 台 provides a good example of a phonetic constituent that is used for two distinct readings. Yí 怡 and dài 怠 are both analysed as 从心台聲. According to our interpretation, the meaning of the formula 台聲 is exactly the same in both cases: ``台 is phonetic". The ambiguity in the pronunciation of the phonetic constituent is irrelevant to Xǔ Shèn's analysis.See n° 134 dài 怠below.
台 [與之切; LH j\#, OCM l\#] phonetic in 怡 [與之切; LH j\#, OCM l\#]



43. SW 10B 408: 043; DXB 217 (10B 12a); GL vol 11: 4668b; Duàn 504 (10B 28a); TKJ 1445; Ozaki vol. 5: 1010.
HG:  (慈), cí, ``show protective loving care"

G:	愛也。
SSF: 从心,
PIF: 兹聲。

is 'to show loving care'. [[is '(a way of) showing loving care (scil. towards one's next of kin of the younger generation).]]
It has xīn ``heart'' as a semantic constituent,
zī (GSR 966a:"black") is the phonetic constituent.

G  Duàn rewrites 愛 as 形 without providing sound reasons. He takes Xǔ Shèn not to be entitled to use the graph 愛 in any other meaning than that indicated by the gloss which is part of the graphological analysis. But Xǔ Shèn cannot possibly have overlooked that his analysis of ài 愛was counterintuitive to the general reader. Xǔ Shèn's point was that it was graphologically correct. Duàn, on the other hand, confuses graphological with semantic analysis here as so often elsewhere.
PIF   兹 [子之切; LH tsi\#, OCM ts\# ] phonetic in 慈 [疾之切; LH dzi\#, OCM dz\#]

44. SW 10B 408: 044; DXB 217 (10B 12a); GL vol 11: 4674; Duàn 504 (10B 28a); TKJ 1445; Ozaki vol. 5: 1010, see
variants in WGY: 449.


HG:
G: SSF: PIF:

  (䓗),
愛也。从心, 氏聲。


qí,
is 'to show loving care'.
It has xīn ``heart'' as a semantic constituent,
shì (GSR 867a: `` surname, ...") is the phonetic constituent.

G	This illustrates a third type of enchained definitions in which a series of entries with identical glosses are listed consecutively.
Again, Duàn rewrites the gloss as 形也, as if Xǔ Shèn was not entitled to use the graph 愛 in
any other meaning than that indicated by the gloss which is part of the graphological analysis.
PIF     氏[承旨切; LH dzheB < gieB, OCM ge)] phonetic in 䓗 [巨支切; LH gie, OCM ge]
***

45. SW 10B 408: 045; DXB 217 (10B 12a); GL vol 11: 4674; Duàn 504 (10B 28a); TKJ 1445; Ozaki vol. 5: 1011.

HG:	(忑),	y(,
EBF:   䓗忑，	(as in) qíy(,
G:	不憂事也。 is not to worry about matters.

SSF: PIF:

从心, 虒聲。

It has xīn ``heart'' as a semantic constituent,
sī (SW: ``tiger with horns, ...") is the phonetic constituent.



SPI:   讀若移。 pronounced like yí.
G   Duàn omits the lower part of 憂 and writes something like 微 or rather ⤏ defined by Xǔ Shèn below as 愁也, again without providing good argument. As in the case of ài 愛 Duàn fails to understand the principles of Xǔ Shèn's graphological analysis according to which a graph is glossed not by the main meaning of the word written by that graph but by the meaning that is relevant to the graphological composition of the character. Thus Xǔ Shèn is manifestly entitled to use y\$u 憂 in its current meaning different from the gloss he gives for it in Shuōwén.
PIF 虒 [息移切; LH sie, OCM sle] phonetic in 忑 [移尔切; LH je, OCM le] Duàn tacitly and gratuitously rewrites 尔 as 爾, which may be harmless, but is irritating.
PR   Note the radical difference between the initials.
SPI      Y+ 忑 [移尔切; LH je, OCM le] read like yí. 移 [SW: 弋支切, LH je < jâi, OCM lai]
    In some cases Xǔ Shèn finds that identifying the phonetic constituent is insufficient, and he then adds dúruò 讀若 glosses. See Bottéro & Harbsmeier (2008) and the extensive literature on dúruò quoted in Coblin (1978 &1979).


46. SW 10B 408: 046; DXB 217 (10B 12b); GL vol 11: 4674b; Duàn 504 (10B 28a); TKJ 1445; Ozaki vol. 5: 1011.

HG: < (恮),
G:	謹也。
SSF: 从心,
PIF: 全聲。


qu+n,
is 'diligent'.
It has xīn ``heart'' as a semantic constituent,
quán (GSR 234a: ``complete, ...") is the phonetic constituent.

G The wealth of vocabulary in the area of diligence matches perfectly the prominence of the cultural injunction encouraging diligence in society. But note that many of the words listed are not at all common in the literature.
PIF    全[疾緣切; LH dzyan, OCM dzon] phonetic in 恮 [此緣切; LH tshyan, OCM tshon]

47. SW 10B 408: 047; DXB 217 (10B 12b); GL vol 11: 4674b; Duàn 504 (10B 28a); TKJ 1445; Ozaki vol. 5: 1011.


HG:

8 (恩), *n, ``beneficence"

G:

SSF: PIF:

惠也。

从心, 因聲。

is 'generosity'. [[is '(a kind of) generosity (scil. of an elevated and concrete kind).]]
It has xīn ``heart'' as a semantic constituent,
yīn (GSR 370a: ``rely on , ...") is the phonetic constituent.

SFF     Duàn gratuitously rewrites this: 从心因,因亦聲, as if Shuōwén did not regularly disregard concurrent semantic functions of phonetic constituents.  In any case, if Duàn had  been right




rewriting in this place he should have gone on to rewrite, systematically, all other similar passages in Shuōwén. And there is no way of knowing whether Xǔ Shèn did or did not agree with Duàn's judgements on what is or is not semantic in function.
PIF     因[於真切; LH )in, OCM )in] phonetic in 恩 [烏痕切; LH )\#n, OCM )\#n]
PR	Note the difference in vowel quality. The endings did not rhyme and thus constitute yet another exception to Duàn's rule 同聲必同部.

48. SW 10B 408: 048; DXB 217 (10B 12b); GL vol 11: 4675; Duàn 504 (10B 28b); TKJ 1446; Ozaki vol. 5: 1012, see
variants in WGY: 449.


HG:
G: AG1:

AG2:

SSF: PIF:

(慸),
高也。一曰: 極也。一曰:
困劣也。从心,
帶聲。

dì,
is 'high'.
An (alternative) source says:
it is 'extreme'.
(Yet) another (alternative) source says:
it is 'to be in trouble (and in weak position)'. It has xīn ``heart'' as a semantic constituent,
dài (GSR 315a: ``girdle, ...") is the phonetic constituent.

AG1 Yī yu* 一曰 is ambiguous, as we have seen: it can either refer to ``a certain source", which corroborates what Xǔ Shèn has been saying, or it can refer to ``an additional account'' which conflicts with Xǔ Shèn's own. When there is a sequence of two yī yu* 一曰 supplementary glosses, the meaning intended is always the latter.
AG2 It is rare for Xǔ Shèn to provide three glosses for one given character. What his practice here shows up is that he was writing his dictionary on the basis of a variety of sources from different traditions. Thus he clearly had access to as many as three sources on an exceedingly rare word like dì 慸. It is as if accumulated quotations of alternative sources indicate doubtfulness of the gloss as opposed to qu* 闕 which acknowledges lack of relevant information. One is tempted to imagine that Xǔ Shèn was just casting about for a well attested meaning that was suitable for an explana- tion of the graph.
PIF    帶[當蓋切; LH tâs, OCM tâs] phonetic in 慸 [特計切; LH des, OCM dês]
PR    Note the contrasting vowel quality. Strictly speaking, and given this example, we should not be entitled to argue for the distinction between the vowels here on the basis of xiésh*ng 諧聲series. Quite generally, whenever a phonological constituent of a syllable is found to vary in any 諧聲 series one cannot exclude the possibility that it will also have varied in other members of the same series. Clearly membership in the relevant series cannot be used as an argument for what in fact is known to vary within that series.



49. SW 10B 408: 049; DXB 217 (10B 12b); GL vol 11: 4675; Duàn 504 (10B 28b); TKJ 1446; Ozaki vol. 5: 1012, see
variants in WGY: 449.


HG:
G1:
G2: SSF: PIF:

› (憖), 問也。 謹敬也。从心,
猌聲。


yìn,
is (firstly) 'to ask',
is (secondly) 'to show diligent respect'.
It has xīn ``heart'' as a semantic constituent,
yìn (SW: ``of dogs: show one's teeth in anger...") is the phonetic constituent.

CAS1: 一曰:
說也。
CAS2: 一曰:
甘也。

An (alternative) source says: It is 'to be pleased (yuè)'.
a (second alternative) source says: It is 'to be willing'.

IQ:

《春秋傳》曰： The Springs and Autumns Tradition (i.e. Zuǒzhuàn, SSJZS: 2177b) says:

昊天不憖。又曰：

"Bright Heaven is not willing (to support?).'' It also says (SSJZS: 1852a):

兩君之士皆未"The   gentlemen   of   the   two   rulers   were   not   quite
憖。	satisfied."
G1 Most of the words Xǔ Shèn defines are ambiguous. Xǔ Shèn chooses the one meaning that he deems most relevant to the graphological analysis. We need to know why he chooses in this case to provide two meanings for the rare word yìn. From the size of the entry it appears that this may simply have been yet another case, like that immediately preceding, in which Xǔ Shèn was unable to resolve his uncertainty and preferred to list up whatever he found in his sources.
Duàn rewrites 問 as (k\%n) \$. The motivation for this remains unclear.
PIF     Note that Xǔ Shèn takes lái 來 to be phonetic in yìn 猌 (Shuōwén 10A 13a), where it is a manifest f*i sh*ng 非聲 ``wrong phonetic".
    Yìn 猌 [魚僅切; LH ng\$nC, OCM ngrins or OCM ng\#ns] phonetic in yìn 憖 [魚覲切; LH ng\$nC, OCM ngrins or ng\#ns]
CAS2 Duàn rewrites 甘 as 且 on the basis of Yùpi&n 玉篇 and a gloss in Zuǒzhuàn.
IQ Given the existence of three commentaries to the Annals, it is remarkable that Xǔ Shèn regularly refers to the Zuǒ commentary as the Chūnqiū zhuàn.
The received Zuǒzhuàn (0i 哀 16) text reads: 旻天不弔，不憖遺一老. ``Compassionate
Heaven vouchsafes me no comfort, and has not left me the aged man to support me.'' (tr. Legge)
    Zuǒzhuàn (Wén 文 12.6) (SSJZS: 1852a) ``兩君之士皆未憖也 ``The soldiers of our two armies are not yet satisfied'' (tr. Legge).



50. SW 10B 408: 050; DXB 217 (10B 12b); GL vol 11: 4677b; Duàn 504 (10B 29a); TKJ 1446; Ozaki vol. 5: 1013, see
variants in WGY: 449.


HG:
G: AG1:

AG2:

(懬),
闊也。
一曰: 廣也， 大也。一曰:
Ⓒ也。

kuàng, ``broad-minded"
is 'expansive'. [[is '(a way of being) expansive (scil. in a figurative or psychological way).]]
A (first alternative) source says: it is 'broad',
it is 'big'.
A (second alternative) source says: it is 'to be extensive'.

SSF1: 从心,
SSF2: 从廣，
PIF:	廣亦聲。

It has xīn ``heart'' as a semantic constituent,
it (also) has gu&ng ``broad'' as a semantic constituent,
gu&ng (GSR 707h:"... broad, ...") is at the same time the phonetic constituent.

G The addition of the heart radical to make presumably psychologising derived words or meanings of words is discussed in Páng Pǔ 龐樸 2004. The present word, that is in this case rewritten with a heart radical, is kuàng 曠 'broad' and while kuàng 懬 as well as kuàng 懭 are rarely attested in surviving literature, it does seem natural to surmise that the addition of the heart radical instead of the sun radical was not unrelated to the figurative uses of the word kuàng.
AG1 As in the last lexical entry, Xǔ Shèn is uncharacteristically talkative here; one might even say he is being redundant. The clustered appearance of this phenomenon of redundance raises interesting questions regarding the composition of the work. The second alternative gloss ku&n Ⓒ does not seem to add any new perspective on the semantics of the word.
    It is rare for Xǔ Shèn to quote alternative sources as giving a sequence of meanings, but in the present case the addition of the gloss dà 大 serves to indicate that the relevant meaning intended by gu(ng 廣 is a more abstract one than the literal 'broad'.
    Duàn rewrites this as 廣大也 without providing sound reason, as if Xǔ Shèn did not also elsewhere quote double glosses as alternatives.
AG2 Duàn tacitly and gratuitously moves our AG2 to the end of this dictionary entry, as if two alternative sources did not often precede the graphic analysis in Xǔ Shèn's system.
PIF Gu(ng 廣 [古晃切; LH kuângB, OCM kwâng)] phonetic in kuàng 懬 [苦謗切; LH khuângB/C, OCM khwâng)/h]



51. SW 10B 408: 051; DXB 217 (10B 12b); GL vol 11: 4678b; Duàn 504 (10B 29a); TKJ 1446; Ozaki vol. 5: 1014, see
variants in WGY: 449.

HG: B (悈),
G:	飾(read 飭) 也。
SSF: 从心,
PIF: 戒聲。


jiè, ``warn"
is 'admonish'. [[is '(a way of) admonishing (scil. by way of warning against disorder.]]
It has xīn ``heart'' as a semantic constituent,
jiè (GSR 990a: ``to guard against, ...") is the phonetic constituent.

IQ:

《司馬法》曰： 有虞氏悈於中國。

The Sīm& F& says:
"Shùn published warnings in the central states."

G Duàn rewrites 飾 as 飭, and since he notes that the two characters are often used interchange- ably he thereby acknowledges that his emendation might seem superfluous.
PIF  戒[居拜切; LH k/iC , OCM kr*+ h] phonetic in 悈 [古拜切; LH k/iC , OCM kr*+ h]
IQ The number of quotations in Shuōwén from Sīm( F( is surprising. It is almost as if Xǔ Shèn happened to have a copy of his book at hand wherever he composed his Shuōwén. In fact, the received text of Sīm( F( omits the heart radical and writes jiè 戒. When the heart radical does not make a word with a new pronunciation, its presence must naturally remain optional since it only makes explicit a psychological meaning which the word in question would have had even when written without the heart radical.
    The Sīm( F( text reads as follows: 有虞氏戒于國中，欲民體其命也。 (Lǐ Líng 1992, Sīm( F( yì zhù 司馬法譯注, ch. 2 Ti&nz+ zhī yì 天子之義) ``Shun issued an admonishment to the all in the States/ Capital hoping that the people would understand his instructions."
PR   Tóngyīn 同音 OCM and LH are homophones.

52. SW 10B 408: 052; DXB 217 (10B 12b); GL vol 11: 4679; Duàn 504 (10B 29b); TKJ 1447; Ozaki vol. 5: 1014, see
variants in WGY: 449-50.

HG:  (䓫), y(n,

G:	謹也。
SSF: 从心,
PIF: 䛑聲。

is 'diligent'.
It has xīn ``heart'' as a semantic constituent,
?? (SW: ``what one leans on, ...") is the phonetic constituent.

PIF   䛑[於謹切; LH )\$nC, OCM )\#ns] phonetic in 䓫 [於靳切; LH )\$nB, OCM )\#n)]
PR The pronunciation of 䛑 is not attested outside the compounds in which it occurs. Xǔ Shèn is forced to assign a pronunciation as well as a meaning to a graph which he must have known was not used on its own and therefore had no attested pronunciation on its own. Two points may be




made in his defence: 1. for all we know, Xǔ Shèn may in fact have seen the graph as an indepen- dent character in texts that we do not know today; 2. Xǔ Shèn may have regarded the pronuncia- tion as reconstructed on the basis of the compound characters in which it clearly serves as a phonetic. These points having been made, the fact remains that the assignment of a meaning to the graph remains completely opaque to us. We dwell on this case because it recurs at many points throughout his dictionary.

53. SW 10B 408: 053; DXB 217 (10B 12b); GL vol 11: 4679b; Duàn 504 (10B 29b); TKJ 1447; Ozaki vol. 5: 1015, see
variants in WGY: 450.

HG:	“ (慶),
G:	行賀人也。
SSF1: 从心,
SSF2: 从夂。


qìng,
is 'to go and congratulate others'.
It has xīn ``heart'' as a semantic constituent,
it (also) has zh( (Karlgren: ``SW to arrive from behind, ...") as a semantic constituent.

GI:	吉禮以鹿皮In auspicious ritual one takes the skin of a deer as a ritual

爲贄，
SSF3: 故从鹿省。

present,
hence [the graph] has an abbreviated form of deer as a seman- tic constituent.

HG The radical is neither on the left nor below. In fact it is not anywhere on the ``fringe'' of the character. Such cases are not common, but their existence makes one wonder why Xǔ Shèn was not tempted to take the same approach in cases like y\$u 憂 and ài 愛.
G    The use of xíng 行 to mean 'go and VERB' is unusual and one might even suspect that the idea expressed by xíng here is not that of locomotion but that of deliberate and elaborate conduct as in xíngshì 行事.
SSF2 It is not clear what semantic relevance there should be for this constituent 夂. There is also a whole range of radicals in the Shuōwén for which the semantic specification is either diffuse or irrelevant to the graphs under the radical.
GI It is uncommon for Xǔ Shèn to provide explanations for his assignment of semantic constituents. These are much more common in Sòng comments on the composition of characters.
SFF3 There are 9 cases of explanatory Semantic subsumption formulae with the formula 故从 in which Xǔ Shèn justifies his statement that X 从 Y . The formula 故从 is of crucial importance for a proper understanding of the ubiquitous simple formula 从. The justification of 从 in the formula 故从 is quite regularly in terms of semantic relevance. In other words, one has reasons to say that a graph is 从X to the extent that the meaning of X is relevant to it. This may seem trivial and obvious, but as we shall see there are cases where one might be in doubt about the semantic significance of 从.



HP	Qìng 慶 [丘竟切; LH kh\$angC, OCM khrang(h)] No phonetic constituent identified

54. SW 10B 408: 054; DXB 218 (10B 13a); GL vol 11: 4680b; Duàn 504 (10B 29b); TKJ 1447; Ozaki vol. 5: 1015, see
variants in WGY: 450.


HG:
G: SSF: PIF: IQ:

ƒ (愃),
Ⓒ嫺心腹皃。从心,
宣聲。
《詩》曰： 赫兮愃兮。


xu&n,
descriptive of an expansive frame of mind. It has xīn ``heart'' as a semantic constituent,
xu+n (GSR 164t: ``spread, ...") is the phonetic constituent.
The Book of Odes (SSJZS: 321b) says:
tr. Karlgren: ``how imposing, how conspicuous".

G The technical term mào 皃 is one of the few grammatically orientated semantic terms employed by Xǔ Shèn. The word is often used to characterise insensitive or lively ``adjectival'' meanings. The range of words defined in terms of mào 皃 constitute an indigenous Chinese grammatical category in its own right. See Sūn Liángmíng 2005: 54 (note, however, that the standard character in Shuōwén is 皃, not 貌.)
There is no y\% 也 after mào 皃. It turns out that there are only three exceptions in Shuōwén to
the general rule that mào is not followed by y\%.
    Note the colloquialism xīnfù 心腹 ``mindset'' which Xǔ Shèn seems to use in order to make a euphonic four-character phrase together with the equally colloquial ku&n xián Ⓒ嫺.
Duàn gratuitously rewrites 嫺 as 閒.
PIF     宣[須緣切; LH syan, OCM swan] phonetic in 愃 [況晩切; LH hyanB, OCM hwan)]
IQ	Shījīng (Qí ào 淇奧 55.2). The SSJZS edition has: 赫兮‰兮.

55. SW 10B 408: 055; DXB 218 (10B 13a); GL vol 11: 4681b; Duàn 504 (10B 29b); TKJ 1447; Ozaki vol. 5: 1016.

HG:  z (愻),
G:	順也。
SSF: 从心,
PIF: 孫聲。


xùn,
is 'compliant'. [[is '(a way of being) compliant (scil. in psy- chological attitude and not in action only).]]
It has xīn ``heart'' as a semantic constituent,
sūn (GSR 434a: ``descendent, ...") is the phonetic constituent.

IQ:

《唐書》曰： The Book of Tang (SSJZS: 130c) says:

五品不愻。 ``The five classes (sc. fathers, mothers, elder and younger brothers, sons) are not compliant (in their attitude?)."



G Compare xùn 遜 which refers to compliance in action. One is tempted to say that these two characters write one word which has behavioural as well as psychological meanings.
PIF   孫[思魂切; LH su\#n, OCM sûn] phonetic in 愻 [蘇困切; LH su\#nC, OCM sûns]
IQ The SSJZS edition of Shūjīng Shùndi(n 舜典 writes 五品不遜. The graph with the heart radical is rarely attested and from a modern point of view one is very much inclined to regard the graph with the heart radical as a mere graphic variant of the latter. The interesting question remains to what extend Xǔ Shèn himself shared this modern approach.
Gu\%diàn Zīyī 緇衣 26 has the graph:    (Zhāng Shòuzh\%ng 2000: 143).

56. SW 10B 408: 056; DXB 218 (10B 13a); GL vol 11: 4683; Duàn 505 (10B 30a); TKJ 1448; Ozaki vol. 5: 1016.


HG:
G: SSF: PIF:

IQ:

(忓),
實也。从心,
塞省聲。
《虞書》曰： 剛而忓。

sè,
is '(filled, not empty>) dependable (?)'
It has xīn ``heart'' as a semantic constituent,
It has an abbreviated form of sài (GSR 908a:"... fill, ...") as the phonetic.
The Book of Yu (SSJZS: 138c) says:
"He is firm and trusty."

HG As in the case of the preceding xùn 愻, the standard way of writing the word sè 塞/忓 has replaced the more abstract 心 with the semantically more concrete 土.
PIF   塞 [先代切; LH s\#k, OCM s*+ k] phonetic in 忓 [先則切; LH s\#k, OCM seek]
    This case of graphic abbreviation illustrates a recurrent phenomenon: the element which one would be inclined to regard as the phonetic constituent does not have an independent pronuncia- tion as a separate word, but it recurs elsewhere in other characters with a similar pronunciation. Xǔ Shèn was faced with two choices: either he could arbitrarily assign a reading to a phonetic constituent which does not have an independent reading, or he needed to declare the phonetic constituent an abbreviated form of the character in which it occurs and for which there is a relevant reading. In this instance Xǔ Shèn chooses the latter alternative.
IQ The SSJZS edition of the Shūjīng, G&o Yáo mó 皋陶謨 (4.3.2), writes: 剛而塞 ``he is hard and yet (sincere, true=) just'' (tr. Karlgren).
Note that Xǔ Shèn illustrates the meaning with a clearly psychologising usage of the word.


57. SW 10B 408: 057; DXB 218 (10B 13a); GL vol 11: 4684b; Duàn 505 (10B 30a); TKJ 1448; Ozaki vol. 5: 1017.


HG:
G:

. (恂),
信心也。


xún,
is 'an attitude of trustiness'.




SSF: PIF:

从心, 旬聲。

It has xīn ``heart'' as a semantic constituent,
xún (GSR 392a: ``ten days, ...") is the phonetic constituent.

G Xǔ Shèn emphasises that the good faith in question is not a matter of conduct but of psychological attitude. Xǔ Shèn's gloss, in this instance, happens to provide an implicit justifica- tion of the presence of the radical without using the formula gù cóng 故从.
PIF    旬[詳遵切; LH zuin, OCM s-win] phonetic in 恂 [相倫切; LH suin, OCM swin]
Duàn gratuitously rewrites (and seriously miswrites, it would appear) the fǎnqiè gloss as 相
論.

58. SW 10B 408: 058; DXB 218 (10B 13a); GL vol 11: 4685; Duàn 505 (10B 30a); TKJ 1448; Ozaki vol. 5: 1017.

HG:             (忱),
G:	誠也。
SSF: 从心,


chén,
is 'earnest'. [[is '(a way of being) earnest (scil. in an especially dignified way).]]
It has xīn ``heart'' as a semantic constituent,

PIF: IQ:

䎘聲。
《詩》曰： 天命匪忱。

yín (GSR 656a: ``walk (only Hàn time text ex.), ...") is the phonetic constituent.
The Book of Odes (SSJZS: 604c) says:
"The mandate of Heaven is not earnest>constant."

PIF  䎘[余箴切; LH j\#m, OCM l\#m] phonetic in 忱 [氏任切; iLH dzh\#m, OCM d\#m ?]
IQ The SSJZS edition of Shījīng (Huán 桓 294.1) writes: 天命匪解 ``Heaven's charge (did not slacken=) was never remitted'' (tr. Karlgren). We have to assume that Xǔ Shèn's version of the Shījīng had a text in which the change in the mandate of Heaven was in focus.
PR   Note the non-homogenic initials.

***
THINK SERIES: (59-(61)-63)

59. SW 10B 408: 059; DXB 218 (10B 13a); GL vol 11: 4685b; Duàn 505 (10B 30b); TKJ 1448; Ozaki vol. 5: 1018.
HG: a (惟), wéi,

G:	凡思也。
SSF: 从心,
PIF: 隹聲。

is 'to think generally'.
It has xīn ``heart'' as a semantic constituent,
zhuī (GSR 575a: ``a kind of dove, ...") is the phonetic constituent.



G The use of fán 凡 to indicate generality is striking and not attested for verbs elsewhere in Shuōwén. It is as if the technical term fán 凡 can do for verbs what the expression N 之總稱 ``is a general designation for N'' does for nouns.
PIF    隹[職追切; LH tshui, OCM tui] phonetic in 惟 [以追切; LH wi, OCM wi]
PR     F*i sh*ng 非聲. We find not only unrelated initial consonants, but also different rhymes.

60. SW 10B 408: 060; DXB 217 (10B x); GL vol 11: 4686; Duàn 505 (10B 30b); TKJ 1448; Ozaki vol. 5: 1018.


HG:
G: SSF: PIF:

¶ (懷), 念思也。从心,
褱聲。


huái,
is 'to think about persistently'.
It has xīn ``heart'' as a semantic constituent,
huái (GSR 600a: ``bosom, ...") is the phonetic constituent.

G The opposition fánsī 凡思 niànsī 念思 is a neat instance of synonym distinction in Shuōwén. Duàn noticed the sustained analytic approach Xǔ Shèn takes to the semantic field of 思.
    The old Wénxu(n commentary (Liù Chén zhù Wénxu(n 六臣注文選) quotes Shuōwén 說文曰：懷，藏也 ``huái is to hide'' as well as 說文曰：懷，念思也 ``huái is to think about persis- tently'' (Gǔlín ibidem).
PIF     褱[戶乖切; LH gu/i, OCM grûi] phonetic in 懷 [戶乖切; LH gu/i, OCM grûi]
    In many cases the history of the Chinese writing system was such that a word that came to be written with a radical was originally only written with the ``phonetic'' itself, a point justly emphasised in the work of William Boltz 1994 and probably even more common in unpreserved epigraphy than preserved epigraphy allows us to see today.
PR   Tóngyīn 同音 OCM and LH are homophones.

61. SW 10B 408: 061; DXB 218 (10B 13a); GL vol 11: 4686b; Duàn 505 (10B 31a); TKJ 1448; Ozaki vol. 5: 1019.
HG:   U (惀),	lún,
G:	欲知之皃。 descriptive of the desire to understand.

SSF: PIF:

从心, 侖聲。

It has xīn ``heart'' as a semantic constituent,
lún (GSR 470a: ``Shuowen says: to think, ponder, ...") is the phonetic constituent.

HG This entry is an intrusion into the ``think'' series.
G	The ``desire to understand", close to the Greek philosophia, was not in fact a current notion in ancient Chinese.
PIF     侖[力屯切; LH luin, OCM run] phonetic in 惀 [盧昆切; LH luin, OCM run]
PR	Tóngyīn 同音 OCM and LH are homophones.



62. SW 10B 408: 062; DXB 218 (10B 13a); GL vol 11: 4686b; Duàn 505 (10B 31a); TKJ 1449; Ozaki vol. 5: 1019.


HG:
G: SSF: PIF:

e (想), 冀思也。从心,
相聲。


xi&ng,
is 'to look forward to and think of'.
It has xīn ``heart'' as a semantic constituent,
xi+ng (GSR 731a: ``look at, ...") is the phonetic constituent.

G Xǔ Shèn resumes his sī 思 synonym series after an interruption in the preceding item which is about ``understanding". The notion of understanding is close enough to that of thinking, but it is not explicitly described in terms of 思.
Duàn rewrites jì 冀 as jì 覬. This is because, as we noted before, he holds the mistaken
opinion that Xǔ Shèn must use characters in meanings he attributes to them in his dictionary. However, his definition of jì 冀 is 北方州也. Whereas jì 覬 is glossed as meaning ``look forward to".
PIF    相 [息良切; LH siâng, OCM sang] phonetic in 想 [息兩切; LH siângB, OCM sang)]

63. SW 10B 408: 063; DXB 218 (10B 13b); GL vol 11: 4687; Duàn 505 (10B 31a); TKJ 1449; Ozaki vol. 5: 1019.


HG:
G: SSF: PIF:

(徻),
深也。从心, 嬕聲。

suì,
is (figuratively ?) 'deep'.
It has xīn ``heart'' as a semantic constituent,
suì (SW: ``follow one's thoughts, ...") is the phonetic con- stituent.

HG This word is almost only attested in dictionaries. Words of this sort are uncomfortably common in many sections of Shu\$ wén, and they do raise serious issues of methodology. Xǔ Shèn must have had a variety of sources on which to base himself in writing about these characters, but few of his lexicographic sources have come down to us. Moreover, none of these seem to have left any trace in the received literature. Striking series of cases of this sort dominate in the sections on salt  and  on  horses.  Compare  the  pioneering  M1  Z\%nghuò  馬宗霍 (1959)  Shu\$  wén  ji\%  zì  y+n qúnshū k(o 說文解字引群書考.
PIF     嬕 [徐醉切; LH zuis, OCM s-wi(t)s] phonetic in 徻 [徐醉切; LH zuis, OCM s-wi(t)s]
***

64. SW 10B 408: 064; DXB 218 (10B 13b); GL vol 11: 4687; Duàn 505 (10B 31a); TKJ 1449; Ozaki vol. 5: 1020.
HG: ‚ (慉),	xù, ``support"




G:	起也。
SSF: 从心,

is 'to make someone rise > grow up'. [[is '(a way of) making someone rise > grow up (by providing proper loving support).]] It has xīn ``heart'' as a semantic constituent,

PIF: IQ:

畜聲。
《詩》曰： 能不我慉。

xù (GSR 1018a: ``...nourish, ...") is the phonetic constituent.
The Book of Odes (SSJZS: 304c) says:
"It is possible that you do not cherish me.'' (tr. Karlgren)

G	The addition of the heart radical here does, very clearly, indicate an internalised psychologis- ing nuance in the word in question.
PIF     xù 畜 [丑六切; LH h\$uk (or huk), OCM huk] phonetic in xù 慉 [許六切; LH h\$uk (or huk), OCM huk]
The phonetic xù 畜 is also a semantic constituent if ever there was one in any character.
PR	Tóngyīn 同音 OCM and LH are homophones.
IQ	The SSJZS edition of the Shījīng (Gǔf*ng 谷風 35.5) has: 不我能慉, ``You could not cherish me".


65. SW 10B 408: 065; DXB 218 (10B 13b); GL vol 11: 4688; Duàn 505 (10B 31a); TKJ 1449; Ozaki vol. 5: 1020.


HG:
G: SSF:
PIF:
CAS:

(⨈),
滿也。从心,
ོ聲。一曰:

yì,
is 'full'.
It has xīn ``heart'' as a semantic constituent,
yì (SW: ``be in high spirits, ...") is the phonetic constituent. An (alternative) source says:

十萬曰⨈。 A hundred thousand are called yì.

SA:
    
[ 悥] ， 籀文省。

is a large seal style graph and it is abbreviated.

G Fullness or filling something up apparently has nothing to do with the heart or mind. One wonders why this graph was written with the heart semantic constituent.
PIF The graph 訲 is taken to be equivalent in HYDZD: 3946 to ོ Shuōwén 3A 56.021.
ོ [於力切; LH sheC, OCM lheh ?] phonetic in ⨈ [於力切; LH )\$\#k OCM: unattested]
CAS Having given his explanation, Xǔ Shèn goes on to quote an entirely separate interpretation which in this case does seem to exclude any psychologising extended meaning. Alternative sources in Shuōwén are of at least three types:
1. corroborative ways of making basically the same semantic interpretation.
2. competitive ways of interpreting the graph in a way that allows one to understand why the semantic constituent is present.




3. alternative ways of interpreting the graph in a way that seems irrelevant to the presence of the semantic constituent.

66. SW 10B 408: 066; DXB 217 (10B x); GL vol 11: 4689b; Duàn 505 (10B 31b); TKJ 1449; Ozaki vol. 5: 1021.


HG:
G: SSF: PIF:

R (悹),
憂也。从心, 官聲。


guàn,
is 'to worry'.
It has xīn ``heart'' as a semantic constituent,
gu+n (GSR 157a: ``office, ...") is the phonetic constituent.

G Duàn rewrites 憂 as ⤏. As we have noticed under SW 10B 408: 062 (想) for jì 冀 Duàn labours under the misconception that Xǔ Shèn must have used characters in the meanings he attributes to them. Since the ordinary y\$u 憂 is not glossed in Shuōwén as ``to worry", Duàn replaces it throughout the text with the other y\$u ⤏ which is so glossed. On this matter Duàn is at least consistent, though surely profoundly mistaken.
PIF  Bureaucrats would surely insist that the graph 官 is a semantic constituent!
Gu&n 官 [古丸切; LH kuân, OCM kôn] phonetic in guàn 悹 [古玩切; LH kuânC, OCM kôns]

67. SW 10B 408: 067; DXB 218 (10B 13b); GL vol 11: 4690/10353; Duàn 505 (10B 31b); TKJ 1450; Ozaki vol. 5: 1021.


HG:
G: SSF: PIF:

• (憀), 憀然也。从心,
翏聲。


liáo,
is (as in) 'with clear understanding'.
It has xīn ``heart'' as a semantic constituent,
liù (GSR 1069a: ``whistling of the wind (Chuang), ...") is the phonetic constituent.

G What, one wonders, is the difference between this gloss and liáo mào 憀皃 ``descriptive for"? The answer is that the present formula indicates the pattern in which the descriptive word tends to occur. Thus the gloss does not give the meaning of the character but a phrase in which it is idiomatic.
    Overtly circular glosses like this one are fairly rare in Shuōwén, and they are of course not very informative. Xǔ Shèn assumes that we know what the word means and specifies what kind of construction it typically enters.
PIF  䟠 [力救切; LH liuC, OCM riuh] phonetic in 憀 [洛蕭切; LH leu, OCM riû]



68. SW 10B 408: 068; DXB 218 (10B 13b); GL vol 11: 4690b; Duàn 505 (10B 31b); TKJ 1450; Ozaki vol. 5: 1022.



HG:
G: SSF: PIF: IQ:

(䓨),
敬也。从心, 客聲。
《春秋傳》曰：
以陳備三䓨。

kè,
is 'to respect'.
It has xīn ``heart'' as a semantic constituent,
kè (GSR 766d': ``guest") is the phonetic constituent.
The Springs and Autumns Tradition (i.e. Zuǒzhuàn, SSJZS: 1985b) says:
"With Chen he made up the number of the three ``re- spected entities"."

HG     臣鉉等曰：今俗作恪。 ``Xuán and others comment: today the current form is 恪."
G	The number of heart radical terms glossed by jìng 敬 is remarkable: we count 7.
PIF     Duàn suggests that the text should read 从心客,客亦聲, as if Xǔ Shèn did not often disregard the concurrent semantic functions of phonetic constituents.
客 [苦格切; LH khak, OCM khrâk] phonetic in 䓨 [苦各切; LH khâk, OCM khâk]
IQ The SSJZS edition of Zuǒzuàn (Xiāng 襄 25.10) writes: 庸以元女大姬配胡公，而封諸陳， 以備三恪。"… gave his own oldest daughter, T'ae-ke, in marriage to (his son), duke Hoo, and invested him with Ch'in, thus completing the number of the 'three honoured States'.'' (tr. Legge).


***
Minor FEAR SERIES 1 (69-70)
[Note that there is another more extended FEAR SERIES in n° 239 below.]

69. SW 10B 408: 069; DXB 218 (10B 13b); GL vol 11: 4693; Duàn 506 (10B 32a); TKJ 1450; Ozaki vol. 5: 1023.

HG:  w (愯),
G:	懼也。
SSF: 从心,
PIF: 雙省聲。


sǒng,
is 'to fear'.
It has xīn ``heart'' as a semantic constituent,
It has an abbreviated form of shu+ng (GSR 1200a:"a pair") as the phonetic.

IQ:

《春秋傳》曰： The Springs and Autumns Tradition (i.e. Zuǒzhuàn, SSJZS: 2087c) says:

駟氏愯。

"The Sì clan got frightened."

PIF     雙[所江切; LH srcng, OCM srông] phonetic in 愯 [息拱切; LH siongB, OCM song)]



IQ	The SSJZS edition of the Zuǒzhuàn (Zhāo 昭 19.8) reads: 駟氏聳 ``hereby alarming the Sze family.'' (tr. Legge)
    Here as elsewhere, it appears that Xǔ Shèn describes a use of radicals in characters to write words that is significantly more systematic than that in our received texts.

70. SW 10B 408: 070; DXB 218 (10B 13b); GL vol 11: 4694/10361; Duàn 506 (10B 32a); TKJ 1450; Ozaki vol. 5:
1023.


HG:
G:

SSF: PIF: SA:

· (懼),
恐也。
从心, 瞿聲。
[愳], 古文


jù,
is 'to be terrified'. [[is '(a way of) being terrified (scil. but relatively mildly).]]
It has xīn ``heart'' as a semantic constituent,
jù (GSR 96c: ``anxious, ...") is the phonetic constituent.  is an ancient style graph.

G	Jù 懼 and kǒng 恐 are hūzhù 互註 characters in Shuōwén without being synonyms.
PIF     瞿[九遇切。又音衢; LH kyâC, OCM kwah] phonetic in 懼 [其遇切; LH gyâC, OCM gwah]
***

Minor DEPENDENCE SERIES (71-72)

71. SW 10B 408: 071; DXB 218 (10B 13b); GL vol 11: 4694b; Duàn 506 (10B 32a); TKJ 1451; Ozaki vol. 5: 1023.


HG:
G: SSF: PIF:

  (怙),
恃也。从心, 古聲。


hù,
is 'to rely on'.
It has xīn ``heart'' as a semantic constituent,
gǔ (GSR 49a: ``ancient, ...") is the phonetic constituent.

PIF     古[公戶切; LH kâB, OCM kâ)] phonetic in 怙 [矦古切; LH gâB, OCM gâ)]

72. SW 10B 408: 072; DXB 219 (10B 14a); GL vol 11: 4695; Duàn 506 (10B 32a); TKJ 1451; Ozaki vol. 5: 1024.


HG:
G: SSF: PIF:

, (恃),
賴也。从心, 寺聲。


shì,
is 'to depend on'.
It has xīn ``heart'' as a semantic constituent,
sì (GSR 961m: ``hall, ...") is the phonetic constituent.



G Xǔ Shèn juxtaposes two near synonyms and glosses them not by the same character, as often elsewhere, but he decides to gloss the first character by the second and then to provide an near synonym gloss for the second character in this mini series. One would like to know why he decides against hùzhù 互注 in cases like these. Conceivably, part of the motivation is that he treats binomes always in the order of their occurrence in the phrase and attaches an external gloss only to the second member of the binome.
PIF     寺[祥吏切; LH zi\#C, OCM s-l\#h ?] phonetic in 恃 [時止切; LH gâB, OCM gâ)]
***

73. SW 10B 408: 073; DXB 219 (10B 14a); GL vol 11: 4695; Duàn 506 (10B 32a); TKJ 1451; Ozaki vol. 5: 1024.


HG:
G: SSF: PIF:

‡ (慒), 慮也。 从心,
曹聲。


cóng,
is 'to make plans'.
It has xīn ``heart'' as a semantic constituent,
cáo (GSR 1053a: `` servant, ...") is the phonetic constituent.

PIF     曹[昨牢切; LH dzou, OCM dzû] phonetic in 慒 [藏宗切; LH dzoung, OCM dzûng]
PR	F*i sh*ng 非聲.

74. SW 10B 408: 074; DXB 219 (10B 14a); GL vol 11: 4695b; Duàn 506 (10B 32b); TKJ 1451; Ozaki vol. 5: 1024.


HG:
G:

SSF: PIF: SA:

I (悟),
覺也。
从心, 吾聲。
[	]，古文悟。


wù,
is 'to become aware'. [[EP: ``is (a way of) becoming aware (scil. suddenly or abruptly)'.]]
It has xīn ``heart'' as a semantic constituent,
wú (GSR 58f: ``I, ...") is the phonetic constituent. is an ancient style graph for wù.

G Jué 覺 is in turn glossed as 寤也, and one notes that both these psychological terms are written with heart-less graphs. For wù 寤 we have: 寐覺而有信曰寤 'to wake up from sleep so as to come to have beliefs>conscience is called wù'. Xǔ Shèn's glossing system comes close to the application of hùzhù 互注 to the members of a binome.
PIF    吾[五乎切; LH ngâ, OCM ngâ] phonetic in 悟 [五故切; LH ngâC, OCM ngâh]



75. SW 10B 408: 075; DXB 219 (10B 14a); GL vol 11: 4695b; Duàn 506 (10B 32b); TKJ 1451; Ozaki vol. 5: 1025.


HG:
G:
DG:
AG:

£ (憮),
愛也。
韓鄭曰憮。: 一曰:
不動。


wǔ, ``cherish"
is 'to show loving care'.
In Hán and in Zhèng they use the word wǔ. An (alternative) source says:
not move.

SSF: 从心,
PIF: 無聲。

It has xīn ``heart'' as a semantic constituent,
wú (GSR 103a: ``not have, ...") is the phonetic constituent.

G Compare the word fǔ 撫 ``caress'' which corresponds nicely as a manual pendant to the psychological wǔ 憮.
    Duàn rewrites this, here as elsewhere, as 形也, as if Xǔ Shèn was not entitled to use the graph 愛 in any other meaning than that indicated by the gloss which is part of the graphological analysis.
DG This gloss reccurs in F&ngyán 方言. The systematic relation between F&ngyán and Shuōwén needs exploration: if Xǔ Shèn was working with a copy of F&ngyán, one wonders how he decided in favour of the relatively few dialectal glosses he quotes from the book. See M1 Z\%nghuò, Shu\$ wén ji\% zì y+n F&ngyán k(o (1959,ch 3: 26).
    Immitating the style of the F&ngyán, Xǔ Shèn offers dialect information instead of a gloss. Dialectal glosses like these are concerned with the usage of words rather than graphs. The highly interesting field of graphic dialectology usefully developped by Lǐ Pǔ 李圃, does not enter the purview of Xǔ Shèn's analysis. When Xǔ Shèn offers dialectal glosses we must assume that the semantic gloss embedded in the dialectal gloss contains the analytically relevant semantic information.
AG	This gloss looks implausible and it has even been suggested to rewrite the text with xīn 心
instead of bù 不.
PIF     無[武扶切; LH muâ, OCM ma] phonetic in 憮 [文甫切; LH muâB, OCM ma)]

76. SW 10B 408: 076; DXB 219 (10B 14a); GL vol 11: 4696b; Duàn 506 (10B 32b); TKJ 1451; Ozaki vol. 5: 1025.


HG:	(形),
G:	惠也。
SSF: 从心,
PIF: 先聲。

ài,
is 'generous'. [[is 'a (way of being) generous towards others (scil. in the loving mode).]]
It has xīn ``heart'' as a semantic constituent,
xi+n (GSR 478a: ``preceding, ...") is the phonetic con- stituent.



SA:    [⦯]，古文。  is an ancient style graph.
HG    On the basis of his graphological analysis Xǔ Shèn seems to have arrived at the conclusion that the standard graph for ài 愛 'love' is in fact a phonetic loan. See TKJ: 728.
G Xǔ Shèn's gloss identifies generosity as the hypernym of love, and in so doing he makes an important point: loving care without the element of free generosity is not love. For example, parental love for their children does not really count as love when it is based not on generosity but on expectation of future reward. One hastens to add that Xǔ Shèn's main purpose in glossing his characters was not this kind of conceptual analysis.
PIF Duàn tacitly rewrites this phonetic gloss by providing a seal graph version of xi&n 先. This represents a distortion of Xǔ Shèn's text because Xǔ Shèn's analysis of graphs is not graphic but graphemic. Xǔ Shèn is concerned not with the graph that is phonetic but with the grapheme and he regularly identifies his graphemes in ways that are more abstract than writing out the seal version of a graph that is physically at issue.
先[穌前切; LH k\$s, OCM k\#ts] phonetic in 形 [烏代切; LH )\#s, OCM )*+ ts]

77. SW 10B 408: 077; DXB 219 (10B 14a); GL vol 11: 4697b; Duàn 506 (10B 32b); TKJ 1452; Ozaki vol. 5: 1026.


HG:	(徽),
G:	知也。
SSF: 从心,
PIF: 胥聲。

xǔ,
is 'to understand'. [[is '(a way of having) understanding (scil. as a stable ability or talent).]]
It has xīn ``heart'' as a semantic constituent,
xū (GSR 90e: ``together, ...") is the phonetic constituent.

PIF     胥[相居切; LH sia, OCM sa] phonetic in 徽 [私呂切; LH siaB, OCM sa)]

78. SW 10B 408: 078; DXB 219 (10B 14a); GL vol 11: 4697b; Duàn 506 (10B 32b); TKJ 1452; Ozaki vol. 5: 1026.


HG:
G:

SSF: PIF:

 (慰),
安也。
从心, 尉聲。


wèi, ``to console"
is '(to make) peaceful'. [[EP: (a way of) making peaceful (scil. by providing comfort)'.]]
It has xīn ``heart'' as a semantic constituent,
wèi (GSR 525b: ``comptroller, ...") is the phonetic constituent.

CAS: 一曰:
恚怒也。

An (alternative) source says: 'it is anger'.

PIF    Duàn tacitly and gratuitously rewrites the standard graph to take the seal form, as if there was a convention in Shuōwén to the effect that phonetic constituents are listed in their seal form. See Shuōwén 10B 408: 076; ài 形.



尉[於胃切; LH ?us, OCM uts] phonetic in 慰 [於胃切; LH ?us, OCM uts]
PR	Tóngyīn 同音 OCM and LH are homophones.
CAS The current binome in the alternative gloss gives a separate meaning which seems unattested elsewhere in the literature.

79. SW 10B 408: 079; DXB 219 (10B 14a); GL vol 11: 4698; Duàn 506 (10B 33a); TKJ 1452; Ozaki vol. 5: 1026.


HG:
G: SSF: PIF: SPI:

(⨍),
謹也。从心, പ聲。
讀若毳。

cuì,
is 'diligent'.
It has xīn ``heart'' as a semantic constituent,
zhuì (SW: ``to divine, ...") is the phonetic constituent. Pronounced like cuì.

PIF     zhuì പ[之芮切; LH suis, OCM suts] phonetic in cuì ⨍ [此芮切; LH tshyâs, OCM tshots]
PR	This marginal case of a f*i sh*ng 非聲 had the same rhymes according to MC reconstructions, but if the OC and LH reconstructions were right the vowels were different.
SPI      Cuì ⨍ [此芮切; LH tshyâs, OCM tshots] read like cuì 毳 [此芮切; LH tshyas, OCM tshots!]

80. SW 10B 408: 080; DXB 219 (10B 14a); GL vol 11: 4698b; Duàn 506 (10B 33a); TKJ 1452; Ozaki vol. 5: 1027.


HG:
(Š),
chóu,
G:
Š箸也。
is (as in) 'to hesitate'.
SSF:
从心,
It has xīn ``heart'' as a semantic constituent,
PIF:
聲。
chóu (GSR (1090m: 籌): ``counting stick...") is the pho- netic constituent.
G This is another way of writing chóuchú 躊躇 ``hesitate". The definition of a character by the typical synonym binome in which it occurs is a standard procedure in Shuōwén. (The graphic shape of liánmiánzì 連綿字 is notoriously unstable in classical Chinese orthography, a extreme example of this instability being the word w*iyí 逶迤 for which we have dozens of recorded variant graphs. The question remains whether in this case Xǔ Shèn is intending to indicate, by the format he uses for his definition, that the word Š can be used independently to mean the same as Š箸.
PIF The phonetic constituent is not a Head graph in Shuōwén.
      [直由切; LH dru, OCM dru ?]) phonetic in Š [直由切; LH dru, OMC dru] (see the allograph ʱ for 疇 TKJ: 1983, or 籌 TKJ: 640).
The phonetic constituent is homophonous with the character.



81. SW 10B 408: 081; DXB 219 (10B 14a); GL vol 11: 4699; Duàn 506 (10B 33a); TKJ 1452; Ozaki vol. 5: 1027.


HG:
G: SSF: PIF: IQ:

  (怞),
朗(=悢?)也。从心,
由聲。
《詩》曰： 憂心且怞。


chóu,
is 'resentment' (?).
It has xīn ``heart'' as a semantic constituent,
yóu (GSR 1079a: ``from, ...") is the phonetic constituent.
The Book of Odes (SSJZS: 466c) says:
"I am worried in my heart and resentful ??."

G    Xǔ Shèn's definition is at variance with early glosses.
    Duàn miswrites the character by inverting the elements, and then goes on to suggest the emendation hèn 恨 for the character he has just misrepresented. His conjecture is tempting nontheless.
IQ	The SSJZS edition of Shījīng (Gǔ zh\$ng 鼓鐘 208.3) has the woman radical in our graph: 憂心且䭬 ``I am worried in my heart and agitated.'' (tr. Karlgren).
PIF     由[以周切; LH ju, OCM liu] phonetic in 怞 [直又切 (《廣韻》: 直由切); dru, OCM d-liu ?]
PR   Note the non-homogenous initials.


82. SW 10B 408: 082; DXB 219 (10B 14b); GL vol 11: 4700; Duàn 506 (10B 33a); TKJ 1453; Ozaki vol. 5: 1027, see
variants in WGY: 451.


HG:
G: SSF: PIF:

SPI:

(⣳),
⣳撫也。从心,
某聲。
讀若侮。

wǔ,
is (as in)'to fondle affectionately'.
It has xīn ``heart'' as a semantic constituent,
mǒu (GSR 948a: ``a certain ...", SW: méi 'prune') is the pho- netic constituent.
Pronounced like wǔ.

G Duàn rewrites 撫 as 憮 and proposes to omit the repeated ⣳. Duàn would have to omit all the repeated definienda if he wanted to be consistent, and even that would be unacceptable because it destroys a possibly significant trace of the sources for the compilation of Shuōwén.
PIF In his entry for mǒu 某 (Shuōwén 6A 206:146, TKJ: 767, méi), Xǔ Shèn's gloss shows that he derives the graph from the meaning méi 梅 'prune' and his reading of the character was manifestly méi (謨杯切). However, when this constituent enters into other graphs as a phonetic it appears that its value can be, as in the present case, mǒu (莫厚切) rather than méi (謨杯切). Note that the phonetic 某 functions in méi 媒 as well as in móu 謀.
    One given phonetic constituent can have several distinct values when used as a phonetic con- stituent in different characters.



    某 [莫厚切（《集韻》: 謨杯切）; LH moB, OCM mô)] phonetic in ⣳ [亡甫切; LH muoB, OCM mo)]
SPI     Wǔ ⣳ [亡甫切; LH muoB, OCM mo)] read like wǔ 侮 [文甫切; LH muoB, OCM mo)]

***
EFFORT SERIES (83-88)

83. SW 10B 408: 083; DXB 219 (10B 14b); GL vol 11: 4701; Duàn 506 (10B 33a); TKJ 1453; Ozaki vol. 5: 1028, see
variants in WGY: 451.

HG: ✭ (忞),
G:	彊也。
SSF: 从心,
PIF: 文聲。


mín,
is 'strong'. [[is '(a way of making) strong (scil. in a figurative sense) effort.]]
It has xīn ``heart'' as a semantic constituent,
wén (GSR 475a: ``ornate, ...") is the phonetic constituent.

IQ:

《周書》曰： In The book of Zhǒu (SSJZS: 231a) it says:


SPI:

在受德忞。讀若旻。

"When it came to Shòu, his virtue was powerful.'' pronounced like mín

G Duàn congenially but also wantonly expands this gloss to 自勉彊也 on the late basis of a Yùn huì 韻會 quotation.
IQ The SSJZS edition of the Shūjīng 書經 (Lì zhèng 立政 47.5) writes: 鳴呼。其在受。德䯭惟羞刑暴德之人同于厥邦 ``Oh, (when it rested with Shòu=) when the turn came to Shòu, his character was impetuous; (shamed=) disgraced (punished ones=) criminals and men of a violent character were his associates in his state.'' (tr. Karlgren)
PIF   文[無分切; LH mun, OCM m\#n] phonetic in 忞 [武巾切; LH m\$n, OCM mr\#n]
SPI The present dúruò 讀若 gloss provides no unexpected or new information, one is tempted to speculate that in the preceding entry Xǔ Shèn has got into a temporary habit of adding a dúruò 讀若 gloss.
Mín 忞 [武巾切; LH m\$n, OCM mr\#n] read like mín 旻 [武巾切; LH m\$n, OCM mr\#n]

84. SW 10B 408: 084; DXB 219 (10B 14a); GL vol 11: 4701b; Duàn 506 (10B 33b); TKJ 1453; Ozaki vol. 5: 1028.
HG:  \$(慔), mù,

G:	勉也。
SSF: 从心,

is 'to make an earnest effort'.
It has xīn ``heart'' as a semantic constituent,



PIF: 莫聲。	mù (GSR 802a: ``evening, ...") is the phonetic constituent.
G	Xǔ Shèn is taking over the gloss from )ry( 爾雅 (Shì xùn 釋訓 3.28, Xú Cháohuá 徐朝華
1994: 136: 懋懋, 慔慔,勉也).
PIF     Duàn gratuitously rewrites the phonetic with 廾 instead of 大: 温.
    莫 [莫故切，又，慕各切; LH mâh, OCM mâkh] phonetic in 慔 [莫故切; LH mâh, OCM mâkh]
PR	Tóngyīn 同音 OCM and LH are homophones.

85. SW 10B 408: 085; DXB 219 (10B 14a); GL vol 11: 4701b; Duàn 506 (10B 33b); TKJ 1453; Ozaki vol. 5: 1028.
HG: q (䓧), mi&n,

G:	勉也。
SSF: 从心,
PIF: 面聲。

is 'to be make a strenuous effort'. [[is '(a way of making) a stre- nous effort (scil. by way of self-reflection).]]
It has xīn ``heart'' as a semantic constituent,
miàn (GSR 223a: ``face...") is the phonetic constituent.

G	Shuōwén's gloss looks strange unless one brings to bear on it an explanation such as the conjectural one proposed in our explanatory paraphrase.
PIF     面[彌箭切; LH mianC or mienC, OCM mens] phonetic in 䓧 [弥殄切; LH mianB or mienB, OCM men)]


86. SW 10B 408: 086; DXB 219 (10B 14a); GL vol 11: 4702; Duàn 506 (10B 33b); TKJ 1453; Ozaki vol. 5: 1029.


HG:
G: SSF: PIF:

(⠂),
習也。从心, 曳聲。

yì,
is 'to carry into (diligent) practice'.
It has xīn ``heart'' as a semantic constituent,
yè (GSR 338a: ``to drag, ...") is the phonetic constituent.

HG Duàn argues that Shuōwén must have had a different head graph here, he proposes: 忕.
G	Here, as often elsewhere, an irrelevant lexical entry intervenes in a semantic series. One can only speculate about the reasons for this irregularity.
PIF Duàn rewrites the phonetic to suit his initial conjecture. In effect he takes it upon himself to rewrite most of the present lexical entry!
曳 [余制切; LH jas, OCM lats or jats] phonetic in ⠂ [余制切; LH jas, OCM lats or jats]
PR	Tóngyīn 同音 OCM and LH are homophones.



87. SW 10B 408: 087; DXB 219 (10B 14a); GL vol 11: 4703b; Duàn 507 (10B 34a); TKJ 1453; Ozaki vol. 5: 1030.

HG:  ¯ (懋),
G:	勉也。
SSF: 从心,


mào,
is 'to make strenuous effort'.
It has xīn ``heart'' as a semantic constituent,

PIF: IQ:
A:

楙聲。
《虞書》曰： 時惟懋哉。
 [※]，或省.

mào (GSR 1109d: ``Shuowen says: winter peach tree") is the phonetic constituent.
The Book of Yú (SSJZS: 130b) says:
"In this be energetic."
When written   (the word mào) is alternatively written in an abbreviated way.

G	The old Wénxu(n commentary (Liù Chén zhù Wénxu(n 六臣注文選 1987: ju(n 10, 190) quotes: 說文曰：懋，盛懋 ``is abundant".
    Compare )ry( (Shì xùn 釋訓 3.28, Xú Cháohuá 徐朝華 1994: 136): 懋懋, 慔慔,勉也. It is probably not by coincidence that the gloss mù 慔 of the )ry( has been defined in n° 84 above.
PIF     楙 [莫候切(《廣韻》: 莫袍切); LH mouC, OCM mûh] phonetic in 懋 [莫候切; LH mouC, OCM mûh]
IQ The Shùndi(n (SSJZS: 130b) has: 汝平水土惟時懋哉 ``you shall regulate water and land, in this be energetic!'' (tr. Karlgren). Shàngshū wénzì jiào gǔ 尚書文字校詁 (ed. Zàng Kèhé, 1999: 58) has no alternative reading and has no comment on the character.
PR   Tóngyīn 同音 OCM and LH are homophones.

88. SW 10B 408: 088; DXB 219 (10B 14b); GL vol 11: 4704; Duàn 507 (10B 34a); TKJ 1454; Ozaki vol. 5: 1030.
HG: Š (⧂;慕), mù,

G:	習也。
SSF: 从心,
PIF: 莫聲。

is 'to put into practice' [[is '(a way of) engaging in a practice (scil. so as to emulate what one admires)'.]]
It has xīn ``heart'' as a semantic constituent,
mù (sic) (GSR 802a: ``evening, ...") is the phonetic constituent.

PIF     莫 [莫故切，又，慕各切; LH mâC, OCM mâkh] phonetic in 慕 [莫故切; LH mâC, OCM mâkh]
PR	Tóngyīn 同音 OCM and LH are homophones.
***



89. SW 10B 408: 089; DXB 219 (10B 14b); GL vol 11: 4704b; Duàn 507 (10B 34a); TKJ 1454; Ozaki vol. 5: 1030.
HG: G (悛), qu+n,

G:	止也。
SSF: 从心,

is 'to cease'. [[is '(a way of) ceasing to do something (scil. because one regrets what one has done before)]].
It has xīn ``heart'' as a semantic constituent,

PIF:

夋聲。

qūn (GSR 468p: ``SW says 'to run' and 'to squat' (no text)") is the phonetic constituent.

PIF     夋[七倫切; LH tshuin, OCM tshun ?] phonetic in 悛 [此緣切; LH tshyan, OCM tshon]
PR	Note the vowel difference. Here we have a case where both the initials and the finals are the same, but where the crucial main vowel is different. F*i sh*ng 非聲.

90. SW 10B 408: 090; DXB 219 (10B 14b); GL vol 11: 4704b-5; Duàn 507 (10B 34a); TKJ 1454; Ozaki vol. 5: 1031.


HG:
G: SSF: PIF:

(徣),
肆也。从心, 隶聲。

tuì,
is 'to give free rein '.
It has xīn ``heart'' as a semantic constituent,
dài (GSR 509a: ``SW says: reach to, thus taking it to be the primary form of 逮, ...") is the phonetic constituent.

PIF     隶[徒耐切; LH d\#s or j\#s, OCM l*+ s or l\#s] phonetic in 徣 [他骨切（《廣韻》: 他内切); LH thu\#s, OCM lhus]
PR	F*i sh*ng 非聲?

91. SW 10B 408: 091; DXB 219 (10B 14b); GL vol 11: 4705; Duàn 507 (10B 34a); TKJ 1454; Ozaki vol. 5: 1031.


HG:	(忸),
G:	趣步
忸忸也。
SSF: 从心,
PIF: 與聲。

yǔ,
walk quickly>be in a hurry,
is as in (the reduplicated form) yǔyǔ.
It has xīn ``heart'' as a semantic constituent,
yǔ (GSR89b: ``associate with, ...") is the phonetic constituent.

G	Walking quickly seems unrelated to the heart or mind. But note that being in a hurry would be a psychological state.
    The order of the statements here is confusing: one would expect the reduplicated ``as in"- phrase to precede the gloss.
PIF     與[余呂切; LH jâB, OCM la)] phonetic in 忸 [余呂切; LH jâB, OCM la)]



PR	Tóngyīn 同音 OCM and LH homophones.

92. SW 10B 408: 092; DXB 219 (10B 15a); GL vol 11: 4705b; Duàn 507 (10B 34a); TKJ 1454; Ozaki vol. 5: 1031.


HG:
G: SSF: PIF:

 (䓭), 說也。 从心,
䴊聲。


t+o,
is 'to be pleased'.
It has xīn ``heart'' as a semantic constituent,
y&o (GSR 1078a: ``to scoop, ...") is the phonetic constituent.

G Note that Xǔ Shèn writes shu\$ 說 for yuè 悅 in spite of the fact that he glosses shu\$ 說 in terms of explaining and unravelling. This is a neat reminder of the fact that Xǔ Shèn feels entitled to use graphs with readings and interpretations other that those he gives in his own dictionary. Duàn would have avoided many of his less fortunate emendations if he had recognised this elementary point.
PIF    䴊[以沼切; LH jau, OCM lau] phonetic in 䓭 [土刀切; LH thâu, OCM lhâu]

          *** PEACEFUL SERIES (93-95)

93. SW 10B 408: 093; DXB 219 (10B 15a); GL vol 11: 4705b; Duàn 507 (10B 34a); TKJ 1454; Ozaki vol. 5: 1032.


HG:   (䓶),
G:	安也。
SSF: 从心,

y+n,
is 'peaceful'. [[is '(scil. poetic/archaic for being) peaceful (scil. in one's mind)'.]]
It has xīn ``heart'' as a semantic constituent,

PIF: IQ:

厭聲。
《詩》曰: 䓶䓶夜飲。

y+/ yàn (TKJ 1280) (GSR 616c: ``fed up, ...")is the phonetic constituent.
The Book of Odes (SSJZS: 421b) says: ``Peacefully (we) drink in the night.'' (tr. Karlgren)

G	The gloss is taken from )ry( (Shì xùn 釋訓, Xú Cháohuá 徐朝華 1994: 135, 3.24)
PIF     厭[於輒切。又，一琰切; LH ?iam, OCM ?em] phonetic in 䓶 [於鹽切; LH ?iam, OCM ? em]
IQ	The SSJZS edition of Shījīng (Zhàn lù 湛露 174.1) writes: 厭厭夜飲 ``peacefully we drink in the night.'' (tr. Karlgren).



    Duàn surreptitiously rewrites the character 飲 in an archaising way 抠. It is of course quite possible that he was looking at some Shījīng edition with this graph, but Duàn has no excuse for introducing his preferred character into Xǔ Shèn's text.


94. SW 10B 408: 094; DXB 219 (10B 15a); GL vol 11: 4707; Duàn 507 (10B 34b); TKJ 1455; Ozaki vol. 5: 1032.
HG: ¨ (憺), dàn,

G:	安也。
SSF: 从心,
PIF: 詹聲。

is 'peaceful'. [[is '(a way of being) peaceful (scil. in a superior often metaphysical way)'.]]
It has xīn ``heart'' as a semantic constituent,
zh+n (GSR 619a: ``garrulous, ...") is the phonetic constituent.

HG      The binome dànbó 澹泊 ``live tranquilly'' is attested, for example, in Móuzi 牟子 L+ huò lùn
理惑論22421, etc.
PIF     詹[職廉切; LH tsham, OCM tam] phonetic in 憺 [徒敢切; LH dâmB /C, OCM dâm?/h]

95. SW 10B 408: 095; DXB 219 (10B 15a); GL vol 11: 4707b; Duàn 507 (10B 34b); TKJ 1455; Ozaki vol. 5: 1032.
HG:             (怕), bó,

G:	無爲也。
SSF: 从心,
PIF: 白聲。

is 'to be unassertive'. [[is '(a way of being) unassertive (scil. this leading to a peaceful state of mind)'.]]
It has xīn ``heart'' as a semantic constituent,
bái//bó (GSR 782a: ``...white, ...") is the phonetic constituent.

G   The word Xǔ Shèn defines is more commonly written as bó 泊 and is used by S(m1 Xiāngrú in  his  famous  Z+  xū  fù  子虛賦 in  the  phrase  泊乎無為.  It  appears  that  in  Xǔ  Shèn's  system  of orthography the correct 'spelling' of this word is with the heart radical and not with the water radical.
    Duàn argues that 為 needs to be rewritten as 偽, on the basis of Xǔ Shèn's glosses of these words. This again shows his deep misunderstanding of the nature of the Shuōwén, as if Xǔ Shèn felt obliged or indeed even as much as inclined to use words in the meanings indicated by the glosses in his Shuōwén. If Duàn's methodology were right, Xǔ Shèn would have to use 所 to refer to the hissing sound of an axe.
PIF     白[旁陌切; LH bak, OCM brâk] phonetic in 怕 [匹白切又，葩亞切; LH phak, OCM phrâk]
***



96. SW 10B 408: 096; DXB 219 (10B 15a); GL vol 11: 4708b; Duàn 507 (10B 34b); TKJ 1455; Ozaki vol. 5: 1033.


HG:
G1:

G2: SSF: PIF:

4 (恤),
憂也。
收也。从心, 血聲。


xù,
is 'to be worried'. [[is '(a way of) being worried (scil. empatheti- cally, on behalf of someone)'.]]
It is to take in.
It has xīn ``heart'' as a semantic constituent,
xuè (GSR 410a: ``blood, ...") is the phonetic constituent.

G1   Duàn again miswrites 憂 as ⤏.
G2 The translation of 收也 has to be tentative as is the case with many Xǔ Shèn's glosses lacking context. What makes this case particularly worrying is the fact that the first gloss would seem to be entirely satisfactory and sufficient from a graphological point of view. If Xǔ Shèn had been purely concerned with graphological analysis the second gloss, even if semantically informative, would still have been strictly redundant. Analytically redundant glosses like these do show that Xǔ Shèn's predominant methodology does not prevent him, occasionally, from taking note of semantic glosses that have no direct or indirect relevance to his graphological analysis. Thus, although Xǔ Shèn never begins to aspire to cover anything like the whole semantic range of the meanings of a word, he finds himself tempted, on occasion, to record graphologically irrelevant glosses for rare words.
SSF Duàn rewrites 血 in an archaising form, as if there was a convention in Shuōwén to the effect that phonetic constituents were listed in archaising form.
PIF    血[呼決切; LH huet, OCM hwǐt] phonetic in 恤 [辛聿切; LH suit, OCM swit]

97. SW 10B 408: 097; DXB 219 (10B 15a); GL vol 11: 4709; Duàn 507 (10B 35a); TKJ 1455; Ozaki vol. 5: 1033.


HG:   (䓓),
G:	極也。
SSF: 从心,
PIF: 干聲。

g+n,
is 'extreme (ly concerned ?)'.
It has xīn ``heart'' as a semantic constituent,
g+n (GSR 139a: ``shield, ...") is the phonetic constituent.

G      When a rare word is glossed by a common word with a wide range of meanings, translating the gloss can, as in this case, become a most uncomfortably arbitrary exercise: we have no way of knowing what kind of extremes Xǔ Shèn may have been thinking of.
PIF    干[古寒切; LH kân, OCM kân] phonetic in 䓓 [古寒切; LH kân, OCM kân] (perfect fǎnqiè
match)
PR	Tóngyīn 同音 OCM and LH are homophones.



***
Minor JOYFUL SERIES (98-99)

98. SW 10B 408: 098; DXB 219 (10B 15a); GL vol 11: 4709; Duàn 507 (10B 35a); TKJ 1455; Ozaki vol. 5: 1034.

HG:  ¸ (懽),
G:	喜പ也。
SSF: 从心,

guàn,
is 'to be full of joy and anticipation'.
It has xīn ``heart'' as a semantic constituent,

PIF: IQ:

䧙聲。
《爾雅》曰：

guàn (GSR 158a: ``heron, ...") is the phonetic constituent. The )ry& (SSJZS: 2591a) says:

懽懽、愮愮 憂"Guànguàn and yáoyáo refer to being worried and having no
無告也。	one to expressed one's worry to."
G	Dà Xú b\%n's second graph is very rare and all other early editions of Shuōwén have the other graph 択.
PIF     䧙[工奐切; LH kuânC, OCM kwâns] phonetic in 懽 [古玩切; LH huân, OCM hwân]
IQ	Xǔ Shèn quotes )ry( (Shì xùn 釋訓 3.70) (Xú Cháohuá 徐朝華 1994: 146) for an alternative definition of the reduplicated form of the character he is interested in.
Yáo 愮 is one of the character used in Shuōwén for which there is no head graph.


99. SW 10B 408: 099; DXB 219 (10B 15a); GL vol 11: 4710; Duàn 507 (10B 35a); TKJ 1455; Ozaki vol. 5: 1034.


HG:
G:

(忂),
懽也。

yú,
is 'joyful'.

SNT:

琅邪朱虛有忂亭。 In Zhūxū in the Lángyé district there is this Yú village.

SSF: PIF:

从心, 禺聲。

It has xīn ``heart'' as a semantic constituent,
yù TKJ 1251 (GSR 124a: ``monkey, ...") is the phonetic constituent.

HG Dictionary word, no text or Warring States Inscriptions. Compare yú 愚 'stupid' (n° 123) in which the position of the heart radical differs but the pronunciation is the same, so that this case differs interestingly from the strictly comparable pair yí 怡 (n° 42), dài 怠 (n° 134) in which the pronunciation his differentiated.
G The gloss, here as often elsewhere in Shuōwén, is the head graph of the preceding entry. Even without any textual evidence to support our conjecture we feel entitled to a strong suspicion that in Xǔ Shèn's times there was a binome hu&nyú 懽忂 meaning 'joyful'.



SNT Toponymic notes like this kind show Xǔ Shèn's awareness of the importance of place names in lexicography.
PIF     禺[牛具切; LH ng\$oC , OCM ngoh] phonetic in 忂 [噳倶切; LH ng\$o, OCM ngo]
***

100. SW 10B 408: 100; DXB 219 (10B 15a); GL vol 11: 4710; Duàn 507 (10B 35a); TKJ 1456; Ozaki vol. 5: 1035, see
variants in WGY: 452.


HG:
G:
AG:

y (䮢), 飢餓也。一曰:
憂也。


nì,
is 'to suffer hunger'.
An (alternative) source says:
It is 'trouble'. [[EP: it is (a kind of) trouble (scil. in the form of lack of food)]]

SSF: 从心,

It has xīn ``heart'' as a semantic constituent,

PIF: IQ:

叔聲。
《詩》曰： 䮢如朝飢。

shū (GSR 1031b: ``harvest, ...") is the phonetic constituent.
The Book of Odes (SSJZS: 282c) says: ``I am troubled as if morning-hungry".

G Xǔ Shèn's alternative gloss is the current one according to Z\%ng Fúbāng 宗福邦 et al. (Gùxùn Huìzu(n 故訓匯纂 2003: 799). What motivated Xǔ Shèn's gloss is plainly the Shījīng quotation which he does provide. In this instance the origin of Xǔ Shèn's dictionary in the commentarial tradition concerned with the classics is manifest.
AG Note that on our understanding of this gloss, y\$u 憂 does not have the standard meaning of 'to worry' but rather that of 'having reason to worry, be in trouble'.
Duàn again rewrites 憂 as ⤏.
PIF 叔 [式竹切; LH shuk, OCM nhuk ] phonetic in 䮢 [奴歷切; LH nek < neuk, OCM niûk (After < comes an earlier form, i.e. later LH nek, earlier LH and before probably neuk.)]
IQ The SSJZS edition of Shījīng (Rǔ fén 汝墳 10.1) has 調 instead of 朝: 䮢如調飢 ``I am desirous as if morning-hungry.'' (tr. Karlgren).


101. SW 10B 408: 101; DXB 219 (10B 15a); GL vol 11: 4711b; Duàn 507 (10B 35b); TKJ 1456; Ozaki vol. 5: 1035.


HG:
G: SSF: PIF:

(⣲),
勞也。。从心,
卻聲。

j(,
is 'to work hard'.
It has xīn ``heart'' as a semantic constituent,
què (GSR 776b: ``refuse, ...") is the phonetic constituent.



HG Dictionary word, no text. Duàn declares this to be an allograph of ѱ: ``本一字耳", because the two graphs have the same gloss and the same pronunciation. Xǔ Shèn, meanwhile, is writing a graphological dictionary, and for him there are two graphs to analyse. As modern linguists, we would be inclined to say that the graphs write the same word, but we have many examples of exact homophones with contrasting semantic constituents/radicals representing distinct nuances of meaning: huò 獲 'gather, obtain (generally)', 穫 'gather, harvest grain'; zhù 住 'stay (generally)', 駐'stay (of equestrian armies)'.
    In the absence of a classical Chinese word for 'word', that is more specific and unambiguous then Duàn's zì 字, discussions on what word is referred to by a character can become quite inconvenient, although Duàn's discourse on the matter shows that talking about words is not impossible in classical Chinese. What has to be said about the difficulties surrounding the concept of a word applies mutatis mutandis to the concept of a concept. The history of concepts like that of a word, or that of a language, or indeed that of a concept in traditional China has not to our knowledge been explored so far.
PIF 卻 [去約切; LH kh\$âk, OCM khak] phonetic in ⣲ [其虐切(《廣韻》：奇逆切); LH k\$âk, OCM kak ]


102. SW 10B 408: 102; DXB 219 (10B 15b); GL vol 11: 4711b; Duàn 507 (10B 35b); TKJ 1456; Ozaki vol. 5: 1036.


HG:
G:
IQ:

§ (䮮),	xi+n,
䮮詖也。	is (as in) 'fraudulent pretence'.
䮮利於上， If one fraudulently pretends to profit one's superior one is a

佞人也。
SSF: 从心,
PIF: 僉聲。

glib-tongued person.
It has xīn ``heart'' as a semantic constituent,
qi+n (GSR 613a: ``...all (SHU), ...") is the phonetic constituent.

G	Xǔ Shèn defines the word by a synonym compound including the definiendum and the definition may thus be called circular.
Duàn suggests that 䮮 is a mistake for 險, as if Xǔ Shèn did not regularly define rare
characters by compounds into which they enter.
IQ	Unattributed illustrative quotations of this sort are not very common in Shuōwén. See however n° 4 above (xìng 性).
PIF     僉 [七廉切; LH tshiam, OCM tsham < k-sam] phonetic in 䮮 [息廉切; LH tshiam, siam, OCM tsham < OCM k-sam or sam]



103. SW 10B 408: 103; DXB 219 (10B 15b); GL vol 11: 4712; Duàn 507 (10B 35b); TKJ 1456; Ozaki vol. 5: 1036.


HG:
G: SSF: PIF:

r (愒), 息也。从心,
曷聲。


qì,
is 'to rest'.
It has xīn ``heart'' as a semantic constituent,
hé (GSR 313d: ``what, ...") is the phonetic constituent.

HG It appears that this head graph is Xǔ Shèn's orthographic standard for writing the word currently written qì 憩 today. In any case qì 憩 is not a head graph in Shuōwén.
[臣鉉等曰：今俗䫲作憩，非是。 ``Xuán and the others comment: now a vulgar variant is
qì, it is not correct."]
PIF	曷[胡葛切; LH gât, OCM gât ] phonetic in 愒 [去例切; LH kh\$ât, OCM khat]

104. SW 10B 408: 104; DXB 219 (10B 15b); GL vol 11: 4712b; Duàn 507 (10B 35b); TKJ 1457; Ozaki vol. 5: 1036.


HG:
G: SSF: PIF:

(忩),
精䮴也。从心,
毳聲。

hū,
is 'exquisite simplicity'.
It has xīn ``heart'' as a semantic constituent,
cuì (GSR 345a: ``fine hair, ...") is the phonetic constituent.

G Dictionary word (no received text), but note that Hóum( méngshū 侯馬盟書 (1976: 350) does use the graph. Xǔ Shèn's analytic gloss seems to show a certain analytic interest on his part in notions that are close to modern ideas of naiveness.
PIF 毳 [此芮切; LH tshyâs, OCM tshots] phonetic in 忩 [千短切(《廣韻》: 呼骨切); LH tshuânB, OCM tshôn) and LH hu\#t, OCM hût]
PR     F*i sh*ng 非聲.

105. SW 10B 408: 105; DXB 219 (10B 15b); GL vol 11: 4713; Duàn 508 (10B 36a); TKJ 1457; Ozaki vol. 5: 1037, see
variants in WGY: 452.


HG:	(⟆),
G:	疾利口也。
SSF1: 从心,
SSF2: 从冊。
IQ:	《詩》曰：
相時[⟆]民。

xi+n,
is 'quick eloquence'.
It has xīn ``heart'' as a semantic constituent,
it (also) has cè ``volume'' as a semantic constituent.
The Book of Odes (actually the Shàngshū Páng*ng) (SSJZS: 169c) says:
"Look at those smooth-tongued people."



G	Xǔ Shèn's use of postposed kǒu 口 to mean 'with the mouth as an articulatory organ' goes beyond current commentarial style and seems to have a colloquial touch.
IQ	The Shàngshū (Pán g*ng 盤庚 18.12.1) 相時䮮民 ``Look at those dispersed (homeless) people."
HP	xi&n ⟆[息廉切; LH sem(C), OCM sêm(s)]

***
TENSENESS SERIES (106-(110)-112)

106. SW 10B 408: 106; DXB 219 (10B 15b); GL vol 11: 4714b; Duàn 508 (10B 36a); TKJ 1457; Ozaki vol. 5: 1038.

HG: `` (急),
G:	褊也。
SSF: 从心,
PIF: 及聲。


jí, ``in straits"
is 'narrowed-down/tightened'. [[is '(a way of) being narrowed down/tightened (scil. by difficult circumstances)'.]]
It has xīn ``heart'' as a semantic constituent,
jí (GSR 681a: ``reach, ...") is the phonetic constituent.

HG Not infrequently, Xǔ Shèn closes a series with the definition of the common semantic denominator term. Here, on the other hand, he starts the series with the common denominator 急.
G Xǔ Shèn glosses bi(n 褊 as ``for clothes to be narrow'' and 'purists' might complain that this haute-couture reading is inapplicable to the present context and shows Xǔ Shèn to be inconsistent. However, 'purists' of this sort misconstrue Xǔ Shèn's dictionary. His gloss for the graph bi(n 褊 is designed to explain the presence of the semantic constituent yī 衣 and it does not purport to give an account of the variety of meanings of the character bi(n, or indeed to identify the basic meaning of bi(n. What Xǔ Shèn's gloss identifies is the meaning which he thought was relevant to the graphic orthography.
    If we understand Xǔ Shèn properly here, in our tentative explanatory paraphrase, he is providing a psychologically very advanced analysis of the concept of being in trouble: it is true enough that Xǔ Shèn reflects this gloss from the )ry( (Shì yán 釋言 2.24) (Xú Cháohuá 徐朝華1994: 82), but his explanation of the graph here suggests the idea that difficult situations constrain.
PIF   及 [巨立切; LH g\$\#p, OCM g\#p] phonetic in 急 [居立切; LH k\$\#p, OCM k\#p]
PR   Homogeneous initials.



107. SW 10B 408: 107; DXB 219 (10B 15b); GL vol 11: 4714b; Duàn 508 (10B 36a); TKJ 1457; Ozaki vol. 5: 1038.


HG:
G: SSF: PIF:

CAS:

(忷),
憂也。从心, 䶫聲。
一曰: 急也。

bi&n,
is 'to worry'.
It has xīn ``heart'' as a semantic constituent,
bi&n (GSR 219a: ``SW says criminal accusing one another (no text), ...") is the phonetic constituent.
An (alternative) source says: It is 'to be in straits'.

HG Dictionary word, no text.
G    Duàn rewrites 憂 as ⤏.
PIF Note that just as in syntax there are such things as discontinuous words like German anfangen ``begin'' in das fängt erst morgen an ``this only begins tomorrow", so in graphology there are graphically discontinuous elements: the graphic constituent 䶫 is a case in point because what it attaches to is in fact inserted in a way that makes biàn 䶫 into a discontinuous graphic constituent.
䶫[方免切; LH p\$anB, OCM pren] phonetic in 忷 [方䗽切; LH p\$anB, OCM pren ?]
PR   Tóngyīn 同音 OCM and LH homophones.
CAS It is rare in Shuōwén that it is the alternative gloss which establishes a homogloss link to the preceding lexical entry.
Duàn rewrites 急 with an archaising form of that character.

108. SW 10B 408: 108; DXB 219 (10B 15b); GL vol 11: 4715b; Duàn 508 (10B 36b); TKJ 1457; Ozaki vol. 5: 1039.


HG:
G: SSF: PIF:

(徸),
疾也。从心, 亟聲。

jí,
is 'to be over-eager'.
It has xīn ``heart'' as a semantic constituent,
jí (GSR 910a: ``urgently, ...") is the phonetic constituent.

CAS: 一曰:
謹重皃。

An (alternative) source says: descriptive of diligent seriousness.

G	On the basis of Yùnhuì 韻會 Duàn rewrites this as 急性也, and this would indeed have been an easier gloss.
PIF     亟[紀力切,又 去吏切; LH k\$\#k, OCM k\#k] phonetic in 徸 [己力切; LH k\$\#k, OCM k\#k]
PR	Tóngyīn 同音 OCM and LH are homophones.



109. SW 10B 408: 109; DXB 219 (10B 15b); GL vol 11: 4716; Duàn 508 (10B 36b); TKJ 1458; Ozaki vol. 5: 1039.
HG:  « (懁), juàn,

G:	急也。
SSF: 从心,

is 'to be in straits'.
It has xīn ``heart'' as a semantic constituent,

PIF: SPI:

瞏聲。 讀若絹。

qióng (GSR 256h: ``turn round, ...") is the phonetic constituent. Pronounced like juàn.

G	Duàn rewrites 急 in an archaising way at a point where what is at issue is not the graph but the word.
PIF     瞏[渠營切; LH kyenC, OCM kwens] phonetic in 懁 [古縣切; LH Guan, OCM wrên]
PR	F*i sh*ng 非聲
SPI     Juàn 懁 [古縣切; LH Guan, OCM wrên] read like juàn 絹 [吉掾切; LH kyenC, OCM kwens]

110. SW 10B 408: 110; DXB 219 (10B 15b); GL vol 11: 4716b; Duàn 508 (10B 36b); TKJ 1458; Ozaki vol. 5: 1040.


HG:
G: SSF: PIF:

(⡁),
恨也。从心, 巠聲。

xìng,
is 'to resent'.
It has xīn ``heart'' as a semantic constituent,
jīng (GSR 831a: ``Shuowen says: a vein of water (no text)") is the phonetic constituent.

HG This entry is an intrusion into the ``tenseness'' series.
PIF     巠 [GY:古靈切; LH keng, OCM kêng] phonetic in ⡁ [胡頂切; LH gengC, OCM gêngh]

111. SW 10B 408: 111; DXB 219 (10B 15b); GL vol 11: 4716b; Duàn 508 (10B 37a); TKJ 1458; Ozaki vol. 5: 1040, see
variants in WGY: 452.


HG:	(),
G:	急也。
SSF1: 从心,
SSF2: 从弦，
PIF:	弦亦聲。

xián,
is 'to be in straits'.
It has xīn ``heart'' as a semantic constituent,
it (also) has xián ``string'' as a semantic constituent. xián (GSR 366f: ``bow string ...) is at the same time the phonetic constituent.

SNT:

河南密縣有	亭。 In Hénán Mì district there is the village Xián.

HG Dictionary word, no text.



G Duàn again rewrites 急 in an archaising way at a point where what is at issue is not the graph but the word.
SNT At many points in his dictionary, Xǔ Shèn shows a characteristic non-literary interest in place names involving rare characters. Xǔ Shèn is thus not limited in his focus to summarising literary glosses: the dictionary takes due note of non-literary uses of graphs to write vernacular dialects on the one hand and to designate places on the other.
PIF     弦[胡田切; LH gen, OCM gǐn] phonetic in	[胡田切; LH gen, OCM gǐn]
PR	Tóngyīn 同音 OCM and LH are homophones.

112. SW 10B 408: 112; DXB 220 (10B 16a); GL vol 11: 4717; Duàn 508 (10B 37a); TKJ 1458; Ozaki vol. 5: 1041.
HG: ˆ (慓), piào,

G:	疾也。
SSF: 从心,
PIF: 票聲。

is 'to be eager'. [[is '(a way of) being eager (scil. so as to be (psy- chologically) in a hurry)'.]]
It has xīn ``heart'' as a semantic constituent,
piào (GSR 1157b: ``Shuowen says: leaping flames, ...") is the phonetic constituent.

PIF     票[方昭切; LH phiau, OCM phiau] phonetic in 慓 [敷沼切; LH phiauC, OCM phiauh] (our
Shuōwén fǎnqiè interpretation of the phonetic is based on this tentative suggestion).
    Duàn considers that the word written by 票 is registered in Shuōwén under the graph 爂 TKJ: 1384.
***

113. SW 10B 408: 113; DXB 220 (10B 16a); GL vol 11: 4717; Duàn 508 (10B 37a); TKJ 1458; Ozaki vol. 5: 1041.

HG: ƒ (懦),
G:	駑弱者也。
SSF: 从心,
PIF: 需聲。


nuò,
is 'someone who is listless'. [[is '(a way of) feeling listless (scil. because of lack of courage)'.]]
It has xīn ``heart'' as a semantic constituent,
xū (GSR 134a: ``wait, ...") is the phonetic constituent.

HG     Duàn rewrites the head graph as 䓩 and declares 懦 to be the work of ``an unsubtle mind'' 淺人.
G	The gloss seems to aim at a nominal meaning of 懦, but the word is currently verbal, and it is not clear why Xǔ Shèn chooses his nominal gloss. This is not very common in Shuōwén.



    Duàn omits the zh\% 者 without giving reasons, as if there were no other Shuōwén glosses ending in zh\% 者. Since nú 駑 is not a head graph in Shuōwén he rewrites this as nú 奴, as if there were not many cases in which Xǔ Shèn uses graphs which are not glossed in Shuōwén. One may sympathise with Duàn when he wishes there had been no zh\% 者 in this gloss, but one has less sympathy for his surreptitious omission of this character.
PIF 需 [相俞切; LH sio, OCM sno] phonetic in 懦 [人朱切(《廣韻》：奴䎵切); LH ño, UCM no]


114. SW 10B 408: 114; DXB 220 (10B 16a); GL vol 11: 4718; Duàn 508 (10B 37b); TKJ 1458; Ozaki vol. 5: 1042.


HG:
G: SSF: PIF:

½ (恁), 下齎也。从心,
任聲。


rèn,
is 'to be in low spirits' (?).
It has xīn ``heart'' as a semantic constituent,
rèn (GSR 667f: ``endure, ...") is the phonetic constituent.

G      Here, as elsewhere, Xǔ Shèn's gloss is not much less opaque than the character glossed. It is an interesting question whether what is an opaque gloss to us today might have been good colloquial Chinese to Xǔ Shèn's contemporaries. Supposing that it was current Chinese, one would like to know why it left so little trace in the considerable literature of his time. A possible answer, and indeed a tempting assumption, is that cases like these represent Hàn colloquialisms.
PIF     任[如林切; LH ñimB, OCM n\#m)] phonetic in 恁 [如甚切; LH ñimB, OCM n\#m)]
PR	Tóngyīn 同音 OCM and LH are homophones.

115. SW 10B 408: 115; DXB 220 (10B 16a); GL vol 11: 4718b; Duàn 508 (10B 37b); TKJ 1459; Ozaki vol. 5: 1043.


HG:
G: SSF: PIF:

(X),
失常也。从心,
代聲。

tè,
is 'to be abnormal'.
It has xīn ``heart'' as a semantic constituent,
dài (GSR 918i: ``take the place of, ...") is the phonetic con- stituent.

G	This gloss can be taken to be either strictly analytical or indeed as remarkably colloquial. The link to what precede suggests that we do have a colloquialism here.
PIF     代[徒耐切; LH d\#C, OCM l*+ kh] phonetic in X [他得切; LH d\#C, OCM l*+ kh]



116. SW 10B 408: 116; DXB 220 (10B 16a); GL vol 11: 4719; Duàn 508 (10B 37b); TKJ 1459; Ozaki vol. 5: 1043, see
variants in WGY: 453.


HG:
G: SSF: PIF:

  (怚),
驕也。从心, 且聲。


jù,
is 'arrogant'.
It has xīn ``heart'' as a semantic constituent,
jū (GSR 46a: ``moreover, ...") is the phonetic constituent.

PIF This conveniently illustrates that the phonetic value of phonetic constituents does not have to be the ancestor of the most current modern pronunciation. Phonetic constituents can be multivalent.
(jū) 且 [子余切; LH tsiâB, OCM tsâ)] or (qi\%) 且 [千也切; LH tshiaB, OCM tsha)] phonetic in
怚 [子去切(《廣韻》：將預切); LH tsiâC, OCM tsâh]

117. SW 10B 408: 117; DXB 220 (10B 16a); GL vol 11: 4719; Duàn 508 (10B 37b); TKJ 1459; Ozaki vol. 5: 1043.


HG:
G:

SSF: PIF:

E (悒),
不安也。
从心, 邑聲。


yì,
is 'restless'. [[is '(a way of being) restless (scil. psycholo- gically)'.]]
It has xīn ``heart'' as a semantic constituent,
yì (GSR 683a: ``city, ...") is the phonetic constituent.

PIF     邑[於汲切; LH )\$p, OCM )\#p] phonetic in 悒 [於汲切; LH )\$p, OCM )\#p]
PR	Tóngyīn 同音 OCM and LH are homophones.

118. SW 10B 408: 118; DXB 220 (10B 16a); GL vol 11: 4719b; Duàn 509 (10B 38a); TKJ 1459; Ozaki vol. 5: 1043.


HG:
G1:
G2: SSF: PIF: IQ:

GQ:

A (悆),
忘也。嘾也。从心, 余聲。
《周書》曰： 有疾不悆。 悆，喜也。


yù,
is 'to neglect'.
It is (also) 'to keep in one's mouth'.
It has xīn ``heart'' as a semantic constituent,
yú (GSR 82a: ``I, ...") is the phonetic constituent.
In the book of Zhǒu (SSJZS: 196a) it says:
"He is sick and not in high spirits.'' (tr. Karlgren) Yú (, in this example,) is to be glad."

G2	Our translation is based on Xǔ Shèn's gloss: in the entry for dàn 嘾 SW 2A 22. 101.
PIF     余[以諸切; LH jâ, OCM la] phonetic in 悆 [羊茹切; LH jâC , OCM lah]



IQ	The SSJZS edition of Shū (Jīn téng 金縢 34.1) writes: 既克商二年王有疾。弗豫 ``After the victory over the Shāng, in the second year, the king fell ill and was not happy."
GQ This is one of the few cases where Xǔ Shèn adds traditional glosses after his quotation. In such contexts what the quotation illustrates is effectively a postposed alternative gloss.

119. SW 10B 408: 119; DXB 220 (10B 16a); GL vol 11: 4721; Duàn 509 (10B 38a); TKJ 1459; Ozaki vol. 5: 1044.


HG:
G: SSF: PIF:

; (䓒),
更也。从心, 弋聲。


tè,
is 'to cause to change'.
It has xīn ``heart'' as a semantic constituent,
yì (GSR 918a: ``to shoot with arrow and string attached, ...") is the phonetic constituent.

G   Xǔ Shèn does not link the meaning of this word to the semantic constituent, except if we assume that he understood tè 䓒 to involve a deliberate act of changing, which is a standard meaning of g*ng 更, but which is not attested as a meaning for tè 䓒. However, Xǔ Shèn's choice of gloss may indeed represent an attempt to move the meaning of tè 䓒 closer to the realm of the psychological. No matter how we need to understand this gloss in the end, the great ambiguity of the graph 更, here used out of context, cannot be considered to have been as much more helpful in Hàn times than it is today.
    Duàn rewrites 更 with an archaising graph, although what is at issue is the word and not the graph.
PIF   弋[與職切; LH j\#k, OCM l\#k] phonetic in 䓒 [他得切; LH th\#k, OCM lh*+ k]

120. SW 10B 408: 120; DXB 220 (10B 16a); GL vol 11: 4721b; Duàn 509 (10B 38b); TKJ 1459; Ozaki vol. 5: 1045.


HG:
G: SSF: PIF:

  (憪), 愉也。从心,
間聲。


xián,
is 'pleasant'.
It has xīn ``heart'' as a semantic constituent,
ji+n (GSR 191b: ``interstice, ...") is the phonetic constituent.

HG The k(ishū 楷書 form of this graph gives the form of the right hand constituent 閒 which is a head graph in Shuōwén. The seal script form which Xǔ Shèn discusses, is everywhere identified as having the 'moon' constituent in the 'gate' (閒). However, the Dà Xú b\%n seems to identifies the phonetic as the 'sun' constituent in the 'gate' (間) which is not a head graph in Shuōwén. Xǔ Shèn's analysis is indifferent to this vacillation in the specification of the phonetic, but it does show that Xǔ Shèn's orthography is less than consistent. Ji&n 閒 and ji&n 間 are recognised as allographs but not explicitly analysed as such.



PIF     閒 [古閑切 (《廣韻》：古䊵切）; LH k/n, OCM krên] phonetic in 憪 [戶閒切; LH g/n, OCM grên]


121. SW 10B 408: 121; DXB 220 (10B 16a); GL vol 11: 4721b; Duàn 509 (10B 38b); TKJ 1460; Ozaki vol. 5: 1045.


HG:
G:

SSF: PIF: IQ:

m (愉),
薄也。
从心,
俞聲。
《論語》曰：


yú, ``pleasant"
is 'to be close'. [[is '(a way of being) close (scil. so as to become intimately congenial)'.]]
It has xīn ``heart'' as a semantic constituent,
yú (GSR 125a: ``agree, ...") is the phonetic constituent. The Lúnyǔ (SSJZS: 2494b) says:

私覿，愉愉如也。 ``In private audience, he was completely congenial."
G    Here Xǔ Shèn explains a fairly common word with a very common word which normally has an unrelated meaning, báo 薄 ``thin, close". The Lúnyǔ quotation below shows that the intended meaning is indeed the standard meaning ``pleased'' but it gives no help towards an explanation for the gloss báo. Our translation suggests a very tentative solution to this puzzle.
PIF     俞[羊朱切; LH jo, OCM lo] phonetic in 愉 [羊朱切; LH jo, OCM lo]
PR	Tóngyīn 同音 OCM and LH are homophones.
IQ	Lúnyǔ (Xi&ng d(ng 鄉黨 (10.5.2).
One certainly feels Xǔ Shèn needed a quotation to clarify his rather opaque gloss.


122. SW 10B 408: 122; DXB 220 (10B 16b); GL vol 11: 4723; Duàn 509 (10B 39a); TKJ 1460; Ozaki vol. 5: 1046.

HG:   ´(懱),
G:	輕易也。
SSF: 从心,


miè,
is 'make light of> despise'.
It has xīn ``heart'' as a semantic constituent,

PIF: IQ:

蔑聲。
《商書》曰： 以相陵懱。

miè (GSR 311a: ``destroy, ...") is the phonetic constituent. The Book of Sh)ng (SH() says:
"... so as to tyrannise and despise them".

G	Xǔ Shèn's binome gloss would seem to represent a disambiguating colloquialism of his time, since 輕 by itself would be even more ambiguous than 輕易.
PIF     蔑[莫結切; LH met, OCM mêt] phonetic in 懱 [莫結切; LH met, OCM mêt]
PR	Tóngyīn 同音 OCM and LH homophones.
IQ	This passage can not be found in the SSJZS edition of Shūjīng.



***
STUPIDITY SERIES (123-127)

123. SW 10B 408: 123; DXB 220 (10B 16b); GL vol 11: 4723b; Duàn 509 (10B 39a); TKJ 1460; Ozaki vol. 5: 1047, see
variants in WGY: 453.

HG: t (愚),
G:	䮴也。
SSF1: 从心,
SSF2: 从禺。


yú,
is 'simple-minded'. [EP: is 'simple-minded' (scil. so as to lack necessary or normal mental skills)'.]]
It has xīn ``heart'' as a semantic constituent,
it (also) has yú ``monkey'' as a semantic constituent.

GI:

禺，猴屬， Yú is a kind of monkey, the stupidest of the wild animal.
獸之愚者。

SSF2 Xǔ Shèn fails to specify the manifest role of yú 禺 as phonetic. What later came to be called a huì yì 會意 character deriving from two elements would be explained along exactly the same lines as those above. The standard Shuōwén line that is missing according to Xǔ Shèn's conventions would be 禺亦聲. (The question whether the yú 禺 monkey was as stupid as Xǔ Shèn suggests and whether it was currently regarded as stupid is of no concern for us here. What matters to us is the structure of Xǔ Shèn's explanation.)
GI Xǔ Shèn's supplementary note shows his preoccupation with the semantic motivation of the phonetic component in the character.
HP   yú 愚 [麌倶切; LH ng\$o, OCM ngo]

124. SW 10B 408: 124; DXB 220 (10B 16b); GL vol 11: 4724; Duàn 509 (10B 39a); TKJ 1460; Ozaki vol. 5: 1047.
HG: » (䮴), zhuàng,

G:	愚也。
SSF: 从心,
PIF: 贛聲。

is 'stupid'. [[is '(a way of being) stupid (scil. in the innocent simple-minded mode)'.]]
It has xīn ``heart'' as a semantic constituent,
gòng (SW: ``to give, ...") is the phonetic constituent.

G This case of unhelpful hùzhù 互注 shows that on occasion Xǔ Shèn is prepared to pair off quite common standard words with fairly recondite near-equivalents. One's suspicion is that zhuàng 䮴 may have been a colloquial word, in which case there is nothing unhelpful in his procedure.
PIF  贛[古送切; LH kongC , OCM kôngh ?] phonetic in 䮴 [陟絳切; LH trcngC , OCM trûngh]



125. SW 10B 408: 125; DXB 220 (10B 16b); GL vol 11: 4724b; Duàn 509 (10B 39a); TKJ 1460; Ozaki vol. 5: 1047.


HG:	(徯),
G:	姦也。
SSF: 从心,
PIF: 采聲。

c&i,
is 'depraved'.
It has xīn ``heart'' as a semantic constituent,
c&i (GSR 942: ``gather, pluck") is the phonetic constituent.

G	The Chinese notion of stupidity includes that of moral deficiency.
PIF     采[倉宰切; LH tsh\#B/C, OCM tsh*+ )/h] phonetic in 徯 [倉宰切; LH tsh\#B, OCM tsh*+ )]
HG Dictionary word, no text. Moreover the dictionaries vary disconcertingly in the meanings they attribute to the word.

126. SW 10B 408: 126; DXB 220 (10B 16b); GL vol 11: 4724b; Duàn 509 (10B 39a); TKJ 1461; Ozaki vol. 5: 1048.
HG: – (憃), chǒng,

G:	愚也。
SSF: 从心,
PIF: 舂聲。

is 'stupid' [[is '(a way of being) stupid (scil. in a crude, moronic way)'.]]
It has xīn ``heart'' as a semantic constituent,
chǒng (GSR 1192a: ``to hull grain with a pestle, ...") is the phonet- ic constituent.

PIF     舂[書容切; LH shong, OCM lhong] phonetic in 憃 [丑江切; LH trhong, OCM rhong]
PR	Note the non-homogeneous initials


127. SW 10B 408: 127; DXB 220 (10B 16b); GL vol 11: 4725; Duàn 509 (10B 39a); TKJ 1461; Ozaki vol. 5: 1048, see
variants in WGY: 453.

HG:	(䓷), aì,

G:	䌭也。
SSF1: 从心,
SSF2: 从疑，

is 'moronic'.
It has xīn ``heart'' as a semantic constituent,
it (also) has yí ``doubt'' as a semantic constituent.

PIF: CAS:

疑亦聲。
一曰: 惶也。

yí (GSR 956a: ``doubt, ...") is at the same time the phonetic constituent.
An (alternative) source says: is 'to be terrified'.

SSF2 Note that the Cartesian positive perspective on doubt was alien to ancient Chinese thinking where doubt is generally taken as a sign of stupidity rather than intelligence. An interesting



exception to this general rule is presented by the case of Wáng Ch\%ng 王充 who cultivates a gentle attitude of doubt in his writings.
PIF     疑[語其切; LH ng\$\#, OCM ng\#] phonetic in 䓷 [五䰬切; LH ng\#C , OCM ng*+ h]
***

128. SW 10B 408: 128; DXB 220 (10B 16b); GL vol 11: 4725; Duàn 509 (10B 39b); TKJ 1461; Ozaki vol. 5: 1048.


HG:
G: SSF: PIF:

  (䓖),
很也。从心, 支聲。


zhì, ``jealous resentment'' is 'resentful'.
It has xīn ``heart'' as a semantic constituent,
zhī (GSR 864a: ``branch, ...") is the phonetic constituent.

G     When Xǔ Shèn writes 很 he intends hèn 恨 as interpreted in the present section see: 182 below. The Shuōwén gloss for h\%n 很 is 不聽从也 'be recalcitrant' and is hard to fit into this context. Gu(ngy( 廣雅 5 Shìgǔ 釋詁 4 defines 很 as 恨.
PIF   支[章移切; LH tshe < kie, OCM ke] phonetic in 䓖 [之義切; LH tsheC < kieC , OCM keh]

129. SW 10B 408: 129; DXB 217 (10B x); GL vol 11: 4725b; Duàn 509 (10B 39b); TKJ 1461; Ozaki vol. 5: 1049.

HG: D (悍),
G:	勇也。
SSF: 从心,
PIF: 旱聲。


hàn, ``foolhardy"
is 'courageous'. [[is '(a way of being) courageous (scil. to an irresponsibly high degree)'.]]
It has xīn ``heart'' as a semantic constituent,
hàn (GSR 139s: ``drought, ...") is the phonetic constituent.

PIF     旱[矦旰切; LH gânB, OCM gân)] phonetic in 悍 [侯旰切; LH gânC, OCMgâns]

130. SW 10B 408: 130; DXB 220 (10B 16b); GL vol 11: 4726; Duàn 509 (10B 39b); TKJ 1461; Ozaki vol. 5: 1049, see
variants in WGY: 453.

HG:  „ (態),
G:	意也。
SSF1: 从心,
SSF2: 从能。


tài, ``expressive appearance"
is 'thought'. [[is '(a kind of) thought (scil. as manifested in one's appearance)'.]]
It has xīn ``heart'' as a semantic constituent,
it (also) has néng ``able'' as a semantic constituent.



A:	[ 嫔] ， In the graph , [the word tài] is alternatively written with the
或从人。 (graphic) constituent rén.
G       Duàn gratuitously adds a character: 意態也, as if all Xǔ Shèn's definitions were analytic.
SSF2 Duàn rewrites this as 从心能, as if the other formula in the editions he used was not also current throughout Shuōwén. (Xi(o Xú b\%n p. 210 writes 從心能).
    The OCM reconstructions show a clear rhyme which one would expect can justify the treatment of néng 能 as a phonetic constituent. It is tempting to surmise that it was only Xǔ Shèn's insufficient grasp of Old Chinese phonology that prevented him from treating néng 能 as phonetic. But, before accusing of such ignorance, we must remember that Xǔ Shèn commonly omitted to mention the phonophoric function of the semantic constituent.
A	In the context of the discussion of allographs, the formula 从X is purely graphic in meaning and is translatable simply as ``has the graphic constituent".
The word huò 或 in this formula must be taken to mean something like ``alternatively'' (to the standard version).
HP	能 [奴登切; LH n\#(ng), OCM n*+ !] not declared phonetic in 態 [他代切; LH th\#C , OCM nh*+ h]

131. SW 10B 408: 131; DXB 220 (10B 16b); GL vol 11: 4726b; Duàn 509 (10B 39b); TKJ 1461; Ozaki vol. 5: 1049.
HG:  \% (怪), guài, ``extraordinary"

G:	異也。
SSF: 从心,
PIF: 圣聲。

is 'different'. [[is '(a way of being) different (scil. from the ordi- nary or expected)'.]]
It has xīn ``heart'' as a semantic constituent,
kū (SW: ``(dialect) work the land, ...") is the phonetic constituent.

G	While it is true that yì 異 can come to mean ``consider as strange'' and the like, it is important to realise that in this gloss Xǔ Shèn provides a hypernym rather than a supposed synonym.
PIF     圣[苦骨切; LH khu\#t ?, OCM khw*+ t ?] phonetic in 怪 [古壞切; LH ku/C, OCM kwr*+ h]

***
NEGLIGENCE SERIES (132-(137,138)-142)

132. SW 10B 408: 132; DXB 220 (10B 17a); GL vol 11: 4727; Duàn 509 (10B 39b); TKJ 1461; Ozaki vol. 5: 1050.


HG:
G:

(⨂),
放也。

dàng,
is 'unrestrained'.




SSF: PIF:

从心, 象聲。

It has xīn ``heart'' as a semantic constituent,
xiàng (GSR 728a: ``elephant, ...") is the phonetic constituent.

G	Compare, perhaps, the binome fàngdàng 放蕩 'be unrestrained'.
PIF     [徐兩切; LH ziângB, OCM s-jang)] phonetic in [徒朗切; LH dângC, OCM lângh]
PR	Note the non homogeneous initials.


133. SW 10B 408: 133; DXB 220 (10B 17a); GL vol 11: 4727b; Duàn 509 (10B 39b); TKJ 1462; Ozaki vol. 5: 1050.


HG:
G:

SSF: PIF:

Œ(慢),
惰也。
从心,
曼聲。


màn, ``negligent"
is 'indolent'. [[is '(a way of being) indolent (scil. with respect to what requires dedication, diligence and seriousness)'.]]
It has xīn ``heart'' as a semantic constituent,
màn (GSR 266a: ``extended, ...") is the phonetic constituent.

CAS: 一曰:

An (alternative) source says:

慢，不畏也。 màn is 'to fail to show proper reverence'.
G The reading 'slow' attaches to the character mán 渟 in Shuōwén. It appears that these two manifestly related words had different tones, and that at later stages of the language the two words were merged and received the falling tone.
PIF    曼[無販切; LH mânC, OCM mâns] phonetic in 慢 [謀晏切; LH manC, OCM mrâns]

134. SW 10B 408: 134; DXB 220 (10B 17a); GL vol 11: 4728; Duàn 509 (10B 39b); TKJ 1462; Ozaki vol. 5: 1050.

HG:          (怠), dài, ``lazy"

G:	慢也。
SSF: 从心,
PIF: 台聲。

is 'negligent'. [[is '(a way of being) negligent (scil. out of laziness)'.]]
It has xīn ``heart'' as a semantic constituent,
yí (GSR 976p: ``loan for t'ai (globe-fish-like, globular =) rounded...") is the phonetic constituent.

PIF Yí 怡 (n° 42) and dài 怠 are both analysed as 从心台聲. According to our interpretation, the meaning of the formula 台聲 is exactly the same in both cases: ``台 is phonetic". The ambiguity in the pronunciation of the phonetic constituent is irrelevant to Xǔ Shèn's analysis.
台[與之切; LH j\#, OCM l\#] phonetic in 怠 [徒亥切; LH d\#B, OCM l\#)]



135. SW 10B 408: 135; DXB 220 (10B 17a); GL vol 11: 4728; Duàn 509 (10B 39b); TKJ 1462; Ozaki vol. 5: 1050.

HG: ­ (懈),
G:	怠也。
SSF: 从心,
PIF: 解聲。


xiè, ``remiss"
is 'lazy'. [[is '(a way of being) lazy (scil. often for lack of energy, or because one is tired or exhausted)'.]]
It has xīn ``heart'' as a semantic constituent,
ji\$ (GSR 861a: ``dissolve, ...") is the phonetic constituent.

PIF     解[佳買切; LH k/B, OCM krê)] phonetic in 懈 [古隘切; LH k/C, OCM krêh]

136. SW 10B 408: 136; DXB 220 (10B 17a); GL vol 11: 4728b; Duàn 509 (10B 39b); TKJ 1462; Ozaki vol. 5: 1051, see
variants WGY: 454.


HG:
G:


SSF0: IQ:

A:

SA:

(憜),
不敬也。

从心ᗫ省。
《春秋傳》曰：
"執玉憜。"
[惰]，憜
或省��。
[䑠]，古文。

duò,
is 'to fail to show respectful diligence'. [[is '(a way of) failing to show respectful diligence (scil. out of inertia rather than straightforward lack of respect)'.]]
It has xīn ``heart'' and an abbreviated form of duò ᗫ as semantic constituents,
The Springs and Autumns Tradition (i.e. Zuǒzhuàn, SSJZS: 1802b) says:
"They took the jade with unceremonious negligence."
In the variant , (the word) duò is alternatively written omitting fù.
 is an ancient style graph.

G The final gloss in the series, here as often elsewhere, may be taken to recapitulate the common feature in the whole series.
SSF0 Phonetic constituents are often characterised by the formulaic suffix sh\%ng 省. The present instance shows up what is logically evident: also semantic constituents can take abbreviated forms. For example, the head graph duò 隋 is analysed as having two semantic constituents, one unabbreviated and one abbreviated:从肉，从隓省. In view of this formulaic practice there is no structural need to emend the text as it stands. Duàn's addition of sh*ng 聲 is disingenuous because it seems to assume, contrary to fact, that Xǔ Shèn always points out the phonetic functions of his semantic constituents. Moreover, the sequence 从X省聲 is unacceptable as a statement to the effect that X is both a semantic constituent and a phonetic component. Abbreviation is most common in phonophoric constituent. Abbreviated semantic constituents are naturally rare in Shuōwén. Moreover we cannot see that any of these few abbreviated semantic constituents are ever said to be at the same time phonetic constituents.




Xǔ Shèn does not make explicit the phonetic nature of the reduced form of duò ᗫ.
IQ	The SSJZS edition of the Zuǒzhuàn (X( 僖 11.2) writes 受玉惰 ``He received the nephrite with an air of indifference.'' (tr. Legge).
HP	ᗫ[GY: 徒果切; LH duâiB/C, OCM lôi)/h] is not declared phonetic in 憜 [徒果切; LH duâiB, OCM lôi)]


137. SW 10B 408: 137; DXB 220 (10B 17a); GL vol 11: 4729b; Duàn 510 (10B 40a); TKJ 1462; Ozaki vol. 5: 1051.


HG:  (⧵/慫),
G:	驚也。
SSF: 从心,

sǒng,
is 'scared'.
It has xīn ``heart'' as a semantic constituent,

PIF: SPI:

從聲。 讀若悚。

cóng (GSR 1191d: ``follow, ...") is the phonetic constituent. Pronounced like sǒng.

HG This entry is an intrusion into the ``negligence'' series.
    Duàn surreptitiously corrects the Sòng edition head graph, writing the heart radical at the bottom of the character, in accordance with the seal graph. In this case Duàn is perfectly entitled to exercise his free judgement since the k(ishū 楷書 head graphs are additions from the Sòng dynasty and do not in anyway go back to Xǔ Shèn's own time, but his rewriting is arbitrary nonetheless.
PIF    從[慈用切; LH dziong, OCM *dzong] phonetic in ⧵/慫 [息拱切; LH siongB, OCM song)]
SPI     Sǒng ⧵/慫 [息拱切; LH siongB, OCM song)] read like sǒng 悚 [LH siongB, OCM *songʔ]
    Xǔ Shèn provides a homophonous and near-synonymous word after dúruò 讀若. The main purpose does not seem to be to show how exactly a problematic word is to be pronounced, but to bring to the reader's attention another relevant word that may or may not be etymologically related to the present one.
    Duàn suggests that sǒng 悚 should be rewritten as sǒng 竦 because 悚 is not in Shuōwén, as if all the graphs used in Shuōwén were explained in Shuōwén. As it happens the word standardly written sǒng 悚 today is entered under the graph sǒng 愯, but there are many instances where Xǔ Shèn's text does not follow the orthography established in his own system of head graphs. What Duàn fails to realise is that the graph 悚 is used here to identify a word of the spoken language and not as a graph to be discussed. Characters are used in the Shuōwén in two subtly distinct ways: in certain contexts a character is used to refer to the grapheme it is a token of. This is the case in head graphs and also in the identification of phonophoric and semantic constituents as well as in the discussion of variant graphs. In other contexts, notably in glosses and other semantic comments, characters are used to represent words of the language. When Duàn substitutes a homophone he should have been looking for the standard way in Shuōwén to write the given word, i.e. he should have chosen 愯. However, as we have seen above the need to look for an emendation arises only




on the mistaken assumption that Xǔ Shèn's orthography was consistent and that he only used characters which are also explained as head graphs.

138. SW 10B 408: 138; DXB 220 (10B 17a); GL vol 11: 4730; Duàn 510 (10B 40a); TKJ 1462; Ozaki vol. 5: 1052.


HG:
G: SSF: PIF:

& (怫), 鬱也。从心,
弗聲。

fú,
is 'depressed'.
It has xīn ``heart'' as a semantic constituent,
fú (GSR 500a: ``not, ...") is the phonetic constituent.

HG This entry is an intrusion into the ``negligence'' series.
G   For reasons we find difficult to understand Duàn wrongly rewrites this character as 鷯, when in fact the binome fú yù 怫鬱 'depressed' seems to explains Xǔ Shèn's gloss satisfactorily: it is very common in Shuōwén that a graph is glossed in terms of another word with which it forms a binome compound.
PIF   弗[普未切; LH put, OCM p\#t] phonetic in 怫 [符弗切; LH but or bus, OCM b\#t(s)]
PR   Note the homogeneous initials.


139. SW 10B 408: 139; DXB 220 (10B 17a); GL vol 11: 4730; Duàn 510 (10B 40a); TKJ 1463; Ozaki vol. 5: 1052.


HG:
G: SSF: PIF:

IQ:

(3),
忽也。从心, 介聲。
《孟子》曰：
孝子之心不若是3。

xiè,
is 'negligent'.
It has xīn ``heart'' as a semantic constituent,
jiè (GSR 327a: ``armour, ...") is the phonetic con- stituent.
The Mèngzi (SSJZS: 2733c) says:
"The attitude of a filial son is not negligent like this."

PIF     介[古拜切; LH k/t, k/s, OCM krêt(s)] phonetic in 3 [呼介切; LH h/s, OCM hrêts]
PR   Note the homogeneous initials.
IQ It seems that Xǔ Shèn adds an illustrative quotation here because the word he has glossed is rare. What remains unclear is why so many exceedingly rare graphs have no such illustrative gloss on the one hand, and how he chose the head characters for which to add such illustrative quota- tions.
The SSJZS edition of Mencius (Wàn zh&ng 萬章) (5.A.1) writes jiá 恝 instead of 3: 夫公明
高以孝子之心為不若是恝。"Now G\%ngmíng Gāo did not think that a son's mind was such that he would be so complacent."



140. SW 10B 408: 140; DXB 220 (10B 17a); GL vol 11: 4731b; Duàn 510 (10B 40a); TKJ 1463; Ozaki vol. 5: 1052.
HG: ; (忽), hū, ``disregard"

G:	忘也。
SSF: 从心,
PIF: 勿聲。

is 'to forget about'. [[is '(a way of) forgetting about (scil. what may not be worth taking account of)'.]]
It has xīn ``heart'' as a semantic constituent,
wù (GSR 503a: ``don't, ...") is the phonetic constituent.

G	Note that the meaning 'forget' is much rarer in the early sources than the clearly distinct meaning 'forget about'.
PIF     勿[文弗切; LH mut, OCM m\#t] phonetic in 忽 [呼骨切; LH hu\#t, OCM hm*+ t]

141. SW 10B 408: 141; DXB 220 (10B 17b); GL vol 11: 4731b; Duàn 510 (10B 40a); TKJ 1463; Ozaki vol. 5: 1053, see
variants in WGY: 454.
HG: § (忘), wàng, ``to forget (about)"

G:	不識也。
SSF1: 从心,
SSF2: 从亡， PIF: 亡亦聲。

is 'to be unaware'. [[is '(a way of being) unaware (scil. of what one ought to remember)'.]]
It has xīn ``heart'' as a semantic constituent,
it (also) has wáng ``lack'' as a semantic constituent.
Wáng (GSR 742a:"... not exist, ...") is at the same time the pho- netic constituent.

G Xǔ Shèn's gloss shows that he took wàng 忘 not in its most common meaning 'to neglect what one is aware of, fail to pay proper attention to what one is aware of' but in the now more current meaning 'to forget'. If we understand his gloss in n° 140 correctly, he has just used the word in its old acceptance. Note that the notion of forgetting is easier to define than that of neglecting.
SFF     Duàn rewrites this as 从心,亡聲, basing himself on the Yùn huì 韻會 edition of Shuōwén.
PIF     The relevant reading of 亡 is wáng and not wú. It turns out that the reading wú is unattested in early sources, although one feels sure that the reading is of early origin.
    亡 [TKJ: 1812, 武方切; LH muâng, OCM mâng] phonetic in 忘 [武方切; LH muâng(C), OCM mâng]! (level tone)


142. SW 10B 408: 142; DXB 220 (10B 17b); GL vol 11: 4732; Duàn 510 (10B 40b); TKJ 1463; Ozaki vol. 5: 1053.


HG:
G1:

‘ (慲), mán, ``obtuse"
忘也。	is 'oblivious'. [[is '(a way of being) oblivious (scil. of what anyone should know)'.]]




G2: SSF: PIF:

慲兜也。从心,
䎗聲。

It is as in mándǒu ``all obtuse".
It has xīn ``heart'' as a semantic constituent,
mán (GSR 183a: ``even (>plain), ...") is the phonetic constituent.

G2 The alternative gloss seems to give a colloquial association rather than a new definition.
PIF     䎗[母官切; LH mân, OCM mân] phonetic in 慲 [毋官切; LH mân, OCM mân]
PR	Tóngyīn 同音 OCM and LH homophones.
***

Minor INDULGENCE SERIES (143-144)

143. SW 10B 408: 143; DXB 217 (10B x); GL vol 11: 4732b; Duàn 510 (10B 40b); TKJ 1463; Ozaki vol. 5: 1053.
HG: 3 (恣), zì, ``indulgent"

G:	縱也。
SSF: 从心,
PIF: 次聲。

is 'to set free'. [[is '(a way of) setting free (scil. oneself, so as to become unrestrained)'.]]
It has xīn ``heart'' as a semantic constituent,
cì (GSR 555a: ``next in order, ...") is the phonetic constituent.

PIF     次[七四切; LH tshiC, OCM tshih or tshis] phonetic in 恣 [資四切; LH tsiC, OCM tsih]

144. SW 10B 408: 144; DXB 220 (10B 17b); GL vol 11: 4733; Duàn 510 (10B 40b); TKJ 1464; Ozaki vol. 5: 1053.


HG:

g (愓), dàng, ``profligate"

G:

SSF: PIF: CAS:

放也。
从心, 昜聲。一曰: 平也。

is 'to let go'. [[is '(a way of) letting go (scil. oneself, in a self- indulgent way)'.]]
It has xīn ``heart'' as a semantic constituent,
yáng (GSR 720a: ``South side, ...") is the phonetic constituent. An (alternative) source says:
it is 'peaceful'.

HG     Cf. dàng 蕩 ``profligate".
PIF     昜[與章切; LH jâng, OCM lang] phonetic in 愓 [徒朗切; LH dângB, OCM lâng)]
***



145. SW 10B 408: 145; DXB 220 (10B 17b); GL vol 11: 4733; Duàn 510 (10B 40b); TKJ 1464; Ozaki vol. 5: 1054, see
variants in WGY: 454.


HG:
G:
SSF: PIF:

Ÿ (憧),
意不定也。
从心, 童聲。


chǒng,
is 'to be unsettled in one's thought'.
It has xīn ``heart'' as a semantic constituent,
tóng (GSR 1188o: ``boy, ...") is the phonetic constituent.

G	Xǔ Shèn makes the psychological of this explicit in this unusually discursive definition.
PIF     童[徒紅切; LH dong, OCM dông] phonetic in 憧 [尺容切; LH tshhong, OCM thong]

146. SW 10B 408: 146; DXB 220 (10B 17b); GL vol 11: 4733b; Duàn 510 (10B 40b); TKJ 1464; Ozaki vol. 5: 1054.


HG:
G: SSF: PIF:

H (䇡),
䏘也。从心, 里聲。


kuī,
is 'to scoff at'.
It has xīn ``heart'' as a semantic constituent,
l( (GSR 978a: ``village, ...") is the phonetic constituent.

SN:

《 春秋傳》In the Springs and Autumns' Tradition (i.e. Zuǒzhuàn, SSJZS:


CAS:

有孔䇡。一曰:
病也。

2175b) there is K/ng Kuī. An (alternative) source says: It is 'to be ill'.

G     This unusual meaning is also attested in Wénxu(n, but the more current meaning for the word is ``to be worried or sad".
SN Xǔ Shèn's supplementary note on proper names are of particular interest. Quite generally the use of characters in proper names of all kinds should be an integral part of their lexical description, a point that is brought out nicely in P. Chantraine (Dictionnaire étymologique de la langue grecque. Histoire des mots, 1999).
Zuǒzhuàn (0i 哀15).
PIF     里[良止切; LH li\#B, OCM r\#)] phonetic in 䇡 [苦回切; LH khui, OCM ?]
PR	F*i sh*ng 非聲.

147. SW 10B 408: 147; DXB 220 (10B 17b); GL vol 11: 4734; Duàn 510 (10B 41a); TKJ 1464; Ozaki vol. 5: 1055.


HG:
G: SSF:

¥ (憰), 權詐也。从心,


jué,
is 'to take an opportunity to cheat'.
It has xīn ``heart'' as a semantic constituent,



PIF:     p聲。	yù (GSR 507a: ``(Shuowen: 'to pierce' (no text))") is the phonet- ic constituent.
G	The binome quánzhà 權詐 seems surprising, we suspect it is a Hàn colloquialism which is attested in Wáng Ch\%ng.
PIF     p [余律切; LH juit, OCM wit] phonetic in 憰 [古穴切; LH kwet, OCM kwǐt]

148. SW 10B 408: 148; DXB 220 (10B 17b); GL vol 11: 4734b; Duàn 510 (10B 41a); TKJ 1464; Ozaki vol. 5: 1055.


HG:	(律),
G:	誤也。
SSF: 从心,
PIF: 狂聲。

guàng,
is 'to be mistaken'.
It has xīn ``heart'' as a semantic constituent,
kuáng (GSR 739o: ``foolish, ...") is the phonetic constituent.

PIF     Duàn writes an archaising graph for 狂, as if there was a convention in Shuōwén to the effect that phonetic constituents are listed in their seal form.
狂[巨王切; LH gyâng, OCM gwang] phonetic in 律 [居況切; LH kyangC, OCM kwângh]

149. SW 10B 408: 149; DXB 220 (10B 17b); GL vol 11: 4734b; Duàn 510 (10B 41a); TKJ 1464; Ozaki vol. 5: 1055, see
variants in WGY: 454.


HG:
G: SSF: PIF:

( (怳),
狂之皃。从心,
況省聲。


hu&ng,
the appearance of being crazy.
It has xīn ``heart'' as a semantic constituent,
it has an abbreviated form of kuàng (GSR 765g:"increase, ...") as the phonetic.

G The formula X之皃 while not unknown elsewhere is not the preferred form in Shuōwén. We find only nine cases.
PIF Duàn rewrites this as 兄聲 and declares the received text to be the work of someone who does not understand the ancient sounds: 不知古音者所改. The thought that the received text has been changed is a mere conjecture on Duàn's part.
況 [許訪切; LH hyang, OCM hwrang] phonetic in [許往切; LH hyângB, OCM hwang)]

150. SW 10B 408: 150; DXB 220 (10B 17b); GL vol 11: 4735; Duàn 510 (10B 41a); TKJ 1465; Ozaki vol. 5: 1056.

HG:	(恑),      gu(, ``pretend"




G:

SSF: PIF:

變也。
从心, 危聲。

is 'to change'. [[is '(a way of) changing (scil. one's appearances in order to deceive)'.]]
It has xīn ``heart'' as a semantic constituent,
w*i (GSR 29a: ``precipitous, ...") is the phonetic constituent.

PIF     危[魚爲切; LH ngyâi, OCM ngoi or OCM ngwai] phonetic in 恑 [過委切; LH kyâiB, OCM koi) or kwai)]


151. SW 10B 408: 151; DXB 221 (10B 18a); GL vol 11: 4735b; Duàn 510 (10B 41a); TKJ 1465; Ozaki vol. 5: 1056.


HG:
G: SSF: PIF:

(⬝),
有二心也。从心,
巂聲。

xié,
is 'to be in two minds'.
It has xīn ``heart'' as a semantic constituent,
guī (GSR 880a: ``Shuowen: name of a bird (no text)") is the phonetic constituent.

HG Dictionary word, no text.
G	Note the discursive definition in what almost look like a colloquial style.
PIF     巂[戶圭切 (《集韻》: 均窺切); LH Gue, OCM wê] phonetic in ⬝ [戶圭切; LH Gue, OCM wê]
PR	Tóngyīn 同音 OCM and LH homophones.

152. SW 10B 408: 152; DXB 221 (10B 18a); GL vol 11: 4736; Duàn 510 (10B 41a); TKJ 1465; Ozaki vol. 5: 1056.


HG:

g (悸), jì, ``shiver"

G:

SSF: PIF:

心動也。
从心, 季聲。

is 'to be moved in one's heart'. [[is '(a way of) being moved in one's heart (scil. by fear)'.]]
It has xīn ``heart'' as a semantic constituent,
jì (GSR 538a: ``youngest (of brothers etc.), ...") is the phonetic constituent.

G	Note how Xǔ Shèn makes the psychological content of the concept explicit by the use of the word xīn 心.
PIF     季[居悸切; LH kwis, OCM kwis] phonetic in 悸 [其季切; LH gwis, OCM gwis]



153. SW 10B 408: 153; DXB 221 (10B 18a); GL vol 11: 4736b; Duàn 510 (10B 41a); TKJ 1465; Ozaki vol. 5: 1057.


HG:
G: SSF: PIF:

ª (憿), 幸也。从心,  þ聲。


ji+o,
is 'to be lucky'.
It has xīn ``heart'' as a semantic constituent,
ji&o (GSR 1162a: ``SW says: to shine, bright (no text)") is the phonetic constituent.

G	Xǔ Shèn glosses this bound morpheme by the second member of the current compound of which it typically forms a part: ji&o xìng 憿幸. This is a current practice in Shuōwén.
Duàn tacitly and gratuitously rewrites 幸 as ;. Apart from everything else this was not a
very pedagogical thing to do given that the character he uses is very arcane and even seasoned epigraphers cannot be expected to be familiar with it.
PIF     þ [古弔切; LH keu, OCM kiâu] phonetic in 憿 [古堯切; LH keu, OCM kiâu]
PR	Tóngyīn 同音 OCM and LH homophones.

154. SW 10B 408: 154; DXB 221 (10B 18a); GL vol 11: 4737; Duàn 510 (10B 41a); TKJ 1465; Ozaki vol. 5: 1057.


HG:
G: SSF: PIF: IQ:

SA:

° (懖),
善自用之意也。从心,
銛聲。
《商書》曰： 今汝懖懖。
[柃]，古文
从耳。


kuò,
has the meaning of 'having a tendency to be self-willed'. It has xīn ``heart'' as a semantic constituent,
xi+n (GSR 621a: ``sharp, ...") is the phonetic constituent.
The Book of Sh+ng (SSJZS: 169b) says:
"Now you are strong-willed people".
  is an ancient style graph that has \$r (as a semantic constituent).

G   All glosses specify meanings, and Xǔ Shèn rarely (6 times) makes this explicit in the way that he does in this gloss by the formula X之意. His standard way of saying that A means B is to say that A is B. It will be interesting to collect the cases where Xǔ Shèn's definitions involve explicit semantic terminology of this sort.
Duàn adds the character 歫 at the beginning of this line, basing himself on a quotation in a
Shūjīng commentary. His arguments for this are less than convincing.
PIF     銛[息廉切; LH siam, OCM slam or slem] phonetic in 懖 [古活切; LH kuât, OCM kwât]
PR	The initials and finals are different, f*i sh*ng 非聲
IQ	The SSJZS edition of Shū (Pán g*ng 盤庚) (18.7) writes: 今汝聒聒 ``Now you are (making a deafening noise=) clamouring.'' (tr. Karlgren).



***
Minor GREEDY SERIES (155-156)

155. SW 10B 408: 155; DXB 221 (10B 18a); GL vol 11: 4738b; Duàn 510 (10B 41b); TKJ 1466; Ozaki vol. 5: 1058.

HG:             (忨),
G:	貪也。
SSF: 从心,
PIF: 元聲。


wán,
is 'greedy'.
It has xīn ``heart'' as a semantic constituent,
yuán (GSR 257a: ``head, ...") is the phonetic constituent.

IQ:

《春秋傳》曰： The Springs and Autumns Tradition (i.e. Zuǒzhuàn) says:

忨歲而揀日。 ``... trifles about years, and desires (length of) days ...'' (tr.
Legge)
PIF 元 [愚袁切; LH ngyan, OCM ngwan or ngon] phonetic in 忨 [五換切 (《廣韻》: 五丸切); LH nguân(C), OCM ngwân(s)]
IQ The SSJZS edition of Zuǒzhuàn (Zhāo 昭 1.8) writes: 主民，翫歲而愒日，其與幾何 ? ``When the president of the people trifles about years, and desires (length of) days, he cannot endure long", thus replacing 忨 by 翫 and 揀 by 愒. The number of cases where the received texts have readings others than those quoted by Xǔ Shèn is remarkable. (tr. Legge)


156. SW 10B 408: 156; DXB 221 (10B 18a); GL vol 11: 4739b; Duàn 510 (10B 41b); TKJ 1466; Ozaki vol. 5: 1058, see
variants in WGY: 455.

HG: Z (惏),
DG: 河內之北謂貪曰惏。
SSF: 从心,
PIF: 林聲。


lán,
in the north of Hénèi they use lán for t+n ``greedy". It has xīn ``heart'' as a semantic constituent,
lín (GSR 655a: ``forest") is the phonetic constituent.

G    Instead of an ordinary gloss Xǔ Shèn provides a dialectal gloss. See footnote 16.1 in F&ngyán 1 (Zh\%u Zǔmó 1956: 5-6) for 晉魏河內之北謂恣曰殘，楚謂之貪. When Xǔ Shèn follows this procedure we must assume that the semantic gloss imbedded in the dialectal gloss contains the analytically relevant semantic information.
PIF     林[力尋切; LH lim, OCM r\#m] phonetic in 惏 [盧含切; LH l\#m, OCM r*+ m]
***



MORAL AND INTELLECTUAL CONFUSION SERIES (157-(158)-167)

157. SW 10B 408: 157; DXB 221 (10B 18a); GL vol 11: 4740; Duàn 510 (10B 41b); TKJ 1466; Ozaki vol. 5: 1058.


HG:	(䮯),
G:	不明也。
SSF: 从心,
PIF: 夢聲。

mèng, ``be confused"
is 'to fail to understand'. [[is '(a way of) failing to under- stand (scil. as a continuing state of confusion)'.]]
It has xīn ``heart'' as a semantic constituent,
mèng (GSR 902a: ``dream") is the phonetic constituent.

G	Duàn gratuitously writes an archaising graph for 明, although what is at issue is the word and not the graph.
PIF     夢[莫忠切; LH mung(C), OCM m\#ng] phonetic in 䮯 [武亘切; LH m\#ng, OCM m*+ ng]

158. SW 10B 408: 158; DXB 221 (10B 18a); GL vol 11: 4740b; Duàn 510 (10B 41b); TKJ 1466; Ozaki vol. 5: 1059.

HG: l (愆),
G:	過也。
SSF: 从心,
PIF: 衍聲。


qi+n,
is 'transgression'. [[is 'a (a kind of) transgression (scil. a serious mistake that normally deserves punishment).]]
It has xīn ``heart'' as a semantic constituent,
y&n (GSR 197a: ``flow over, ...") is the phonetic constituent.

A:	[忓]，或in (the graph)	, [the word qi+n] is alternatively written with


SA:

从寒省。
    [ ֝ ] ， 籀文。

an abbreviated form of hán. is a large seal style graph.

PIF     衍[以淺切; LH kem(C), OCM kêms] phonetic in 愆 [去伲切; LH khian, OCM khjan]
PR	F*i sh*ng 非聲? Interestingly both the initials and the codae are consistent in this case but not the main vowels.


159. SW 10B 408: 159; DXB 221 (10B 18a); GL vol 11: 4741; Duàn 511 (10B 42a); TKJ 1467; Ozaki vol. 5: 1059.
HG: ƒ (慊), xián, ``suspect that"

G:	疑也。
SSF: 从心,
PIF: 兼聲。

is 'to be in doubt'. [[is '(a way of being in) doubt (scil. typically with an element of resentment)'.]]
It has xīn ``heart'' as a semantic constituent,
ji+n (GSR 627a: ``at the same time, ...") is the phonetic constituent.



HG Compare the modern Chinese xián 嫌 ``to suspect". It is indeed as if the word Xǔ Shèn's graph is referring to is the one written today with the woman radical. The current readings for the character Xǔ Shèn discusses are not traditionally linked to the meaning he attached to the graph.
PIF    兼[古甜切; LH kem, OCM kêm] phonetic in 慊 [戶兼切; LH khem(B), OCM khêm())]

INSERTED CHAOS SERIES (160-163) [Note also n° 167]

160. SW 10B 408: 160; DXB 221 (10B 18b); GL vol 11: 4741b; Duàn 511 (10B 42a); TKJ 1467; Ozaki vol. 5: 1060.
HG: [ (惑), huò, ``disorientated"

G:	亂也。
SSF: 从心,
PIF: 或聲。

is 'to be in chaos'. [[is '(a kind of) chaos (scil. of a mental kind)'.]] It has xīn ``heart'' as a semantic constituent,
huò (GSR 929a: ``territory, ...") is the phonetic constituent.

G     As in the preceding lexical entry, Xǔ Shèn provides an entirely relevant hypernym for the word he glosses. It is implausible that he was unaware of the fact that such glosses could not serve as near-equivalents but were designed to identify the genus of which the concept in question was a species, much in the spirit of Aristotle's theory of definitions.
PIF     或[于逼切; LH Gu\#k, OCM w*+ k] phonetic in 惑 [胡國切; LH Gu\#k, OCM w*+ k]
PR	Tóngyīn 同音 OCM and LH homophones.

161. SW 10B 408: 161; DXB 221 (10B 18b); GL vol 11: 4742; Duàn 511 (10B 42a); TKJ 1467; Ozaki vol. 5: 1060, see
variants in WGY: 455.


HG:
G: SSF: PIF:

  (怋),
怓也。从心, 民聲。


mín,
is 'to be befuddled'.
It has xīn ``heart'' as a semantic constituent,
mín (GSR 457a: ``people") is the phonetic constituent.

G   Xǔ Shèn provides an obscure gloss here for which he gives a clarifying explanation in the lexical entry that follows. The conceptual connection between 'stupidity' and the 'people' is evident in many parts of the vocabulary, and it is idiomatically manifest in the idiom yúmín 愚民 which means not 'of the people the stupid ones', but 'the people, who are stupid'.
PIF 民[彌鄰切; LH min, OCM min] phonetic in 怋 [呼昆切(《廣韻》: 彌鄰切); LH hu\#n, UCM hm*+ n]



162. SW 10B 408: 162; DXB 221 (10B 18b); GL vol 11: 4742b; Duàn 511 (10B 42a); TKJ 1467; Ozaki vol. 5: 1060.


HG:
G:

SSF: PIF: IQ:

g (怓),
亂也。
从心, 奴聲。
《詩》曰： 以謹䮣怓。


náo, ``be in mental chaos"
is 'chaos'. [[is '(a way of being) in chaos (scil. of a mental kind)'.]]
It has xīn ``heart'' as a semantic constituent,
nú (GSR 94l: ``slave, ...") is the phonetic constituent.
The Book of Odes (SSJZS: 548b) says:
"and so make the turbulent and obstreperous careful.'' (tr. Karlgren)

G Here again Shuōwén provides a necessary contextually motivated interpretation of a word that has just been used as a gloss.
IQ The illustration from Shījīng serves two entries at the same time and demonstrates that the individual entries in Shuōwén should not be read outside the context in which they occur. Xǔ Shèn disregards the distinction between these two graphs: 怋 and 䮣. In this instance he does come close to writing as if his concern was not with a graph but with a word (see n° 161).
Shījīng (Mín Láo 民勞 253.2).
PIF     奴[人諸切; LH nâ, OCM nâ] phonetic in 怓 [女交切; LH nau or ncu ?, OCM ?]
PR	F*i sh*ng 非聲?

163. SW 10B 408: 163; DXB 221 (10B 18b); GL vol 11: 4743; Duàn 511 (10B 42b); TKJ 1467; Ozaki vol. 5: 1061, see
variants in WGY: 455.


HG:
G:

SSF: PIF:

IQ:

h (惷),
亂也。
从心, 春聲。
《春秋傳》曰： 王室日惷惷焉。


chǔn,
is 'chaotic'. [EP: is (a way of being) chaotic (scil. men- tally, typically through incapacity or stupidity)'.]]
It has xīn ``heart'' as a semantic constituent,
chūn (GSR 463a: ``spring time") is the phonetic con- stituent.
The Springs and Autumns Tradition (i.e. Zuǒzhuàn, SSJZS: 2106a) says:
"The royal house is getting daily more confused."

CAS: 一曰:
厚也。

An (alternative) source says:
it is 'generous' (?)

PIF     Duàn rewrites 春 with an archaising graph, as if there was a convention in Shuōwén to the effect that phonetic constituents are listed in their seal form.



春[昌純切; LH tshhuin, OCM thun] phonetic in 惷 [尺允切; LH tshhuinB, OCM thun)]
IQ	The SSJZS edition of the Zuǒzhuàn (Zhāo 昭 24.6) has: 今王室實蠢蠢焉, ``The royal House is now indeed shaking". (tr. Legge).
CAS Duàn gratuitously rewrites 厚 with an archaising graph, although what is at issue is the word and not the graph.
***

164. SW 10B 408: 164; DXB 221 (10B 18b); GL vol 11: 4744; Duàn 511 (10B 42b); TKJ 1468; Ozaki vol. 5: 1061, see
variants in WGY: 456.
HG: _ (惛), hūn,
G:	不䓴也。 is 'not to understand'. [[is '(a way of) not understanding (scil. as a constant intellectual limitation)'.]]

SSF: 从心,
PIF: 昬聲。

It has xīn ``heart'' as a semantic constituent,
hūn (GSR 457j: ``dusk, ...") is the phonetic constituent.

PIF Duàn tacitly rewrites 昬 as 昏. Note that the character mín 民 being a Táng taboo was commonly replaced by shì 氏 in orthography. The semantic appositeness of hūn 昬 ``dusk'' must have been plain enough for Xǔ Shèn.
昬[武巾切; LH hu\#n, OCM hm*+ n] phonetic in 惛 [呼昆切; LH hu\#n, UCM hm*+ n]
PR   Tóngyīn 同音 OCM and LH homophones.

165. SW 10B 408: 165; DXB 221 (10B 18b); GL vol 11: 4744; Duàn 511 (10B 42b); TKJ 1468; Ozaki vol. 5: 1061.


HG:
G: SSF: PIF:

  (忥),
癡皃。从心, 气聲。


xì,
descriptive of being moronic.
It has xīn ``heart'' as a semantic constituent,
qì (GSR 517a: ``Shuowen says: cloudy vapours (no text)") is the phonetic constituent.

PIF     气[去旣切; LH kh\$s, OCM kh\#s] phonetic in 忥 [許旣切; LH h\$s, OCM h\#s]

166. SW 10B 408: 166; DXB 221 (10B 18b); GL vol 11: 4744b; Duàn 511 (10B 42b); TKJ 1468; Ozaki vol. 5: 1062.


HG:	(怀),
G:	巎言不慧也。
SSF: 从心,

wèi,
is 'to talk in one's sleep and make no sense'. It has xīn ``heart'' as a semantic constituent,



PIF: 衞聲。       wèi (GSR 342a: ``to guard, ...") is the phonetic constituent.
G The expression mèng yán 巎言 ``talk in one's sleep'' might be compared to mèng yán 夢言 in Hán F*iz+ where the idea is that of disclosing one's secret in one's dream rather than of making no sense while speaking when asleep.
PIF     衞[于歲切; LH wes, OCM we(t)s] phonetic in 怀 [于歲切; LH wes, OCM we(t)s]
PR	Tóngyīn 同音 OCM and LH homophones.

167. SW 10B 408: 167; DXB 221 (10B 18b); GL vol 11: 4744b; Duàn 511 (10B 42b); TKJ 1468; Ozaki vol. 5: 1062.
HG: ‰ (䓰), kuì,

G:	亂也。
SSF: 从心,
PIF: 貴聲。

is 'chaos'. [[is '(a way of) being in chaos (scil. psychologically)'.]] It has xīn ``heart'' as a semantic constituent,
guì (GSR 540b: ``precious, ...") is the phonetic constituent.

G	Note that the commentator Yán Sh(gǔ 顏師古 comments on Hànshū 45/2181-(11): 䓰,心亂也.
PIF     貴[居胃切; LH kuis, OCM kus] phonetic in 䓰 [胡對切; LH ku/s, OCM krûs]
***

RESENTMENT SERIES (168-(177)-180)
168. SW 10B 408: 168; DXB 221 (10B 18b); GL vol 11: 4745; Duàn 511 (10B 42b); TKJ 1468; Ozaki vol. 5: 1062, see
variant in WGY: 456.


HG:
G:

; (忌), jì,
憎惡也。 is 'to resent/ or to hate evil'. [[is '(a way of) resenting (scil. typi- cally for perceived misconduct)'.]]

SSF: 从心,
PIF: 己聲。

It has xīn ``heart'' as a semantic constituent,
j( (GSR 953a: ``self, ...") is the phonetic constituent.

G  The synonym compound z*ngwù 憎惡 presumably serves to disambiguate the intended gloss 惡, but the question remains why z*ng 憎 would not have been a sufficient gloss for Xǔ Shèn's purpose. One is inclined to read this gloss as a colloquialism. On the other hand, there is nothing to prevent one from taking this as a verb-object construction in spite of the fact that the predominant use of the binome in pre-Buddhist times appear to have been as a synonym compound.
PIF      J+ 己[居擬切; LH k\$\#B, OCM k\#)] phonetic in jì 忌 [渠記切; LH g\$\#C, OCM g\#h]



169. SW 10B 408: 169; DXB 221 (10B 19a); GL vol 11: 4745b; Duàn 511 (10B 42b); TKJ 1468; Ozaki vol. 5: 1063.


HG:
G: SSF: PIF:

  (忿),
悁也。从心, 分聲。


fèn, ``pent-up anger'' is 'to be upset'.
It has xīn ``heart'' as a semantic constituent,
f*n (GSR 471: ``divide, ...") is the phonetic constituent.

G The gloss here is singularly unhelpful because it explains a fairly common word by an exceedingly rare one. However, the pattern here is one of adjacent hùzhù 互注, and what the present example shows is that the members of a hùzhù pair can be of widely different currency.
PIF    分[甫文切; LH pun, OCM p\#n] phonetic in 忿 [敷粉切; LH phunB/C, OCM ph\#n)/s]

170. SW 10B 408: 170; DXB 221 (10B 19a); GL vol 11: 4745b; Duàn 511 (10B 42b); TKJ 1468; Ozaki vol. 5: 1063.


HG:
G:

SSF: PIF:

№ (悁),
忿也。
从心, ø聲。

yu+n,
is 'pent-up anger'. [[is '(like) pent-up anger (scil. without the element of aggressive resentment)'.]]
It has xīn ``heart'' as a semantic constituent,
yu+n (GSR 228a: ``Shuowen: small worm (no text).") is the phonetic constituent.

CAS: 一曰:
憂也。

An (alternative) source says: It is 'to be worried'.

SA:

[⥍]，籀文。

is a large seal style graph.

HG This hùzhù 互注 pair is taken here as a pair of undistinguished synonyms, although 悁 is the much rarer and the more poetic word.
G  Duàn distinguishes between fèn 忿 and fèn 憤: 忿與憤義不同。憤以气盈爲義，忿以狷急爲義 ``The force of fèn 忿 and fèn 憤 is not the same, fèn 憤 refers to one being filled with the qì of anger, while fèn 忿 refers to being impatient and upset.'' This gloss illustrates one of the immense strength of Duàn Yùcái as a lexicographer. His commentary on Shuōwén is indeed the first dictionary which pays systematic attention to the need for distinction between synonyms and also to the all-important phenomenon of pregnant versus loose use of words. The distinction is between xī yán zh\% 析言者 and tǒng yán zh\% 統言者 (as under wò 䎵, SWJZZ p. 387). Elsewhere the distinction is between xī yán zh\% 析言者 and hún yán zh\% 渾言者 (as under jiàn 見, SWJZZ p. 407). This area of the systematisation of Chinese lexicography as Duàn's contribution has been justly praised in twentieth century celebrated.
PIF ø [烏玄切 (《廣韻》：烏縣切); LH ?uen, OCM ?wên] phonetic in 悁 [於緣切; LH ?yen, OCM ?wen]
CAS Duàn rewrites 憂 as ⤏.



171. SW 10B 408: 171; DXB 221 (10B 19a); GL vol 11: 4746b; Duàn 511 (10B 43a); TKJ 1469; Ozaki vol. 5: 1063.


HG:
G: SSF: PIF: CAS:

(⫇),
恨也。从心, 黎聲。一曰: 怠也。

lí,
is 'to hate'.
It has xīn ``heart'' as a semantic constituent,
lí (GSR 519k: ``numerious, ...") is the phonetic constituent. An (alternative) source says:
It is 'to be remiss'.

G   Duàn chooses to write hèn with the archaic graph ‰.
PIF    黎[郎奚切; LH lei, OCM rǐ] phonetic in ⫇ [郎尸切; LH li, OCM ri]
CAS As we have noted before, it often looks as if Xǔ Shèn gets into a periodic habit of providing alternative glosses for a series of characters, and he does this even when the alternative glosses contribute little substance to his graphological analysis.

172. SW 10B 408: 172; DXB 221 (10B 19a); GL vol 11: 4746b; Duàn 511 (10B 43a); TKJ 1469; Ozaki vol. 5: 1064.


HG:

2 (恚), huì, ``furious ``

G:

SSF: PIF:

恨也。
从心, 圭聲。

is 'to resent'. [[EP: (a way of) resenting (scil. with a strong admix- ture of acute anger)'.]]
It has xīn ``heart'' as a semantic constituent,
guī (GSR 879a: ``jade tablet as token of rank, ...") is the phonetic constituent.

G	For his own reasons, Duàn writes nù 怒 for hèn 恨. There is no need to enter into the details on this point.
PIF     圭[古畦切; LH kue, OCM kwê] phonetic in 恚 [於避切; LH ?yeC , OCM ?weh]

173. SW 10B 408: 173; DXB 221 (10B 19a); GL vol 11: 4747; Duàn 511 (10B 43a); TKJ 1469; Ozaki vol. 5: 1064.


HG:

\$ (怨), yuàn, ``resent (typically superiors)"

G:

SSF: PIF:

恚也。
从心,
夗聲。

is 'furious'. [[is '(a way of) being furious (scil. typically vis-à-vis those who are more powerful than oneself)'.]]
It has xīn ``heart'' as a semantic constituent,
yuàn (GSR 260a: ``Shuowen: to turn over in bed' (no text)") is the phonetic constituent.




SA:	[⟍]， 
古文。

is an ancient style graph.

PIF     夗[於阮切; LH )yanB, OCM )on)] phonetic in 怨 [於願切; LH )yanC, OCM )ons]

174. SW 10B 408: 174; DXB 221 (10B 19a); GL vol 11: 4747b; Duàn 511 (10B 43a); TKJ 1469; Ozaki vol. 5: 1064.
HG:             (怒), nù, ``angry"

G:	恚也。
SSF: 从心,
PIF: 奴聲。

is 'furious'. [[is '(a way of being) furious (scil. showing this openly)'.]]
It has xīn ``heart'' as a semantic constituent,
nú (GSR 941: ``slave, ...") is the phonetic constituent.

G Note that the word which Xǔ Shèn treats as the logically superordinate is huì 恚, which is not the most common in the synonym group: it is certainly much rarer than yuàn 怨 and nù 怒. Thus the terms that Xǔ Shèn treats as the more abstract or general ones are by no means always the commoner ones. (When Xǔ Shèn uses one and the same word to define a series of other words we often take this as an indication that he treated the defining word as being the abstract common denominator).
PIF    奴[人諸切; LH nâ, OCM nâ] phonetic in 怒 [乃故切; LH nâB/C, OCM nâ)/h]

175. SW 10B 408: 175; DXB 221 (10B 19a); GL vol 11: 4748; Duàn 511 (10B 43a); TKJ 1469; Ozaki vol. 5: 1065.

HG:  (䮭),
G:	怨也。
SSF: 从心,
PIF: 敦聲。
IQ:	《周書》曰：


duì, ``detest'' is 'resent'.
It has xīn ``heart'' as a semantic constituent,
duì (GSR 464p: ``a kind of sacrificial vessel, ...") is the phonetic constituent.
In The book of Zhǒu (SSJZS: 204b) it says:

"凡民罔不䮭。'' Speaking of the people in general, none of them do not detest them."
PIF	Duàn tacitly rewrites 敦 with the archaising character ㇫, as if there was a convention in
Shuōwén to the effect that phonetic constituents are listed in their seal form.
    敦 is a notoriously multivalent phonetic constituent. According to HYDCD this character has the following readings: dūn, duī, duì, dùn, dún, tún, tuán, di&o, dào, zhǔn.
    敦[都昆切。又，丁回切; LH tu\#n, OCM tûn/ LH tui, OCM tûi] phonetic in 䮭 [徒對切; LH du\#iC , OCM dûih]
PR	One must not jump to the conclusion that this is a f*i sh*ng 非聲.



IQ Xǔ Shèn quotes the lexicographically relevant part of Shūjīng (K&ng gào 康誥, SSJZS: 204b): 凡民自得罪寇攘姦Ñ殺越人于貨䯭不畏死,罔弗䮭 ``All people who draw guilt upon themselves, being robbers and thieves and villains and traitors, who kill and (overthrow =) destroy and go for (goods=) spoil, and are forceful and do not fear death, there are none who do not detest them.'' (tr. Legge). It is significant that Xǔ Shèn does not always mechanically quote a whole passage but sometimes limits himself to what matters in his argumentative context.
Duàn tacitly rewrites 䮭 with the archaising character ⬷.

176. SW 10B 408: 176; DXB 221 (10B 19a); GL vol 11: 4748b; Duàn 511 (10B 43a); TKJ 1469; Ozaki vol. 5: 1065.

HG:	(慍), yùn, ``be upset"

G:	怒也。
SSF: 从心,
PIF: 昷聲。

is 'to be angry'. [[is '(a way of being) angry (scil. internally, psy- chologically, without necessarily showing one's anger openly)'.]] It has xīn ``heart'' as a semantic constituent,
w*n (GSR 426a: ``Shuowen says: kind (no text), thus taking it to be the primary graph for 溫") is the phonetic constituent.


G	Duàn rewrites 怒 as 怨 on the basis of later quotations. We translate the reading in the received text.
PIF     昷[烏渾切; LH ?u\#n, OCM ?ûn] phonetic in 慍 [於問切; LH ?unC, OCM uns]

177. SW 10B 408: 177; DXB 221 (10B 19a); GL vol 11: 4749; Duàn 511 (10B 43b); TKJ 1470; Ozaki vol. 5: 1065.
HG: ª (惡), è, ``ugliness (moral and aesthetic)"

G:	過也。
SSF: 从心,
PIF: 亞聲。

is 'deviation'. [[is '(a kind of) deviation (scil. from standards of aesthetic or moral beauty)'.]]
It has xīn ``heart'' as a semantic constituent,
yà (GSR 805a: ``inferior, ...") is the phonetic constituent.

HG This entry is an intrusion into the ``resentment'' series.
G  Xǔ Shèn defines a meaning of 惡 that is unrelated to that of 憎, but right before he turns to 憎. This might suggest that occasionally he made his list of characters for treatment in his dictionary before actually deciding which meanings of the character he was to focus on. Nonethe- less Xǔ Shèn's gloss in this case remains profoundly puzzling.
PIF     亞[衣駕切; LH )aC, OCM )râh] phonetic in 惡 [烏各切; LH )âk, OCM )âk]



178. SW 10B 408: 178; DXB 221 (10B 19a); GL vol 11: 4749b; Duàn 511 (10B 43b); TKJ 1470; Ozaki vol. 5: 1066.
HG: ˜ (憎), z*ng,

G:	惡也。
SSF: 从心,
PIF: 曾聲。

is 'to hate'.
It has xīn ``heart'' as a semantic constituent,
z*ng (GSR 884a: ``remote; etc., ...") is the phonetic constituent.

G This gloss splendidly illustrates how Xǔ Shèn feels entirely free to use characters with meanings other than those identified in his glosses. He has just defined the graph 惡 as it used to write the word è ``bad", but in the present definition he feels free to use the same graph under the reading wù ``to hate".
PIF   曾[昨稜切; LH ts\#ng, OCM ts*+ ng] phonetic in 憎 [作滕切; LH ts\#ng, OCM ts*+ ng]

179. SW 10B 408: 179; DXB 221 (10B 19b); GL vol 11: 4749b; Duàn 511 (10B 43b); TKJ 1470; Ozaki vol. 5: 1066.


HG:
G: SSF: PIF:

IQ:

(彡),
恨怒也。从心,
巿聲。
《詩》曰： 視我彡彡。

pèi,
is 'hateful anger'.
It has xīn ``heart'' as a semantic constituent,
pò (GSR 501a: ``SW: abundant vegetation") is the phonetic constituent.
The Book of Odes (SSJZS: 497a) says: ``They look at me with resentful anger".

G Again Duàn has the archaising graph ‰ for hèn 恨 and continues this idiosyncratic practice below.
    One is inclined to regard Xǔ Shèn's binome here as a Hàn dynasty colloquialism, but at the same time this discursive definition does make explicit the semantic nuance Xǔ Shèn thought the defined word expressed.
PIF    巿[普活切; LH put, OCM p\#t] (TKJ: 831) phonetic in 彡 [蒲昧切; LH phus, OCM ph\#s]
IQ The SSJZS edition of Shījīng (Bái huá 白華 229.5) writes: 視我邁邁 ``you look at me with disfavour.'' (tr. Karlgren).


180. SW 10B 408: 180; DXB 221 (10B 19b); GL vol 11: 4750b; Duàn 511 (10B 43b); TKJ 1470; Ozaki vol. 5: 1066, see
variants in WGY: 456.


HG:
G: SSF:

(❰),
怒也。从心,

yì,
is 'angry'.
It has xīn ``heart'' as a semantic constituent,




PIF: SPI:

刀聲。 讀若顡。

d+o (GSR 1131a: ``knife") is the phonetic constituent. Pronounced like yì.

HG Dictionary word, no received text.
PIF 刀 [都牢切; LH tâu, OCM tâu] phonetic in ❰ [魚旣切; LH ng\$s or ng\$ih, OCM ng\#s or ng\#ih]
    Duàn finds that 聲 is corrupt and omits the character. According to him the text read 从心刀, in which case Xǔ Shèn would have declared this character to by composed of two semantic constituents only.
PR     F*i sh*ng 非聲.
SPI Having declared d&o 刀 to be phonetic (an extraordinary statement if ever there was one, and one vividly understands why Duàn was tempted to throw it out) the text goes on to explain the reading of the head graph by an excruciatingly rare character which is notorious for its four equally uncommon readings.
    yì ❰ [魚旣切; LH ng\$s or ng\$ih, OCM ng\#s or ng\#ih] read like yì 顡 [DXB 9A 5b; TKJ: 1213: 五怪切] LH 3\$s or 3\$iC ?, OCM 3\#s ? or 3\#ih (< 3\#ls)?
                      *** DISSATISFACTION/RESENTMENT SERIES (181-189)
181. SW 10B 408: 181; DXB 221 (10B 19b); GL vol 11: 4751; Duàn 511 (10B 43b); TKJ 1470; Ozaki vol. 5: 1067.


HG:
G: SSF: PIF: SPI:

(⥻),
怨恨也。从心,
彖聲。 讀若膎。

xié,
is 'resentful hatred'.
It has xīn ``heart'' as a semantic constituent,
sh( ("Shuowen: pig") is the phonetic constituent. Pronounced like xié.

HG      [臣鉉等曰：彖非聲，未详。 ``Xuán and the others comment: sh+ is a wrong phonetic, we have no explanation."]
PIF 彖 [式視切; LH shiB, OCM lhi) or lh\#i) or hji), or hj\#i) (the OC initial consonant is very uncertain)(《廣韻》: 尺氏切) LH t´shieB (or tshhieB), OCM the) or k-hle) (the OC initial consonant is very uncertain)] phonetic in ⥻ [戶佳切; LH g/ or 4/, OCM grê]
PR     F*i sh*ng 非聲.
SPI As in the preceding entry the dúruò graph in this case xié 膎 is so rare that Xǔ Shèn's phonetic gloss cannot have been very useful to anyone.
Xié ⥻ [戶佳切; LH g/ or 4/, OCM grê] read like xié 膎 [戶皆切; LH h\$ap, OCM hap



182. SW 10B 408: 182; DXB 221 (10B 19b); GL vol 11: 4752; Duàn 511 (10B 44a); TKJ 1470; Ozaki vol. 5: 1067.
HG: 7 (恨), hèn,

G:	怨也。
SSF: 从心,
PIF: 艮聲。

is 'to resent'. [[is '(a kind of) resentment (scil. intense and, unlike
yuàn 怨, not normally directed towards superiors)'.]] It has xīn ``heart'' as a semantic constituent,
gèn (GSR 416a: ``obstinate, ...") is the phonetic constituent.

PIF     艮[古恨切; LH k\#nC, OCM k*+ ns] phonetic in 恨 [胡艮切; LH g\#nC, OCM g*+ ns]

183. SW 10B 408: 183; DXB 221 (10B 19b); GL vol 11: 4752; Duàn 512 (10B 44a); TKJ 1471; Ozaki vol. 5: 1067, see
variants in WGY: 457.


HG:
G: SSF: PIF:

± (懟),
怨也。从心, 對聲。


duì,
is 'to resent'.
It has xīn ``heart'' as a semantic constituent,
duì (GSR 511a: ``respond, ...") is the phonetic constituent.

HG Duàn comments: 今與䮭音義皆同，謂爲一字 ``Today duì 懟 and duì 䮭 are pronounced the same and have the same meaning, one considers them as the same word zì 字.'' This shows very nicely how in Duàn's language the concept zì 字 can already detached from that of the written graph and refers to a monosyllabic word. This shows that the concept of a (monosyllabic) word was current in China before the introduction of western linguistics and was clearly separated from the notion of a character even though the technical term used was the same for both in Duàn's time. The origin of Chinese notions of the word as opposed to the character deserved to be traced carefully in the history of Chinese philology. For polysyllabic expressions the notion of a word as opposed to that of by-syllabic compound of characters became relevant particularly in the context of liánmiánzì 連綿字 which are notorious for their orthographic indeterminacy: Cít\$ng 辭通 by Zhū Qǐfèng 朱起鳳. (1875-1948) is a compilation that is particularly concerned with the varying representations of one and the same binome.
PIF    對[都隊切; LH tu\#s, OCM tûts] phonetic in 懟 [丈淚切; LH druis, OCM druts]

184. SW 10B 408: 184; DXB 221 (10B 19b); GL vol 11: 4752; Duàn 512 (10B 44a); TKJ 1471; Ozaki vol. 5: 1068, see
variants in WGY: 457.


HG:
G:

F (悔), hu(,
悔恨也。 is '(remorseful) resentment'. [[is '(a kind of) (remorseful) resent- ment (scil. of oneself)'.]]




SSF: 从心,
PIF: 每聲。

It has xīn ``heart'' as a semantic constituent,
m\$i (GSR 947i: ``flourishing, ...") is the phonetic constituent.

G      The repetition of the head graph, a practice which may derive from the fact that Xǔ Shèn or his collaborators was mechanically copying a gloss from his source, in this case allows for an interpretation along the lines of our EP.
PIF   每[荒內切; LH m\#B, OCM m*+ )] phonetic in 悔 [荒內切; LH hu\#h, OCM hm*+ h]

185. SW 10B 408: 185; DXB 221 (10B 19b); GL vol 11: 4752b; Duàn 512 (10B 44a); TKJ 1471; Ozaki vol. 5: 1068.


HG:	(⣸),
G:	小怒也。
SSF: 从心,
PIF: 䐻聲。

chì, ``irritation'' is 'slight anger'.
It has xīn ``heart'' as a semantic constituent,
zhù (GSR 127a: ``post or stand on which to suspend musical instruments, ...") is the phonetic constituent.

G    In this gloss Xǔ Shèn does the work of our Explanatory Paraphrase for us. He could simply have written 怒也. In this case, it is our quiet hope that we might have been able to remember to add the EP ``is a (kind of) anger (scil. of slight intensity)". The tragic fact is that many of the characters which Xǔ Shèn discusses are so rarely used that we are unable to do the work of the EP for him. And there is a lingering doubt that Xǔ Shèn himself would have been uncertain of the semantic nuances of the words he discusses. This comes out most clearly in the many instances where he lines up competing glosses for arcane characters without taking a clear stand on whether he agrees with any or all of them.
PIF     䐻[中句切; LH troC , OCM troh] phonetic in ⣸ [充世切; LH tshias, OCM ta(t)s]
PR	F*i sh*ng 非聲.

186. SW 10B 408: 186; DXB 221 (10B 19b); GL vol 11: 4753; Duàn 512 (10B 44a); TKJ 1471; Ozaki vol. 5: 1069.


HG:
G: SSF: PIF:

  (怏),
不服懟也。从心,
央聲。


yàng, ``indignant"
is 'recalcitrant resentment'.
It has xīn ``heart'' as a semantic constituent,
y+ng (GSR 718a: ``centre, ...") is the phonetic constituent.

G Negative adjectives are not common in classical Chinese, and also fairly rare in Shuōwén. We suspect a mild colloquialism here. One notes, however, that we have here a sequence of analytic glosses from n° 184 to 186.



    Duàn refuses to contemplate negative adjectives and rewrites the text adding 也 after 不服. He then congratulates himself on his emendation with the remark 奪一也字，遂不可解矣 ``if one takes away just the one character y\% than this cannot be interpreted anymore."
PIF   央[於良切; LH ?\$âng, OCM ?ang] phonetic in 怏 [於亮切; LH ?\$ângB/C, OCM ?ang?/h]

187. SW 10B 408: 187; DXB 221 (10B 19b); GL vol 11: 4753; Duàn 512 (10B 44b); TKJ 1471; Ozaki vol. 5: 1069.


HG:
G: SSF1: SSF2:

² (懣), 煩也。从心,
从滿。


mèn, ``encumbered'' is 'to feel vexed'.
It has xīn ``heart'' as a semantic constituent,
it (also) has m&n ``... full ...'' as a semantic constituent.

SSF2 Duàn tacitly rewrites this as 从心滿, as if the formula the received text uses was not current, and he (rightly) points out that m(n 滿 is also phonetic.
HP	滿 [莫旱切;  LH mânB, OCM mân)]  not declared phonetic in 懣 [莫困切; LH m\#nC, OCM m*+ ns]


188. SW 10B 408: 188; DXB 221 (10B 19b); GL vol 11: 4753b; Duàn 512 (10B 44b); TKJ 1471; Ozaki vol. 5: 1069.
HG: ž (憤), fèn, ``full of pent-up anger"

G:	懣也。
SSF: 从心,
PIF: 賁聲。

is 'encumbered'. [[is '(a way of) being encumbered (scil. by pent- up resentment)'.]]
It has xīn ``heart'' as a semantic constituent,
b*n (GSR 437a: ``ardent, ...") is the phonetic constituent.

G	For very good reasons Xǔ Shèn avoids the standard association of fèn 憤 with nù 怒. Open anger is very different from concealed festering resentment.
PIF     賁[彼義切; LH pu\#n, OCM p*+ n]?? phonetic in 憤 [房吻切; LH bunB, OCM b\#n)]
PR	According to the Shuōwén fǎnqiè of Táng date this would turn out to be a f*i sh*ng 非聲. However, the current reading b*n will solve the problem.


189. SW 10B 408: 189; DXB 222 (10B 20a); GL vol 11: 4754; Duàn 512 (10B 44b); TKJ 1471; Ozaki vol. 5: 1070.
HG: P (悶), mèn,
G:	懣也。	is 'encumbered'. [[is '(a way of being) encumbered (scil. by a sense of being closed-in and without hope)'.]]




SSF: 从心,
PIF: 門聲。

It has xīn ``heart'' as a semantic constituent,
mén (GSR 441a: ``gate, ...") is the phonetic constituent.

G It may be objected that mèn 懣 is a fairly rare word which may appear to be of little explana- tory usefulness. However upon second thought it seems quite as likely that Xǔ Shèn has discov- ered for himself that this rare word denominates, as it were, a semantic common denominator of a series of words. One hastens to add that Shuōwén is not a book in which such systematic semantic analysis predominates. On the other hand, it appears that in the context of his graphological dictionary Xǔ Shèn does introduce semantic series by placing them together and providing them with systematically related definitions.
PIF     門[莫奔切; LH m\#n, OCM m*+ n] phonetic in 悶 [莫困切; LH m\#nC, OCM m*+ ns]
***

190. SW 10B 408: 190; DXB 222 (10B 20a); GL vol 11: 4754; Duàn 512 (10B 44b); TKJ 1472; Ozaki vol. 5: 1070.

HG: X (惆),
G:	失意也。
SSF: 从心,
PIF: 周聲。


chóu,
is 'to have lost one's orientation'> 'be depressed'. It has xīn ``heart'' as a semantic constituent,
zhǒu (GSR 1083a: ``all round...") is the phonetic constituent.

HG Chóu is discussed immediately before the second member of the compound in which it is most common: chóuchàng 惆悵.
G	In cases like these it is as if Xǔ Shèn chooses to provide colloquial glosses when in fact a wide range of more literary classical idioms would have been available.
PIF     周[職䛔切; LH tshu, OCM tiu] phonetic in 惆 [敕鳩切; LH thru, OCM thru or thriu]

191. SW 10B 408: 191; DXB 222 (10B 20a); GL vol 11: 4754b; Duàn 512 (10B 44b); TKJ 1472; Ozaki vol. 5: 1070.


HG:
G: SSF: PIF:

O (悵),
望恨也。从心,
長聲。


chàng,
is 'to be full of self-recriminatory resentment'. It has xīn ``heart'' as a semantic constituent,
cháng (GSR 721a: ``long, ...") is the phonetic constituent.

HG It is not a coincidence that the members of the binome chóuchàng 惆悵 are dealt with adjacently. However, in this case Xǔ Shèn does not treat these words as bound forms by any of his standard patterns such as defining the members by the compound. Indeed, there are some examples of independent use of the constituent characters.
G    Duàn also rewrites wàng 望 in his archaising style.



PIF     長[直良切; LH drâng, OCM drang] phonetic in 悵 [丑亮切; LH thrângC, OCM thrangh]

192. SW 10B 408: 192; DXB 222 (10B 20a); GL vol 11: 4754b; Duàn 512 (10B 44b); TKJ 1472; Ozaki vol. 5: 1070, see
variants in WGY: 457.

HG:	{ (愾),
G:	大息也。
SSF1: 从心,
SSF2: 从氣，


xì,
is 'to heave a deep sigh'.
It has xīn ``heart'' as a semantic constituent, it (also) has qì as a semantic constituent.

PIF: IQ:

氣亦聲。
《詩》曰： 愾我寤歎。

Xì (GSR 517c: ``to present food (Tso ap. Shuowen)") is at the same time the phonetic constituent.
The Book of Odes (SSJZS: 386a) says: ``moaning I awake and sigh.'' (tr. Karlgren)

G	Duàn claims that all editions have 太 for 大, but the Dà Xú b\%n does not.
Duàn rewrites 也 as mào 皃.
PIF     氣[許旣切; LH kh\$s, 0CM kh\#(t)s] phonetic in 愾 [許旣切; LH kh\$s, OCM kh\#(t)s]
PR	Tóngyīn 同音 OCM and LH homophones.
IQ	The SSJZS edition of Shījīng (Xià quán下泉 153.3) has the mouth radical in our graph 嘆: 愾我寤嘆 ``moaning I awake and sigh."

193. SW 10B 408: 193; DXB 222 (10B 20a); GL vol 11: 4755b; Duàn 512 (10B 44b); TKJ 1472; Ozaki vol. 5: 1071.

HG: ¬ (懆),
G:	愁不安也。
SSF: 从心,


c&o,
is 'to be restless because of low spirits'.
It has xīn ``heart'' as a semantic constituent,

PIF: IQ:

‰聲。
《詩》曰： 念子懆懆。

zào (GSR 1134a: ``Shuowen says: a crowd of birds chirping (no text)") is the phonetic constituent.
The Book of Odes (SSJZS: 497a) says: ``When I think of you I get all depressed"

G	This is the third composite gloss in a row.
PIF     ‰ [穌到切; LH sâuC, OCM sâuh] phonetic in 懆 [七早切; LH ts-hâuB, OCM ts-hâu) < OCM k-sâu]
IQ	Shījīng (Bái huá 白華 229.5) tr. Karlgren: ``I think of you and I am grieved.



***
DEJECTION SERIES (194-205)

194. SW 10B 408: 194; DXB 222 (10B 20a); GL vol 11: 4756b; Duàn 512 (10B 45a); TKJ 1472; Ozaki vol. 5: 1072.
HG:  x (愴), chuàng, ``disheartened"

G:	傷也。
SSF: 从心,
PIF: 倉聲。

is 'to be hurt'. [[is '(a way of) being hurt (scil. feeling hurt, psycho- logically)'.]]
It has xīn ``heart'' as a semantic constituent,
c+ng (GSR 703a: ``granary, ...") is the phonetic constituent.

G      Sh&ng 傷 basically means ``to injure/hurt'' and ``to be injured/hurt". Our translation takes Xǔ Shèn to subsume the concept of being disheartened under the more general concept of being hurt. However, we are well aware that sh&ng 傷 by itself is also current in the meaning ``feel (psycho- logically) hurt". But perhaps this derived meaning is not the one that is relevant here.
PIF 倉 [七岡切; LH ts-hâng, OCM ts-hâng] phonetic in 愴 [初亮切; LH ts-hrângC , OCM ts- hrângh]


195. SW 10B 408: 195; DXB 222 (10B 20a); GL vol 11: 4757; Duàn 512 (10B 45a); TKJ 1472; Ozaki vol. 5: 1072.

HG: \% (怛),
G:	憯也。
SSF: 从心,
PIF: 旦聲。


dá, ``dejected'' is 'to be sad'.
It has xīn ``heart'' as a semantic constituent,
dàn (GSR 149a: ``dawn, ...") is the phonetic constituent.

A:	[⟌ ]，或从in the graph , [the word dá] is alternatively written with he


IQ:

心在旦下。
《詩》曰： 信誓⟌⟌。

graphic constituent xīn ``heart'' under dàn. The Book of Odes (SSJZS: 325b) says:
"we were sworn to good faith (painfully =) earnestly.'' (tr. Karlgren)

PIF    旦[得案切; LH tânC, OCM tâns] phonetic in 怛 [得案切,又當割切; LH tât, OCM tât]
    The addition of the alternative fǎnqiè 當割切 spelling to the standard 得案切 spelling looks like a fine example of rationalising reconstructions of pronunciations. It appears that the observed reading has always been 得案切; but since this would create flagrant instances of f*i sh*ng 非聲in the xiésh*ng 諧聲 system, it looks as if Chinese philologists, hoping to explain the composition of Chinese characters, may have been tempted to postulate an alternative reading 當割切 which would open for a standard analysis of the phonetic constituents of the characters. A great deal of philological detective work will be needed to establish the crucial distinction between observation-




based records of how words are pronounced when written with a certain character on the one hand, and theory-based hypothesis on how graphs must have been read (even when they represent no known words of the spoken language) if they are to function coherently within the xiésh*ng 諧聲system.
PR     F*i sh*ng 非聲.
A   Surprisingly, there are few instances where Xǔ Shèn specifies what is above, below, to the left or to the right of what. The heart radical is notorious for having two variants: one on the position on the left, and the other below the rest of the graph. Only occasionally, this causes Xǔ Shèn to specify the position of the radical.
The illustrative quotation illustrates not the gloss but the documentation for the allograph.
IQ The SSJZS edition of Shījīng (58.6) has 信誓旦旦 ``we were sworn to good faith (painfully =) earnestly.'' (tr. Karlgren).


196. SW 10B 408: 196; DXB 222 (10B 20a); GL vol 11: 4758b; Duàn 512 (10B 45a); TKJ 1473; Ozaki vol. 5: 1073.
HG: ¤ (憯), c&n, ``distressed"

G:	痛也。
SSF: 从心,
PIF: 䕟聲。

is 'to suffer pain'. [[is '(a way of) suffering pain (scil. of a certain psychological kind)'.]]
It has xīn ``heart'' as a semantic constituent,
c&n (GSR 660c:"particle")(SW: ``at one time in the past") is the phonetic constituent.

HG This entry defines the term used to gloss the preceding entry. As we have seen before such interlocking structures in the arrangement of the entries in Shuōwén are common throughout the book.
PIF 䕟 [七感切; LH tshh\#mB, OCM tshh*+ m)] phonetic in 憯 [七感切; LH tshh\#mB, OCM tshh*+ m)]


197. SW 10B 408: 197; DXB 222 (10B 20a); GL vol 11: 4759; Duàn 512 (10B 45a); TKJ 1473; Ozaki vol. 5: 1073.
HG: ‹ (慘), c&n, ``depressed"

G:	毒也。
SSF: 从心,
PIF: 參聲。

is 'to be poisoned'. [[is '(a way of being) poisoned (scil. in a certain psychological way and figuratively)'.]]
It has xīn ``heart'' as a semantic constituent,
c+n (GSR 647a: ``triad, ...") is the phonetic constituent.



HG Near-synonyms that are homophonous must have been the source of great confusion in linguistic communication in ancient China as they continue to be in modern times. Suffice it to mention the example of cí 詞 and cí 辭.
PIF Shuōwén takes account of the graph c&n only as a variant of the head graph sh*n 曑 (see Shuōwén 7A 8b). Xǔ Shèn must have known that on the one hand the graph c&n has many more readings than the graph 曑 and that moreover the most common reading of 參 is not a possible reading for 曑. It appears that Xǔ Shèn got into this position of having to declare phonetic a graph that was not a head graph in his system but only a graphic variant, because the writing system at this time made productive use of the variant but not of what he regarded as the basic form of the graph.
    c&n 參 [所今切 Wáng Lì 2000: 倉含切/蘇甘切; LH tshh\#m, OCM tshh*+ m] phonetic in [七感切; LH tshh\#mB, OCM tshh*+ m)]
    The Shuōwén fǎnqiè for 參 is connected to the fact that Xǔ Shèn's analysis of the graph is based on the ancestor of the reading sh*n. As we have seen before, nothing suggests that the graphologically relevant readings assigned to a graph in Shuōwén should represent the phonetic value of the phonetic constituent.


198. SW 10B 408: 198; DXB 222 (10B 20b); GL vol 11: 4759; Duàn 512 (10B 45a); TKJ 1473; Ozaki vol. 5: 1073.


HG:

T(悽), qī, ``despondent"

G:	痛也。

SSF: 从心,

is 'to feel pain'. [[is '(a way of) feeling pain (scil. in a psychologi- cal way)'.]]
It has xīn ``heart'' as a semantic constituent,

PIF:

妻聲。

qī (GSR 592a: ``wife, ...") is the phonetic constituent.

G	After one interruption in n° 197, Xǔ Shèn reverts to the generic gloss tòng 痛也 which continues to dominate the dejection series from n° 199 to 204.
PIF     妻[七稽切; LH tshhei, OCM tshh*+ i] phonetic in 悽 [七稽切; LH tshheiC , OCM tshh*+ ih]

199. SW 10B 408: 199; DXB 222 (10B 20b); GL vol 11: 4759b; Duàn 512 (10B 45a); TKJ 1473; Ozaki vol. 5: 1073, see
variants in WGY: 457.


HG:
G:

AG:

9 (恫),
痛也。
一曰:
呻吟也。


tǒng, ``dejected"
is 'to feel pain'. [[is '(a way of) feeling pain (scil. in a psycho- logical way)'.]]
An (alternative) source says: It is 'to moan'.




SSF: PIF:

从心, 同聲。

It has xīn ``heart'' as a semantic constituent,
tóng (GSR 1176a: ``same, ...") is the phonetic constituent.

AG	Duàn tacitly moves this AG to the end of the entry.
PIF     同 [徒紅切; LH dong, OCM dông] phonetic in 恫 [他紅切; LH thong, OCM thông]

200. SW 10B 408: 200; DXB 222 (10B 20b); GL vol 11: 4760; Duàn 512 (10B 45b); TKJ 1473; Ozaki vol. 5: 1074.
HG: M (悲), b*i, ``sad"

G:	痛也。
SSF: 从心,
PIF: 非聲。

is 'to feel pain'. [[is '(a way of) feeling pain (scil. psycho- logically)'.]]
It has xīn ``heart'' as a semantic constituent,
f*i (GSR 579a: ``it is not, ...") is the phonetic constituent.

G    Under this heading, it is Duàn who does our work in EP: 按憯者，痛之深者也。恫者，痛之專者也。悲者，痛之上騰者也。各從其聲而得之 ``c(n is the deep pain, t\$ng is the focused pain, b*i is the soaring pain: in each case one gets the meaning from the phonetic constituent."
PIF   非 [芳微切; LH pui, OCM p\#i] phonetic in 悲 [府眉切; LH p\$, OCM pr\#i]

201. SW 10B 408: 201; DXB 222 (10B 20b); GL vol 11: 4760; Duàn 512 (10B 45b); TKJ 1473; Ozaki vol. 5: 1074.
HG: Ⓒ (惻), cè, ``despondent"

G:	痛也。
SSF: 从心,
PIF: 則聲。

is 'to feel pain'. [[is '(a way of) feeling pain (scil. psychologically, often in sympathy with someone else)'.]]
It has xīn ``heart'' as a semantic constituent,
zé (GSR 906a: ``rule, ...") is the phonetic constituent.

PIF     則[子德切; LH ts\#k, OCM ts*+ k] phonetic in 惻 [初力切; LH tshhr\$k, OCM tshhr\#k]

202. SW 10B 408: 202; DXB 222 (10B 20b); GL vol 11: 4760b; Duàn 512 (10B 45b); TKJ 1473; Ozaki vol. 5: 1074.
HG: g (惜), xī, ``wistful"

G:	痛也。
SSF: 从心,
PIF: 昔聲。

is 'to feel pain'. [[is '(a way of) feeling pain (scil. in a certain almost joyful way)'.]]
It has xīn ``heart'' as a semantic constituent,
xī (GSR 798a: ``anciently, ...") is the phonetic constituent.



PIF     Duàn tacitly rewrites xī 昔 with an archaising graph B, as if there was a convention in
Shuōwén to the effect that phonetic constituents are listed in their seal form.
昔 [思積切; LH siak, OCM sak] phonetic in 惜 [思積切; LH siak, OCM sak]
PR	Tóngyīn 同音 OCM and LH homophones.

203. SW 10B 408: 203; DXB 222 (10B 20b); GL vol 11: 4760b; Duàn 512 (10B 45b); TKJ 1474; Ozaki vol. 5: 1075.
HG:  o (愍), m(n, ``despondent"

G:	痛也。
SSF: 从心,
PIF: 敃聲。

is 'to feel pain'. [[is '(a way of) feeling pain (scil. psychologically and intensely)'.]]
It has xīn ``heart'' as a semantic constituent,
m(n (GSR 457G: ``strong (Shuowen, same as 457y)") is the pho- netic constituent.

PIF     敃[眉殞切; LH m\$nB, OCM mr\#n)] phonetic in 愍 [眉殞切; LH m\$n, OCM mr\#n]

204. SW 10B 408: 204; DXB 222 (10B 20b); GL vol 11: 4761; Duàn 512 (10B 45b); TKJ 1474; Ozaki vol. 5: 1075.
HG: € (慇), yīn, ``despondent"

G:	痛也。
SSF: 从心,
PIF: 殷聲。

is 'to feel pain'. [[is '(a way of) feeling pain (scil. psychologically and lastingly)'.]]
It has xīn ``heart'' as a semantic constituent,
yīn (GSR 448a: ``great, ...") is the phonetic constituent.

PIF     殷[於身切; LH ?\$n, OCM ?\#n] phonetic in 慇 [於巾切; LH ?\$n, OCM ?\#n]
PR	Tóngyīn 同音 OCM and LH homophones.

205. SW 10B 408: 205; DXB 222 (10B 20b); GL vol 11: 4761; Duàn 512 (10B 45b); TKJ 1474; Ozaki vol. 5: 1075.


HG:	(3),
G:	痛聲也。
SSF: 从心,

yī,
is 'an onomatopoeic word for feeling pain'. It has xīn ``heart'' as a semantic constituent,

PIF: IQ:

依聲。
《孝經》曰： 哭不3。

yī (GSR 550f: ``lean upon, ...") is the phonetic constituent. The Xiàojīng (SSJZS: 2561a) says:
"When lamenting he does not sob".



G  The last member of this long series of words defined by the general colourless tòng 痛 which we have attempted to render as colourlessly as possible by ``pained'' is deviant through its onomatopoeic character.
    The suffix sh*ng 聲 marking onomatopoeic elements is not frequent, but it remains important as a part of Xǔ Shèn's conceptual system.
PIF     依[於稀切; LH ?\$i, OCM \#i] phonetic in 3 [於豈切; LH ?\$i, OCM \#i]
PR	Tóngyīn 同音 OCM and LH homophones.
IQ	The SSJZS edition of Xiàojīng (S&ng qīn 喪親 18.1) has y+ 偯 instead of yī 3:孝子之喪親也,哭不偯 ``As for a filial son mourning for his parents, his lamentation is without sobbing".
***

206. SW 10B 408: 206; DXB 222 (10B 20b); GL vol 11: 4762; Duàn 513 (10B 46a); TKJ 1474; Ozaki vol. 5: 1076, see
variants in WGY: 458.


HG:
G: SSF:
PIF:

SPI:

(梪),
梪，存也。从心,
簡省聲。
讀若簡。

ji&n,
Ji&n is 'to investigate'.
It has xīn ``heart'' as a semantic constituent，
It has an abbreviated form of ji&n (GSR 191d: ``slip or tablet of bamboo, for writing") as the phonetic.
pronounced like ji&n.

HG One might, obviously, look for this character under the bamboo radical, as it is indeed classified by the K&ngxī zìdi(n system. Xǔ Shèn consistently aims at classifying characters by taking careful account of what he considers relevant meanings of the words these characters are used to write.
G Xǔ Shèn rarely repeats the head graph in the gloss, and when he does he gives the impression that he is simply quoting from a commentary or handbook. The gloss takes the form of a complete quotation of another gloss. And moreover, the definiens is taken in a highly unusual way.
Duàn rewrites this as 梪梪,在也. According to his well founded emendation in the text (see
)ry( Shì xùn 釋訓曰：存存、梪梪，在也, which reads in fact: 存存、萌萌，在也 (Xú Cháohuá 1994).) This entry in Shuōwén remains opaque to us.
PIF     簡[古限切; LH k/nB, OCM krên)] phonetic in 梪 [古限切; LH k/nB, OCM krên)]
PR	Tóngyīn 同音 OCM and LH homophones.
SPI     Ji(n 梪 [古限切; LH k/nB, OCM krên)] read like ji(n 簡 [古限切][LH k/nB, OCM krên)]



             *** ``MOVEMENT'' SERIES (207-209)

207. SW 10B 408: 207; DXB 222 (10B 20b); GL vol 11: 4763b; Duàn 513 (10B 46a); TKJ 1474; Ozaki vol. 5: 1077.


HG:

ƒ (慅), s+o, ``anxious"

G:

SSF: PIF: CAS:

動也。
从心, 蚤聲。一曰: 起也。

is 'to be moved'. [[is '(a way of being) moved (scil. psychologi- cally so as to become agitated and unsettled)]].
It has xīn ``heart'' as a semantic constituent,
z&o (GSR 1112d: ``flea, ...") is the phonetic constituent. An (alternative) source says:
is 'to arouse'.

PIF     蚤 [子皓切; LH tsouB, OCM tsû)] (graphic variant of z(o 畳 in Shuōwén; TKJ: 1930) phonetic in 慅 [穌遭切; LH sou(B), OCM sû())]

208. SW 10B 408: 208; DXB 222 (10B 21a); GL vol 11: 4763; Duàn 513 (10B 46a); TKJ 1475; Ozaki vol. 5: 1077.


HG:

u (感),


g&n, ``cause to be moved/changed"

G:	動人心也。 is 'to move a person's heart/mind'.

SSF: PIF:

从心, 咸聲。

It has xīn ``heart'' as a semantic constituent,
xián (GSR 671a: ``all, ...") is the phonetic constituent.

G Note the anthropocentric perspective which we already know from the Shuōwén entry xīn 心itself. Xǔ Shèn was very interested in cosmology and he was certainly familiar with the theory of resonance (g(n 感) in Chinese cosmology (Le Blanc 1985).
Xǔ Shèn's gloss here specifically links the character g(n 感 to the human heart in accordance
with his definition of the xīn 心 radical.
PIF  咸[胡監切; LH g/m, OCM gr*+ m] phonetic in 感 [古䲤切; LH k\#mB, OCM k*+ m)]

209. SW 10B 408: 209; DXB 222 (10B 21a); GL vol 11: 4764; Duàn 513 (10B 46b); TKJ 1475; Ozaki vol. 5: 1077, see
variants in WGY: 458.


HG:
G:

(忧),
不(read 心)動也。

yòu,
is 'to be moved in one's heart'.

SSF: 从心,
PIF: 尤聲。

It has xīn ``heart'' as a semantic constituent,
yóu (GSR 996a: ``fault, ...") is the phonetic constituent.



SPI: 讀若祐。	pronounced like yòu.
G	Duàn 's emendation reading xīn 心 for bù 不 is tempting. It does seem that the text as it stands is corrupt. Our translation follows Duàn's conjecture.
PIF     尤[羽求切; LH wu, OCM w\#] phonetic in 忧 [于救切; LH wuC , OCM w\#h]
SPI     Yòu 忧 [于救切; LH wuC , OCM w\#h] read like yòu 祐[于救切; LH wuC, OCM w\#h]
***

210. SW 10B 408: 210; DXB 222 (10B 21a); GL vol 11: 4764; Duàn 513 (10B 46b); TKJ 1475; Ozaki vol. 5: 1078, see
variants in WGY: 458.


HG:
G: SSF: PIF:

(⢈),
怨仇也。从心,
咎聲。

qiú,
is 'resentful hostility'.
It has xīn ``heart'' as a semantic constituent,
jiù (GSR 1068a: ``blame, ...") is the phonetic constituent.

G	Duàn rewrites 仇 as ⢈, on the basis of Gu(ngyùn.
PIF     咎[其久切; LH guB, OCM gu)] phonetic in ⢈ [其久切(《廣韻》: 巨鳩切); LH kuB or guB, OCM ku) or gu)]


***
WORRY SERIES (211-233)

211. SW 10B 408: 211; DXB 222 (10B 21a); GL vol 11: 4764b; Duàn 513 (10B 46b); TKJ 1475; Ozaki vol. 5: 1078.



HG:
G: SSF: PIF:

(䮦),
憂皃。从心, 員聲。

yún,
descriptive of worry.
It has xīn ``heart'' as a semantic constituent,
yuán (GSR 227a: ``circumference, ...") is the phonetic con- stituent.

G    Duàn again rewrites 憂 as ⤏.
    It is at this point that Duàn comments on his rewriting of the character throughout this section: 各本皆作憂，淺人用俗行字改之也。許造此書，依形立解，斷非此形彼義、牛頭馬脯，以自爲矛盾者 (…) 他書可用叚借，許自爲書不可用叚借 ``All editions write 憂, but a superficial person used the vulgar current character to rewrite the text. When Xǔ created this book, once he has established an interpretation on the basis of the shape/graph, he certainly did not used




this shape/graph for different meanings. 'The head of an ox and the torso of a horse', thus contradicting himself. (…)'' A little further on Duàn elaborates his view: ``Other books may use loan characters, Xǔ Shèn in making his book does not use loan characters.'' Duàn comes com- mendably clear on this point, and he is manifestly mistaken. For one thing the idea that Xǔ Shèn 'established an interpretation on the basis of the shape/graph' is profoundly misleading if it is to suggest that what the glosses establish are the basic meaning of the words written by the characters. To make our point entirely explicit: Xǔ Shèn is not using characters only in the meaning which he takes to be relevant to their graphological structure. To repeat a striking example, he never uses qí 其 to refer to a basket, nor does he use the particle suǒ 所 to refer to the sound of an axe hitting a tree as he should have done if Duàn had been right on this very basic point.
    The hypernym of this series y\$u 憂 cannot come at its head because it is construed in Shuōwén as not having the xīn 心 radical in the first place. The graph y\$u ⤏ is defined first towards the end of this radical (n° 232).
PIF    員[王權切; LH wian, OCM wen] phonetic in 䮦 [王分切; LH wun, OCM w\#n]

212. SW 10B 408: 212; DXB 222 (10B 21a); GL vol 11: 4764b; Duàn 513 (10B 46b); TKJ 1475; Ozaki vol. 5: 1078.


HG:
G: SSF: PIF:

(怮),
憂皃。从心, 幼聲。

yǒu,
descriptive of worry.
It has xīn ``heart'' as a semantic constituent,
yòu (GSR 1115f: ``young") is the phonetic constituent.

G    Duàn again rewrites 憂 as ⤏.
    In this series Xǔ Shèn begins with the descriptive terminology defined in terms of the linguistic technical mào 皃 ``descriptive of". In the following items the definition changes from 憂皃 to 憂也, and we note that there is no such formula as 憂皃也. It is as if 皃 takes the place of the 也, creating a statement of a different analytic type. (Sūn Liángmíng 孫良明 2005: 54).
PIF     幼[伊謬切; LH )iuC, OCM )iuh] phonetic in 怮 [於虯切; LH )eu, OCM iû]

213. SW 10B 408: 213; DXB 222 (10B 21a); GL vol 11: 4765; Duàn 513 (10B 46b); TKJ 1475; Ozaki vol. 5: 1079.


HG:
G: SSF: PIF:

(忦),
憂也。从心, 介聲。

jiá,
is 'to worry'.
It has xīn ``heart'' as a semantic constituent,
jiè (GSR 327a: ``armour, ...") is the phonetic constituent.



HG Compare also n° 139 xiè 3, which shows that Xǔ Shèn can declare one and the same phonetic constituent to refer to different but similar pronunciations.
Duàn again rewrites 憂 as ⤏.
PIF     介[古拜切; LH k/s, OCM krêts] phonetic in 忦 [五介切(《集韻》: 訖黠切); LH kat, OCM krât (or OCM krêt ?)]


214. SW 10B 408: 214; DXB 222 (10B 21a); GL vol 11: 4765; Duàn 513 (10B 46b); TKJ 1475; Ozaki vol. 5: 1079.
HG: 1 (恙), yàng,

G:	憂也。
SSF: 从心,
PIF: 羊聲。

is 'to worry'. [[is '(a way of) worrying (scil. particularly to have worries about one's health)'.]]
It has xīn ``heart'' as a semantic constituent,
yáng (GSR 732a: ``sheep") is the phonetic constituent.

G	The gloss comes from )ry( (Shìgǔ 釋詁, Xú Cháohuá 1994: 38).
Duàn again rewrites 憂 as ⤏.
PIF     羊[式羊切; LH jâng, OCM lang] phonetic in 恙 [余亮切; LH jângC, OCM langh]

215. SW 10B 408: 215; DXB 222 (10B 21a); GL vol 11: 4765b; Duàn 513 (10B 46b); TKJ 1475; Ozaki vol. 5: 1079.


HG:
G: SSF: PIF: IQ:

f (惴), 憂懼也。从心,
耑聲。
《詩》曰： 惴惴其慄。


zhuì,
is 'to worry and fear'.
It has xīn ``heart'' as a semantic constituent,
du+n (GSR 168a: ``tip, ...") is the phonetic constituent.
The Book of Odes (SSJZS: 373a) says: ``terrified is his trembling'' (tr. Karlgren).

G	Note Xǔ Shèn's analytic gloss. Duàn again rewrites 憂 as ⤏.
PIF     耑[多官切; LH tuân, OCM tôn] phonetic in 惴 [之瑞切; LH tshyaih, OCM toih]
PR	F*i sh*ng 非聲
IQ	Shījīng (Huáng ni(o 黃鳥 131.1). On two occasions Xǔ Shèn quotes texts involving the rather common character lì 慄 ``tremble with fear'' which he forgot to define in Shuōwén.



216. SW 10B 408: 216; DXB 222 (10B 21a); GL vol 11: 4766; Duàn 513 (10B 47a); TKJ 1476; Ozaki vol. 5: 1080.

HG: Š (憌),
G:	憂也。
SSF: 从心,
PIF: 鈞聲。


qióng,
is 'to worry'.
It has xīn ``heart'' as a semantic constituent,
jūn (GSR 391e: ``potter's wheel, ...") is the phonetic constituent.

G	Duàn again rewrites 憂 as ⤏.
PIF     鈞 [居fi切; LH kwin, OCM kwin] phonetic in 憌 [常倫切(《廣韻》: 渠營切); LH gyeng, OCM gweng]
PR	F*i sh*ng 非聲

217. SW 10B 408: 217; DXB 222 (10B 21a); GL vol 11: 4766; Duàn 513 (10B 47a); TKJ 1476; Ozaki vol. 5: 1080.

HG:  g (怲),
G:	憂也。
SSF: 从心,


b(ng,
is 'to worry'. [[is '(a way of) worrying (scil. in a poetic context)]]
It has xīn ``heart'' as a semantic constituent,

PIF: IQ:

丙聲。
《詩》曰： 憂心怲怲。

b(ng (GSR 757a: ``cyclical character") is the phonetic con- stituent.
The Book of Odes (SSJZS: 481c) says:
"the grief of the heart is intensive'' (tr. Karlgren).

G	Duàn again rewrites 憂 as ⤏.
PIF     丙[兵永切; LH pangB, OCM prang)] phonetic in 怲 [兵永切; LH pangB/C, OCM prang)/h]
IQ	Shījīng (Ku+ biàn ‰弁 217.4).

218. SW 10B 408: 218; DXB 222 (10B 21b); GL vol 11: 4766b; Duàn 513 (10B 47a); TKJ 1476; Ozaki vol. 5: 1080.


HG:
G: SSF: PIF: IQ:

\ (惔), 憂也。
从心, 炎聲。
《詩》曰： 憂心如惔。

tán, ``troubled'' is 'to worry'.
It has xīn ``heart'' as a semantic constituent,
yán (GSR 617a: ``blaze, ...") is the phonetic constituent.
The Book of Odes (SSJZS: 440a) says: ``Their worried hearts were as if troubled".

G	Duàn again rewrites 憂 as ⤏.



PIF    炎[于廉切; LH jam, OCM lam] phonetic in 䇢 [徒甘切; LH dâm, OCM lâm]
IQ Ancient commentators have discussed this character in Shījīng Jié nán sh&n 節南山 (191.1) and Zhèng Xuán 鄭玄 maintained that the reference was to the glow of fire. Karlgren follows Zhèng Xuán and translates: ``the grieved hearts are as if burning". This quotation is also cited under the character chán [䎌] (SW 10A 19a; TKJ: 1375). The coexistence of different versions of the same text in the illustrative quotations is probably to be explained by the fact that Xǔ Shèn was using different second hand sources which in turn were based on different editions of the Shījīng. This use of secondary sources does not exclude the likelihood that on many other occasions Xǔ Shèn is quoting the classics from memory.


219. SW 10B 408: 219; DXB 222 (10B 21b); GL vol 11: 4767; Duàn 513 (10B 47a); TKJ 1476; Ozaki vol. 5: 1081.


HG:
G:

SSF: PIF:

IQ:

^ (䓤),
憂也。
从心, 䎿聲。
《詩》曰：


chuò, ``grieved, sad"
is 'to worry'. [[is '(a way of) worrying (scil. in a poetic con- text)'.]]
It has xīn ``heart'' as a semantic constituent,
zhuó (GSR 295a: ``Shuowen says: to connect") is the phonetic constituent.
The Book of Odes (SSJZS: 286b) says:

"憂心䓤䓤。'' ``My sorrowful heart is very sad.'' (tr. Karlgren)

CAS: 一曰:
意不定也。

An (alternative) source says:
It is to be unsettled in one's mind.

G Here Xǔ Shèn fails to identify the use of chuò as a descriptive verb in the examples he himself provides. Thus it appears that the absence of mào 皃 in a gloss is in no way inconsistent with Xǔ Shèn's awareness that a word is currently used as a reduplicative descriptive word.
Duàn again rewrites 憂 as ⤏.
PIF     䎿[陟劣切; LH truât, OCM trot] phonetic in 䓤 [陟劣切; LH truât, OCM trot]
PR	Tóngyīn 同音 OCM and LH homophones.
IQ	Shījīng (C(o chóng 草蟲 14.2).

220. SW 10B 408: 220; DXB 222 (10B 21b); GL vol 11: 4767b; Duàn 513 (10B 47a); TKJ 1476; Ozaki vol. 5: 1081, see
variants in WGY: 458.


HG:
G: SSF:

(慯),
憂也。从心,


sh+ng,
is 'to worry'.
It has xīn ``heart'' as a semantic constituent,



PIF: 殤省聲。 it has an abbreviated form of sh+ng (GSR 720k: ``to die in childhood") as the phonetic.
G    Duàn again rewrites 憂 as ⤏.
PIF Duàn arbitrarily rewrites 殤 as 傷. In cases like these, when a current phonetic constituent has no independent reading as a graph for a word, Xǔ Shèn is forced to use his sh\%ng sh*ng 省聲formula. When there is a series of homophonous candidates to be used in this formula his choice of the graph he claims to have been abbreviated becomes arbitrary. Thus it makes no difference to his system whether the graph to have been abbreviated is claimed to be the common sh&ng 傷 or the slightly less common sh&ng 殤, or for that matter the rarer characters 觴, 漡, 鬺.
殤 [式陽切; LH shâng, OCM lhang] phonetic in 慯 [式亮切 (《廣韻》: 式羊切); LH shâng,
OCM lhang]


221. SW 10B 408: 221; DXB 222 (10B 21b); GL vol 11: 4768; Duàn 513 (10B 47b); TKJ 1477; Ozaki vol. 5: 1081.

HG: j (愁),
G:	憂也。
SSF: 从心,
PIF: 秋聲。


chóu, ``be very sad"
is 'to worry'. [[is '(a way of) worrying (scil. typically about a very sad thing of the past or present)'.]]
It has xīn ``heart'' as a semantic constituent,
qiū (GSR 1092a: ``autumn, ...") is the phonetic constituent.

G	Duàn again rewrites 憂 as ⤏.
PIF     秋[七由切; LH ts-hiu, OCM ts-hiu] phonetic in 愁 [士尤切; LH dzru, OCM dzru]

222. SW 10B 408: 222; DXB 222 (10B 21b); GL vol 11: 4768; Duàn 513 (10B 47b); TKJ 1477; Ozaki vol. 5: 1082.


HG:
G: SSF: PIF:

(愵),
憂皃。从心, 弱聲。

nì,
descriptive of being worried.
It has xīn ``heart'' as a semantic constituent,
ruò (GSR 1123a: ``weak, ...") is the phonetic constituent.

SPI:

讀與䮢同。 The pronunciation is the same as that of nì.

G	Duàn again rewrites 憂 as ⤏.
PIF     The phonetic ruò 弱 is commonly present as a phonetic in characters pronounced nì in modern Chinese: 溺, 嫋, 糑.
弱[而勺切; LH ñâk, OCM niauk] phonetic in 愵 [奴歷切; LH nek < neuk, OCM niâuk]
SPI     This unusual formula contrasts with the standard forms dúruò 讀若 and dúrú 讀如. In this instance, it is important to remember that an alternative gloss Xǔ Shèn provides for nì 䮢 is: 一曰



憂也 (n° 100). It is almost as if is suggesting that the characters nì 愵 and nì 䮢 are used to write the same words.
Nì 愵 [奴歷切; LH nek < neuk, OCM niâuk] read like nì 䮢 [奴歷切; LH nek, OCM niûk].

223. SW 10B 408: 223; DXB 222 (10B 21b); GL vol 11: 4768b; Duàn 513 (10B 47b); TKJ 1477; Ozaki vol. 5: 1082.


HG:
G: SSF: PIF:

(惂),
憂困也。从心,
;聲。

k&n,
is 'to worry about getting into trouble'.
It has xīn ``heart'' as a semantic constituent,
xiàn (GSR 672a: ``Shuowen: small pit (no text)") is the phonetic constituent.

G	Duàn again rewrites 憂 as ⤏. This binome gloss, we suspect, represents a colloquialism.
PIF     ; [戶䦤切; LH g/mB, OCM gr*+ m)] phonetic in 惂 [苦感切; LH kh\#mB, OCM kh*+ m)]

224. SW 10B 408: 224; DXB 222 (10B 21b); GL vol 11: 4768b; Duàn 513 (10B 47b); TKJ 1477; Ozaki vol. 5: 1082.

HG: J (悠),
G:	憂也。
SSF: 从心,
PIF: 攸聲。


yǒu,
is 'to worry'. [[is '(a way of) worrying (scil. intensely or wistfully, in a poetic mode)'.]]
It has xīn ``heart'' as a semantic constituent,
yǒu (GSR 1077a: ``place, ...") is the phonetic constituent.

G	Duàn again rewrites 憂 as ⤏.
PIF     攸[以周切; LH ju, OCM ju] phonetic in 悠 [以周切; LH ju, OCM ju]
PR	Tóngyīn 同音 OCM and LH homophones.

225. SW 10B 408: 225; DXB 222 (10B 21b); GL vol 11: 4769; Duàn 513 (10B 47b); TKJ 1477; Ozaki vol. 5: 1083.


HG:
G:

SSF: PIF: SPI:

N (悴),
憂也。
从心, 卒聲。
讀與《易》萃卦同。


cuì, ``despondent"
is 'to worry'. [[is '(a way of) worrying (scil. despondently, in a poetic mode)'.]]
It has xīn ``heart'' as a semantic constituent,
zú (GSR 490a: ``soldier, ...") is the phonetic constituent. The pronunciation is the same as that of the cuì hexagram.



G    Duàn again rewrites 憂 as ⤏.
PIF    卒[臧沒切; LH tsu\#t, OCM tsût] phonetic in 悴 [秦醉切; LH dzuis, OCM dzuts]
SPI    Again, Xǔ Shèn opts for a non-standard Supplementary Pronunciation Instruction formula, and it is as if he has got into this temporary habit on the occasion of the nì 愵 above. This phonetic gloss takes account of the fact that zú 卒 is phonetically ambiguous and Xǔ Shèn decides to disambiguate it in this case by juxtaposing it with another graph where 卒 has the same phonetic value.
    Cuì 悴 [秦醉切; LH dzuis, OCM dzuts] read like cuì 萃 [秦醉切; LH tsuet MOC] tsût] [see zāo2 遭 LH tsou, OCM tsû, in Schuessler 2009]

226. SW 10B 408: 226; DXB 222 (10B 21b); GL vol 11: 4769b; Duàn 513 (10B 47b); TKJ 1477; Ozaki vol. 5: 1083.


HG:
G: SSF: PIF: CAS:

} (õ),
憂也。从心, 圂聲。一曰: 擾也。


hùn,
'to worry'.
It has xīn ``heart'' as a semantic constituent,
hùn (GSR 425a: ``pig-sty, ...") is the phonetic constituent. An (alternative) source says:
It is to disturb.

G	Duàn again rewrites 憂 as ⤏.
PIF     圂[胡困切; LH gu\#nC, OCM gûns] phonetic in õ [胡困切; LH gu\#nC, OCM gûns]
CAS Duàn tacitly rewrites 擾 with an archaising graph ず.
This entry shows how Duàn's notes often have as an important source the K&ngxī zìdi(n
dictionary which among other things provides all the information he gives.


227. SW 10B 408: 227; DXB 222 (10B 21b); GL vol 11: 4770; Duàn 513 (10B 47b); TKJ 1477; Ozaki vol. 5: 1083.


HG:
DG:

SSF: PIF:

(⦩),
楚潁之閒謂憂曰⦩。 从心,
?聲。

lí,
In the area between Ch+ and Y1ng one says lí for yǒu ``wor- ry".
It has xīn ``heart'' as a semantic constituent,
lí (GSR 979a: ``Shuowen: split..") is the phonetic con- stituent.

DG	See F&ngyán 方言 1 (Zh\%u Zǔmó 1956: 3 (11) which writes: 楚潁之間謂之憖.
Duàn again rewrites 憂 as ⤏.



PIF     ? [許其切; LH li\#, OCM r\#] phonetic in ⦩ [力至切 (《廣韻》: 里之切); LH lis, OCM
*rits / LH li\#, OCM r\#]
PR	F*i sh*ng 非聲

228. SW 10B 408: 228; DXB 223 (10B 22a); GL vol 11: 4770b; Duàn 514 (10B 48a); TKJ 1477; Ozaki vol. 5: 1083.


HG:
G: SSF: PIF: SPI:

(❸),
憂也。从心, 于聲。
讀若吁。

xū,
is 'to worry'.
It has xīn ``heart'' as a semantic constituent,
yú (GSR 97a: ``to go, ...") is the phonetic constituent. pronounced like xū.

G	Duàn again rewrites 憂 as ⤏.
PIF     于 [羽倶切; LH wâ, OCM wa] (TKJ: 656 亏) phonetic in ❸ [況于切; LH hyâ, OCM hwa]
PR	F*i sh*ng 非聲
SPI     The character xū 吁 appears twice in Shuōwén under 口 and under 亏 radicals (TKJ: 202, 657).
Xū ❸ [況于切; LH hyâ, OCM hwa] read like xū 吁 [況于切; LH hya; MOC hwa]

229. SW 10B 408: 229; DXB 223 (10B 22a); GL vol 11: 4770b; Duàn 514 (10B 48a); TKJ 1478; Ozaki vol. 5: 1084.

HG: g (忡),
G:	憂也。
SSF: 从心,
PIF: 中聲。


chǒng,
is 'to worry'.
It has xīn ``heart'' as a semantic constituent,
zhǒng (GSR 1007a: ``middle, ...") is the phonetic constituent.

IQ:

《詩》曰： 憂心忡忡。

The Book of Odes (SSJZS: 286a) says:
"my grieved heart is agitated;'' (tr. Karlgren).

G    Duàn again rewrites 憂 as ⤏.
PIF According to Xǔ Shèn's analysis the characters 忡 and 忠 (Shuōwén 10B 408: 011), though differing in their initials, are taken to have the same phonetic. Moreover their grapholological analysis is taken to be identical: the placement of the radical is not considered an analytic issue. This situation is not at all unique.
中[陟弓切; LH trung, OCM trung] phonetic in 忡 [敕中切; LH thrung, OCM thrung]
PR   Homogeneous initials



IQ	The illustrative quotation of Shījīng (C(o chóng 草蟲 14.1) shows that Xǔ Shèn was aware that this is a descriptive word.


230. SW 10B 408: 230; DXB 223 (10B 22a); GL vol 11: 4771; Duàn 514 (10B 48a); TKJ 1478; Ozaki vol. 5: 1084, see
variants in WGY: 459.


HG:
G: SSF: PIF: IQ:

@ (悄), 憂也。从心,
肖聲。
《詩》曰： 憂心悄悄。


qi&o,
is 'to worry'.
It has xīn ``heart'' as a semantic constituent,
xiào (GSR 1149G: ``resemble") is the phonetic constituent.
The Book of Odes (SSJZS: 297a) says: ``My grieved heart is pained.'' (tr. Karlgren)

G	Duàn again rewrites 憂 as ⤏.
PIF     肖[私妙切; LH siâuC, OCM siauh] phonetic in 悄 [親小切; LH ts-hiâuB, OCM ts-hiau)]
IQ	The illustrative quotation of Shījīng (B(i zh\$u 柏舟 26.4) shows that Xǔ Shèn was aware that
qi(o 悄 is a descriptive word that could have been defined as 憂皃.

231. SW 10B 408: 231; DXB 223 (10B 22a); GL vol 11: 4771; Duàn 514 (10B 48a); TKJ 1478; Ozaki vol. 5: 1085.



HG:
G:


SSF: PIF:

(䮫),
憂也。

从心, 戚聲。

qī,
is 'to worry'. [[is '(a way of being?) being worried (scil. typically in the context of a very sad thing of the past, or in poetic contexts)'.]]
It has xīn ``heart'' as a semantic constituent,
qī (GSR 1031f: ``battle axe") is the phonetic constituent.

G	Duàn again rewrites 憂 as ⤏.
PIF     戚[倉歷切; LH ts-hek < ts-heuk, OCM ts-hiûk] phonetic in 䮫 [倉歷切; LH ts-hek < ts-heuk, OCM ts-hiûk]
PR	Tóngyīn Tóngyīn 同音 OCM and LH homophones.

232. SW 10B 408: 232; DXB 223 (10B 22a); GL vol 11: 4771b; Duàn 514 (10B 48a); TKJ 1478; Ozaki vol. 5: 1085, see
variants in WGY: 459.

HG:	(⤏),	yǒu, ``worry"




G: SSF1: SSF2:

愁也。从心, 从頁。

is 'to feel distressed'.
It has xīn ``heart'' as a semantic constituent,
it (also) has xié ``head'' as a semantic constituent,

G Note that the two words chóu 愁 (n° 221) and y\$u ⤏ are inter-defined in the pattern we call hùzhù 互注, and which we might as well have called hùxùn 互訓 in accordance with Duàn's practice, but they are written with inconsistent orthography. Xǔ Shèn inserts a word for unhappi- ness/sadness in the WORRY series.
SSF2 Duàn rewrites this as 从心頁, as if the pattern in the received text was not current in
Shuōwén.
HP      y\$u ⤏ [於求切; LH )u, OCM )u]

233. SW 10B 408: 233; DXB 223 (10B 22a); GL vol 11: 4772b; Duàn 514 (10B 48b); TKJ 1478; Ozaki vol. 5: 1086, see
variants in WGY: 459.

HG:  K (患),
G:	憂也。

SSF: 从心,
SGE: 上貫吅，


huàn,
is 'to worry'. [[is '(a way of) worrying (scil. mostly about prospects resulting from a past event or matters of the future, concrete or abstract)'.]]
It has xīn ``heart'' as a semantic constituent, and above it two mouths perforated,

PIF: SA1: SA2:

吅亦聲。
   [⧻]，古文从關省。[⫗]，亦
古文患。

xu+n (SW: ``shout in amazement") is at the same time the phonetic constituent.
  is ancient style graph and it has the abbreviated form of
gu+n (as phonetic).
is also ancient style graph for huàn.

G	Duàn again rewrites 憂 as ⤏.
SGE Xǔ Shèn graphological analysis is rarely descriptive in any detailed way, this being one of the exceptions to the rule. The present graph exhibits a phonetic constituent that is in fact interfered with by a super-imposed stroke going right through it. Cases like these are not common, nor is it common for the two mouths to be vertically rather than horizontally placed. For once, Xǔ Shèn comes close to giving a graphic description of a graph.
PIF     吅[況袁切; LH hu5n] phonetic in 患 [胡丱切; LH guanC, OCM grôns]
***



234. SW 10B 408: 234; DXB 223 (10B 22a); GL vol 11: 4774; Duàn 514 (10B 49a); TKJ 1479; Ozaki vol. 5: 1087.


HG:
G: SSF: PIF:

- (ƒ),	ku+ng, ``pusillanimous"
怯也。	is 'to be fearful'.
从心、匡， It has xīn ``heart'' and ku+ng ``fearful'' as a semantic constituents,
匡亦聲。 ku+ng (GSR 739m:"... square basket, ...") is at the same time the phonetic constituent.

G	Duàn tacitly rewrites this graph with a dog radical because it is that graph that is defined in
Shuōwén, as if Xǔ Shèn only used graphs which he has defined in his Shuōwén.
PIF     匡[去王切; LH khyâng, OCM khwang] phonetic in ƒ [去王切; LH khyâng, OCM khwang]
PR	Tóngyīn 同音 OCM and LH homophones.

235. SW 10B 408: 235; DXB 223 (10B 22a); GL vol 11: 4774b; Duàn 514 (10B 49a); TKJ 1479; Ozaki vol. 5: 1087, see
variants in WGY: 459.


HG:
G: SSF: PIF:

(徏),
思皃。从心, 夾聲。

qiè,
descriptive of wistful thoughtfulness.
It has xīn ``heart'' as a semantic constituent,
ji+ (GSR 630a: ``tweezers, ...") is the phonetic constituent.

PIF     夾[古狎切; LH k/p, OCM krêp] phonetic in 徏 [苦叶切; LH khep, OCM khêp]

236. SW 10B 408: 236; DXB 223 (10B 22a); GL vol 11: 4774b; Duàn 514 (10B 49a); TKJ 1479; Ozaki vol. 5: 1087.


HG:
G: SSF: PIF: CAS:

¹ (懾), 失气也。从心,
聶聲。一曰: 服也。


shè,
is 'to be disheartened'.
It has xīn ``heart'' as a semantic constituent,
niè (GSR 638a: ``promise") is the phonetic constituent. An (alternative) source says:
It is 'to submit'.

PIF     聶[之涉切; LH nrâp, OCM nrap] phonetic in 懾 [之涉切; LH tshap, OCM tap]
CAS   Duàn needlessly rewrites and expands this alternative gloss as 心服也.



***
FEAR SERIES 2 (237-(244)-248)
[Note that the same semantic field is often represented by several series within the same radical in
Shuōwén.]


237. SW 10B 408: 237; DXB 223 (10B 22b); GL vol 11: 4775; Duàn 514 (10B 49a); TKJ 1479; Ozaki vol. 5: 1087.


HG:
G: SSF: PIF: CAS:

œ (憚), 忌難也。从心,
單聲。一曰: 難也。


dàn,
is 'to abhor and find troublesome'.
It has xīn ``heart'' as a semantic constituent,
d+n (GSR 147a: ``simple, ...") is the phonetic constituent. An (alternative) source says:
It is 'to consider as wrought with trouble'.

G  In Shījīng the Máo commentary has the gloss 難也, and it is entirely possible that Xǔ Shèn is basing himself on this gloss. What is surprising, on the other hand, is that Xǔ Shèn finds it worth his while to add an alternative gloss which contributes nothing new.
PIF    單[都寒切; LH tân, OCM tân] phonetic in 憚 [徒案切; LH dânC, OCM dâns]

238. SW 10B 408: 238; DXB 223 (10B 22b); GL vol 11: 4775; Duàn 514 (10B 49a); TKJ 1480; Ozaki vol. 5: 1088.


HG:
G:

S (悼),
懼也。


dào,
is 'to fear' (!)

DG:

陳楚謂懼曰悼。 Between Chén and Ch+ they do (indeed) refer to fear by the word dào.

SSF: PIF:

从心, 卓聲。

It has xīn ``heart'' as a semantic constituent,
zhuó (GSR 1126a: ``high") is the phonetic constituent.

G Xǔ Shèn's gloss here is confusing because surely one might think that 悼 does not basically mean 'to fear' but 'to mourn for'. However, the meaning 'to fear' is in fact attested in Zhu&ngzi. What remains most puzzling is the question why Xǔ Shèn should have preferred one uncommon psychological meaning to another very common one.
DG	See F&ngyán 方言 1 (Zh\%u Zǔmó 1956: 3 (9)). In this gloss the word jù 懼 is not autonymic,
i.e. does not refer to itself but refers to the general concept expressed by it. Thus the meanings for words we give in quotation marks can ascribe the meaning to: 1. the graph X. 2. The word X. 3. The concept expressed by X.
PIF     卓[竹角切; LH trck, OCM trâuk] phonetic in 悼 [徒到切; LH dâuC , OCM dâukh]
PR	F*i sh*ng 非聲. Note the contrasting finals.



239. SW 10B 408: 239; DXB 223 (10B 22b); GL vol 11: 4776; Duàn 514 (10B 49b); TKJ 1480; Ozaki vol. 5: 1089.

HG:  (恐),
G:	懼也。
SSF: 从心,


kǒng, ``terrified"
is 'to fear'. [[is '(a away of being) afraid (scil. in the inten- sive mode of being terrified)'.]]
It has xīn ``heart'' as a semantic constituent,

PIF: SA:

m聲。
[❻]，古文。

gǒng (GSR 1172p: ``hold, ...") is the phonetic constituent.  is an ancient style graph.

PIF     m [居竦切; LH k\$ongB, OCM kong)] phonetic in 恐 [丘隴切; LH kh\$ongB, OCM khong)]
PR	Tóngyīn 同音 LH homophones but not OCM.

240. SW 10B 408: 240; DXB 223 (10B 22b); GL vol 11: 4776b; Duàn 514 (10B 49b); TKJ 1480; Ozaki vol. 5: 1089.


HG:
G: SSF: PIF: SPI:

’ (慴), 懼也。 从心,
習聲。 讀若疊。


zhé,
is 'to fear'.
It has xīn ``heart'' as a semantic constituent,
xí (GSR 690a: ``to practice, ...") is the phonetic constituent. Pronounced like dié.

PIF    習[似入切; LH zip, OCM s-l\#p] phonetic in 慴 [之涉切; LH tshap, OCM tap or tep]
SPI It is as if Xǔ Shèn notices that the phonetic constituent xí 習 [似入切] is not very helpful towards defining the pronunciation zhé 慴 [之涉切]. What is so curious is that by adding the third character dié 曡 [徒叶切] he does not seem to make things much easier.
Zhé 慴 [之涉切; LH tshap, OCM tap or tep] read like dié 疊 [曡: 徒叶切; LH dep, MOC lêp]

241. SW 10B 408: 241; DXB 223 (10B 22b); GL vol 11: 4777; Duàn 514 (10B 49b); TKJ 1480; Ozaki vol. 5: 1089.

HG: ƒ (䇟),
G:	恐也。
SSF: 从心,
PIF: 术聲。


chù, ``(morally) shocked'' (?)
is 'to be terrified'. [[is '(a way of) being terrified (scil. by the prospect or presence of injustice)'.]]
It has xīn ``heart'' as a semantic constituent,
shú (GSR 497a: ``glutinous millet") is the phonetic constituent.

HG   Note that this head graph is followed by its near synonym with which it forms a current binome (chùtì 䇟惕). Puzzlingly Xǔ Shèn provides a contrasting gloss for the second member of the binome.



PIF     术[食聿切; LH zhuit, OCM m-lut] (graphic variant of 䉽 in SW) phonetic in 䇟 [丑律切; LH thruit, OCM rhut ? or t-lhut ?]


242. SW 10B 408: 242; DXB 223 (10B 22b); GL vol 11: 4777; Duàn 514 (10B 49b); TKJ 1480; Ozaki vol. 5: 1089.


HG:
G:

] (惕),
敬[= read 驚]也。

tì, ``apprehensive'' is 'to be scared'.

SSF: 从心,

It has xīn ``heart'' as a semantic constituent,

PIF: A:

易聲。
[悐]，或从狄。

yì (GSR 850a: ``change, ...") is the phonetic constituent. in the graph , [the word tì] is alternatively written with the graphic constituent dí.

HG Here as often elsewhere, Xǔ Shèn deals with the members of a compound, in their order of occurrence, adjacently.
G    Based on the old Wénxu(n (Liù Chén zhù Wénxu(n 六臣注文選) commentary 說文曰：惕， 驚也 (Gǔlín ibidem) we tentatively suggest that jìng 敬 is to be read as jīng 驚 and rewrite the text accordingly.
    Note that Xǔ Shèn could have used the recurrent gloss 懼也 but did not, because he is in this case specifying the semantic nuance not by an analytic gloss but by the use of a specific rather than general word (if the emendation we have chosen is correct, that is).
PIF    易[以䋬切; LH jek, OCM lek] phonetic in 惕 [他歷切; LH thek, OCM lhêk]

243. SW 10B 408: 243; DXB 223 (10B 22b); GL vol 11: 4778; Duàn 514 (10B 49b); TKJ 1480; Ozaki vol. 5: 1090.


HG:	(待),
G:	戰慄也。
SSF: 从心,
PIF: 共聲。

hóng,
is 'to tremble with fear'.
It has xīn ``heart'' as a semantic constituent,
gòng (GSR 1182c: ``join the hands") is the phonetic constituent.

PIF     共[渠用切; LH k\$ongB, OCM kong)] phonetic in 待 [戶工切 又，工恐切; LH gong, OCM gông]


244. SW 10B 408: 244; DXB 223 (10B 22b); GL vol 11: 4778; Duàn 514 (10B 49b); TKJ 1480; Ozaki vol. 5: 1090.


HG:
G:

(徂),
苦也。

hài, ``troubled"
is 'to find something hard to bear'. [[is '(a way of) finding something hard to bear (scil. so as to come to fear it)'.]]




SSF: 从心,
PIF: 亥聲。

It has xīn ``heart'' as a semantic constituent,
hài (GSR 937a: ``cyclical character") is the phonetic con- stituent.

HG This entry is an intrusion into the ``fear'' series 2.
PIF     亥[古哀切; LH g\#B, OCM g*+ )] phonetic in 徂 [胡Ⓒ切; LH g\#C, OCM g*+ h]

245. SW 10B 408: 245; DXB 223 (10B 22b); GL vol 11: 4778b; Duàn 514 (10B 49b); TKJ 1480; Ozaki vol. 5: 1090.

HG: g (惶),
G:	恐也。
SSF: 从心,
PIF: 皇聲。


huáng, ``panic"
is 'to be terrified'. [[is '(a way of) being terrified (scil. to the point of panicking)'.]]
It has xīn ``heart'' as a semantic constituent,
huáng (GSR 708a: ``sovereign, ...") is the phonetic constituent.

PIF     皇[胡光切; LH 4uâng, OCM wâng] phonetic in 惶 [胡光切; LH 4uâng, OCM wâng]
PR	Tóngyīn 同音 OCM and LH homophones.

246. SW 10B 408: 246; DXB 223 (10B 22b); GL vol 11: 4778b; Duàn 514 (10B 49b); TKJ 1480; Ozaki vol. 5: 1090.


HG:	(悑),
G:	惶也。
SSF: 从心,

bù, ``scared out of one's mind'' is 'to panic'.
It has xīn ``heart'' as a semantic constituent,

PIF: A:

甫聲。
   [ 怖]， 或从布聲。

fǔ (GSR 102n: ``great, ...") is the phonetic constituent.
In the graph , [the word bù] is alternatively written with the graphic constituent bù which is phonetic.

G   The practice of defining a new head in terms of the last head word that has just been defined is common in Shuōwén.
PIF    甫[方矩切; LH puâB, OCM pa)] phonetic in 悑 [普故切; LH phâC, OCM phâh]
A The pattern 或从X聲 in cases where X is manifestly irrelevant semantically provides explicit evidence to prove that the technical term cóng 从 could occasionally introduce, in Shuōwén, elements of a purely phonetic kind.


247. SW 10B 408: 247; DXB 223 (10B 23a); GL vol 11: 4778b; Duàn 514 (10B 49b); TKJ 1481; Ozaki vol. 5: 1091.
HG:  ® (慹),   zhí,




G: SSF: PIF:

悑也。从心, 執聲。

is 'to be scared out of one's mind'.
It has xīn ``heart'' as a semantic constituent,
zhí (GSR 685a: ``grasp, ...") is the phonetic constituent.

G Xǔ Shèn glosses this entry with the terms he has just defined in the preceding entry, this concatenation proving that his text must be read consecutively and not as a conjuries of entries under radicals.
PIF Duàn writes 執 with an archaising graph ᠚, as if there was a convention in Shuōwén to the effect that phonetic constituents are listed in their seal form.
執[之入切; LH tship, OCM t\#p] phonetic in 慹 [之入切; LH tship, OCM t\#p]
PR   Tóngyīn 同音 OCM and LH homophones.

248. SW 10B 408: 248; DXB 223 (10B 23a); GL vol 11: 4779; Duàn 514 (10B 49b); TKJ 1481; Ozaki vol. 5: 1091.


HG:
G: SSF: PIF:

(⫞),
悑也。从心, 㱥聲。

qì,
is 'to be scared out of one's mind'.
It has xīn ``heart'' as a semantic constituent,
jī (GSR 854a: ``beat, ...") is the phonetic constituent.

G	Duàn rewrites bù 悑 as ⦦, whereupon he declares this entry and the following as a case of
zhu(n zhù 轉注.
PIF     毄[古歷切; LH khek, OCM khêk] phonetic in ⫞ [苦計切; LH kheiC, OCM khêh]
***

249. SW 10B 408: 249; DXB 223 (10B 23a); GL vol 11: 4778; Duàn 515 (10B 50a); TKJ 1481; Ozaki vol. 5: 1091, see
variants in WGY:460.


HG:	(⦦),
G:	⫞也。
SSF: 从心,

bèi,
is 'to be tired out'.
It has xīn ``heart'' as a semantic constituent,

PIF: A:

䷍聲。
   [俻]，或从疒。

bèi (GSR 984a: ``SW: prepare") is the phonetic constituent. In the graph , [the word bèi] is alternatively written with the graphic constituent nè.

PIF     ䷍[平祕切; LH b\$\#C, OCM br\#kh] phonetic in ⦦ [蒲拜切; LH b/C, OCM br*+ kh]



250. SW 10B 408: 250; DXB 223 (10B 23a); GL vol 11: 4778b; Duàn 515 (10B 50a); TKJ 1481; Ozaki vol. 5: 1092.

HG:  Y (惎),
G:	毒也。
SSF: 从心,
PIF: 其聲。


jì, ``aggressive"
is 'poisonous'. [[is '(a way of feeling) poisonous about (scil. figuratively so as to wish to harm others)'.]]
It has xīn ``heart'' as a semantic constituent,
qí (GSR 952a: ``Shuowen says: winnowing basket") is the phonetic constituent.

IQ:

《周書》曰： In the book of Zhǒu (SSJZS: 256b) it says:

來就惎惎。 ``They come and harbour hostile thoughts".
PIF   Qí 其 is not a head graph in Shuōwén. Xǔ Shèn clearly has problems analysing the structure of grammatical particles and needs to relate the graphs to more substantial meanings. He cannot have overlooked the fact that 其 is one of the most common characters in the language, but it so happens that from the point of view of graphological analysis the character is declared as no more than a graphic variant of jī 箕 which Xǔ Shèn treats as a radical and not, as one might expect, as containing the semantic constituent 竹 and the phonetic-cum-semantic constituent 其.
qí 其 [居之切; LH g\$\#, OCM g\# [qí: LH g\$\#, OCM g\#] (under jī 箕 in SW;  TKJ: 642),
[Gu(ngyùn qí: 渠之切; ] phonetic in 惎 [渠記切; LH g\$\#C, OCM g\#h]
IQ The SSJZS edition of Shūjīng (Qín Shì 秦誓 58.4) writes: 惟古之謀人則曰未就予忌 ``As to my former councillors, I (said=) considered that they did not accommodate themselves to me.'' (tr. Legge).

                 *** SHAME/HUMILIATION SERIES (251-256)

251. SW 10B 408: 251; DXB 223 (10B 23a); GL vol 11: 4780b; Duàn 515 (10B 50a); TKJ 1481; Ozaki vol. 5: 1092.

HG:	(恥), ch(, ``shame"

G:	辱也。
SSF: 从心,
PIF: 耳聲。

is 'humiliation'. [[is '(a kind of) humiliation (scil. of the subjective, psychological kind, as felt by the person humiliated)'.]]
It has xīn ``heart'' as a semantic constituent,
\$r (GSR 981a: ``ear") is the phonetic constituent.

PIF     耳 [而止切（《廣韻》: 而渉切）; LH ñ\#B, OCM n\#)] phonetic in 恥 [敕里切; LH thr\#B, OCM nhr\#) or rh\#)]



252. SW 10B 408: 252; DXB 223 (10B 23a); GL vol 11: 4781; Duàn 515 (10B 50a); TKJ 1481; Ozaki vol. 5: 1092.


HG:

(徬),

ti&n,

DG: 青徐謂慙曰徬。 In Qīng and Xú they call the feeling of shame ti&n.

SSF: 从心,
PIF: 典聲。

It has xīn ``heart'' as a semantic constituent,
di&n (GSR 476a: ``statute, ...") is the phonetic constituent.

DG Compare F&ngyán 方言 6 (Zh\%u Zǔmó 1956: 40 (5)), and see DG under n° 75 above. Xǔ Shèn's glosses of this type constitute his recognition not so much of dialect words as of local ways of speaking. (We note in passing that F&ngyán 方言 is not a dialect dictionary but a dictionary of localism.)
PIF  典[多殄切; LH tenB, OCM t*+ n)] phonetic in 徬 [他典切; LH thenB, OCM th*+ n)]

253. SW 10B 408: 253; DXB 223 (10B 23a); GL vol 11: 4781; Duàn 515 (10B 50a); TKJ 1482; Ozaki vol. 5: 1093.


HG:
G: SSF: PIF:

(彞),
辱也。从心, 天聲。

ti&n, ``??"
is 'humiliation'.
It has xīn ``heart'' as a semantic constituent,
ti+n (GSR 361a: ``Heaven") is the phonetic constituent.

HG If this word is viewed only on the basis of modern pronunciation it looks like a synonymous homophone of the preceding. The reconstructions show that the pronunciation of ti(n 徬 and ti(n 彞 were different. There are, incidentally, many examples in Shuōwén where the order of lexical entries seem inspired by phonetic associations.
PIF    天[他前切; LH then, OCM thǐn] phonetic in 彞 [他點切; LH thenB, OCM thǐn)]

254. SW 10B 408: 254; DXB 223 (10B 23a); GL vol 11: 4781b; Duàn 515 (10B 50a); TKJ 1482; Ozaki vol. 5: 1093.



HG:
G: SSF: PIF:

 (慙),
䑥也。从心, 斬聲。

cán,
is 'to feel ashamed'.
It has xīn ``heart'' as a semantic constituent,
zh&n (GSR 611a: ``cut off, ...") is the phonetic constituent.

G     Here as often elsewhere, the second member of current synonym compound is used to gloss the first one.
One might suspect that kuì 䑥 stands for kuì 愧, but in Xǔ Shèn's orthography the standard
form of the graph has the female radical and the graph with the heart radical enters the system only as an allograph unrelated to the heart radical and taken to contain an abbreviated form of ch+ 恥 as a graphic component, cf 12B 14a n° 236.
PIF    斬[側減切; LH tsr/mB, OCM tsrâm)] phonetic in 慙 [昨甘切; LH dzâm, OCM dzâm]



255. SW 10B 408: 255; DXB 223 (10B 23a); GL vol 11: 4781b; Duàn 515 (10B 50b); TKJ 1482; Ozaki vol. 5: 1093.


HG:
G: SSF: PIF:

½ (䮠), 慙也。从心,
而聲。


nu,
is 'feel ashamed'.
It has xīn ``heart'' as a semantic constituent,
ér (GSR 982a: ``whiskers (on an animal)") is the phonetic constituent.

PIF     而[如之切; LH ñ\#, OCM n\#] phonetic in 䮠 [女六切; LH nruk, OCM nruk]
PR	F*i sh*ng 非聲. But note the nasals initials.

256. SW 10B 408: 256; DXB 223 (10B 23a); GL vol 11: 4782; Duàn 515 (10B 50b); TKJ 1482; Ozaki vol. 5: 1094.


HG:
G: SSF: PIF:

  (䇜),
慙也。从心,
作省聲。


zuò,
is 'feel ashamed'.
It has xīn ``heart'' as a semantic constituent,
It has an abbreviated form of zuò (GSR 806l:"... make ...'' as the phonetic.

HG    Considering the sequence from 254 onwards we have first a definition of the hypernym cán 慙 and a sequence of words glosses in terms of this hypernym. Thus, we not only have the common phenomenon of a series of semantically related words, but in addition we find that these series themselves may exhibit an important internal structure: the first member of the series may identify the hypernym or most general concept and the items that follow may then subsume new lexical entries under this concept.
PIF     作[則洛切; LH dzraC, OCM dzrâh] phonetic in 䇜 [在各切; LH dzâk, OCM dzâk]
***

257. SW 10B 408: 257; DXB 223 (10B 23b); GL vol 11: 4782b; Duàn 515 (10B 50b); TKJ 1482; Ozaki vol. 5: 1094.
HG: ™ (憐), lián, ``to feel affection for"

G:	哀也。
SSF: 从心,
PIF: 䞑聲。

is 'to sympathise'. [[is '(a way of) sympathising (scil. affectionately)'.]]
It has xīn ``heart'' as a semantic constituent,
lín (GSR 387a: ``Shuowen says: will-o'-the-wisp, thus taking it to be the primary form of 燐") is the phonetic constituent.



PIF     Duàn writes 䞑 with an archaising graph 搽, as if there was a convention in Shuōwén to the effect that phonetic constituents are listed in their seal form.
䞑[良刃切; LH lin, OCM rin] phonetic in 憐 [落賢切; LH len, OCM rǐn]

258. SW 10B 408: 258; DXB 223 (10B 23b); GL vol 11: 4783; Duàn 515 (10B 50b); TKJ 1482; Ozaki vol. 5: 1094.


HG:	(尥),
G:	泣下也。
SSF: 从心,
PIF: 連聲。

lián,
is 'to fall' (of tears).
It has xīn ``heart'' as a semantic constituent,
lián (GSR 213a: ``a kind of carriage") is the phonetic constituent.

IQ:

《易》曰： The Book of Changes (Tún 屯; SSJZS: 20a) says:
泣涕尥如。 ``Tears streaming down".

HG The psychological relevance of the downward movement of tears is, of course, indirect only. One sympathise with the scribes who preferred to write this word with the water radical as does the received text of the Yìjīng.
PIF 連 [力延切; LH lianB, OCM ran? or ren)] phonetic in 忞 [力延切; LH lianB, OCM ran) or ren)]
PR   Tóngyīn 同音 OCM and LH homophones.
IQ The SSJZS edition of Yìjīng (Tún 屯.上六爻)(SSJZS: 20a) writes: (乘馬班如,)泣血漣如. Compare also Shījīng (Máng 氓 58 SSJZS 324c): 泣涕漣漣. It begins to look as if Xǔ Shèn, in his memory may have conflated these two texts.


259. SW 10B 408: 259; DXB 223 (10B 23b); GL vol 11: 4783b; Duàn 515 (10B 50b); TKJ 1482; Ozaki vol. 5: 1095.
HG: ; (忍), r\$n,

G:	能也。
SSF: 从心,
PIF: 刃聲。

is 'to be able'. [[is '(a way of being) able (scil. to put up with things that are hard to tolerate)'.]]
It has xīn ``heart'' as a semantic constituent,
rèn (GSR 456a: ``edge of a blade") is the phonetic constituent.

G	The gloss 'to be able' represent a curiously distant hypernym for r\%n 忍, but cases like this show up Xǔ Shèn's common strategy to gloss words in terms of the generic abstract hypernyms.
PIF     刃[而振切; LH ñ\$nC, OCM n\#ns] phonetic in 忍 [而軫切; LH ñ\$nB, OCM n\#n)]



260. SW 10B 408: 260; DXB 223 (10B 23b); GL vol 11: 4784; Duàn 515 (10B 50b); TKJ 1483; Ozaki vol. 5: 1095.


HG:
G:
AG:

SSF: PIF:


SPI:

(徺),
厲也。一曰: 止也。从心, 弭聲。

讀若䗽。

m(an,
is 'fierce'.
An (alternative) source says:
It is 'to stop'.
It has xīn ``heart'' as a semantic constituent,
m( (GSR 360a: ``bow with the ends not bound with string and lacquered but capped with bone and ivory") is the phonetic constituent.
Pronounced like mi&n.

AG	Duàn does not move this AG!
    For the reading ``stop'' there is no evidence in the transmitted literature. The graph used for that meaning is 弭.
PIF     Mi(n 弭[綿婢切; LH mieB, OCM me)] phonetic in mi(n 徺 [弥兖切(《廣韻》: 綿婢切); LH mieB, OCM me)]
PR Tóngyīn 同音 OCM and LH homophones. Schuessler notes: 弥兖切 = MC mjiwänB is a syllable which is impossible in MC, so either there is something wrong, or perhaps mjiänB was intended.
SPI Xǔ Shèn's dúruò 讀若 gloss demonstrates that he reads this character with a nasal final and the Gu(ngyùn fǎnqiè spelling designed to conform to the reading of the phonetic constituent and to create a plausible ancestor for the modern pronunciation m+ is irrelevant to Xǔ Shèn's analysis. The Gu(ngyùn prefers xiésh*ng consistency to straightforward observation. While there is a consider- able number of f*i sh*ng 非聲 phonetic constituents we still need a clear explanation why the 讀若 characters often had very different readings from the characters to which they are claimed to be similar.
    Mi(n 徺 [弥兖切(《廣韻》: 綿婢切); LH mieB, OCM me)] read as mi(n 䗽 [彌兖切, 《廣韻》: 綿婢切; LH mieB, OCM me)]

***
Minor CHASTISE SERIES (261-262)

261. SW 10B 408: 261; DXB 223 (10B 23b); GL vol 11: 4784b; Duàn 515 (10B 51a); TKJ 1483; Ozaki vol. 5: 1096.


HG:
G:

(※),
懲也。

yì, ``chastise'' is 'to punish'.




SSF: 从心,
PIF: 乂聲。

It has xīn ``heart'' as a semantic constituent,
yì (GSR 347a: ``Shuowen says: to mow, thus taking it to be the primary form of 刈 (no text)") is the phonetic constituent.

G	The head graph and the gloss constitute a case of adjacent hùzhù 互註.
PIF     乂[魚廢切; LH ng\$âs, OCM nga(t)s] phonetic in ※ [魚肺切; LH ng\$âs, OCM nga(t)s]
PR	Tóngyīn 同音 OCM and LH homophones.

262. SW 10B 408: 262; DXB 223 (10B 23b); GL vol 11: 47785; Duàn 515 (10B 51a); TKJ 1483; Ozaki vol. 5: 1096.

HG: µ (懲),
G:	※也。
SSF: 从心,
PIF: 徵聲。


chéng, ``punish'' is 'to chastise'.
It has xīn ``heart'' as a semantic constituent,
zh*ng (GSR 891a: ``verifications, ...") is the phonetic constituent.

PIF     徵[陟陵切; LH tr\$ng, OCM tr\#ng] phonetic in 懲 [直陵切; LH dr\$ng, OCM dr\#ng]
***

263. SW 10B 408: 263; DXB 223 (10B 23b); GL vol 11: 4785b; Duàn 515 (10B 51a); TKJ 1483; Ozaki vol. 5: 1096, see
variants in WGY: 461.

HG:  ¡ (憬),
G:	覺寤也。
SSF: 从心,


j(ng,
is 'to come to have a full understanding of'. It has xīn ``heart'' as a semantic constituent,

PIF: IQ:

景聲。
《詩》曰： 憬彼淮夷。

j(ng (GSR 755d: ``bright, ...") is the phonetic constituent.
The Book of Odes (SSJZS: 612b) says: ``We understand the Huái barbarians'' ?

PIF     景[居影切; LH k\$angB, OCM krang)] phonetic in 憬 [倶永切; LH kyangB, OCM kwrang) ?]
IQ	Karlgren has a less personal interpretation of this line of Shījīng (Lǔ sòng 魯頌, Pàn shu+ 泮水 299.8): ``far away are those Huái tribes".










4. Appendix

Radicals Containing the Grapheme 心not Classified Under the Heart Radical


As mentioned before there are a number of graphs which one would expect to be classified under the heart radical but which for some reason are not so classified. Some of these characters cannot be so classified because they themselves are radicals, there are listed below:

The radical suǒ 惢部

1. SW 409: 001; DXB 223 (10B 23b); GL vol 11: 4791; Duàn 515 (10B 51a); TKJ 1484; Ozaki vol. 5: 1149.

HG:  c (惢),
G:	心疑也。
SSF: 从三心。


suǒ,
is 'to be doubtful in one's mind'.
It has three hearts as a semantic constituents.

SF: SPI:

凡惢之屬皆从惢。
讀若《 易》'' 旅瑣瑣"。

As a matter of principle, all (graphs) classified un- der suǒ have suǒ as a semantic constituent. pronounced like Lu suǒsuǒ in the Yìjīng.

HG This may look like an exquisite example of a radical Xǔ Shèn could easily have done without. But if we consider the possibilities Xǔ Shèn was facing when trying his analytic hands on the following character ru+ 繠 we can see that an immediate constituent consisting of two hearts and the seal radical was not something he would have contemplated as a plausible option. Given the meaning 'to hang' he could surely have declared sī 糸 to be the radical and the three hearts an inscrutable phonetic as he does in many other cases of inscrutable phonetics. Moreover, as he has it, the meaning of the radical is totally unrelated to that of the one character declared to have it as a radical. On the other hand, his radical suǒ itself was amenable to a rather elegant interpretation in terms of its three heart constituents. Thus the presence of suǒ under the heart radical would certainly not have raised any eyebrows but, unfortunately, this suǒ would have been a most inadequate phonetic constituent in ru+, the presence of which under the silk radical, in turn would not have been surprising. To summarise, a character like ru+ 繠 constituted an insolvable analytic




dilemma for Xǔ Shèn and he appears to have escaped from this dilemma by working towards his numerologically desirable number of radicals.
SPI     Suǒ 惢 [又，才規、 才累二切(《玉篇》：桑果切) LH suɑiᴮ, OCM snôiʔ]] read like suǒ 瑣
[蘇果切 LH su5i6, OCM sôi) or snôi) ?]
Xǔ Shèn assumes that the reader would be familiar with the pronunciation of the Yìjīng text.
See Yìjīng Lu guà 旅卦 (SSJZS: 68b).

2 SW 409: 002; DXB 223 (10B 23b); GL vol 11: 4792; Duàn 515 (10B 51a); TKJ 1484; Ozaki vol. 5: 1149.


HG:
¾ (繠),
ru(
G:
垂也。
is 'to hang'.
SSF:
从惢,
It has suǒ as a semantic constituent.
PIF:
糸聲。
mì (SW: ``silk thread...") is the phonetic constituent.
PIF     糸[莫狄切; LH mek < ?] phonetic in 繠 [如壘切; LH ñuiB, OCM nui)]
PR	F*i sh*ng 非聲.


The radical sī 思部

3. SW 10B 407: 001; DXB 216 (10B 9b); GL vol 11: 4643b; Duàn 501 (10B 23a); TKJ 1437; Ozaki vol. 5: 990.

HG:             (思),
G:	容也。
SSF: 从心

sī, ``to think"
is 'to cause to encompass??' [[is '(a way of) causing to encom- pass (scil. subjects of thought in one's mind) ?? ]].
It has xīn ``heart'' as a semantic constituent.

PIF: SF:

4聲。 凡思之屬皆从思。

xìn (SW: ``fontanelle") is the phonetic constituent.
As a matter of principle, all (graphs) classified under sī have sī
as a semantic constituent.

PIF     Note that the radical preceding sī 思 is the very rare xìn 4 which is here declared to be phonetic.
4 [息進切; LH sinC, OCM sins or s\#ns] is phonetic in 思 [息兹切, LH s\$\#, OCM s\#].
PR	F*i sh*ng 非聲


CHINESE LEXICOGRAPHY ON MATTERS OF THE HEART	167
4. SW 10B 407: 002; DXB 217 (10B 10a); GL vol 11: 4645b; Duàn 501 (10B 23b); TKJ 1437; Ozaki vol. 5: 991.
HG:  (慮), lu, ``to think about the future"

G:	謀思也。SSF:    从思PIF: 虍聲。

is 'to engage in planning thought'.
It has xīn ``heart'' as a semantic constituent.
hǔ (GSR 57a: ``Shuowen: the streaks on a tiger (no text)") is the phonetic constituent.

